% ==============================================================================
% CHƯƠNG 5: BIỂU MẪU WEB (HTML5 FORMS) & CÁC Ô NHẬP LIỆU
% ==============================================================================
\chapter{Biểu Mẫu Web (HTML5 Forms) \& Các Ô Nhập Liệu}
\label{chap:forms_accessibility}

\section{Tổng Quan Về HTML Forms \& Thuộc Tính Quan Trọng}
Form là thành phần tương tác quan trọng nhất để người dùng gửi dữ liệu lên máy chủ. Một biểu mẫu bao gồm nhiều loại ô nhập liệu khác nhau phục vụ các mục đích dữ liệu riêng biệt.

% ------------------------------------------------------------------------------
% SUBSECTION: THUỘC TÍNH ACTION TRONG THẺ FORM (FORM ACTION LÀ GÌ?)
% ------------------------------------------------------------------------------
\subsection{Thuộc Tính action Trong Thẻ <form> (Form Action Là Gì?)}
\label{subsec:form_action}

\begin{conceptBox}[Thuộc Tính action Trong HTML Form Là Gì?]
\term{Thuộc tính action} của thẻ \codeinline{<form>} xác định địa chỉ \term{URL} (Uniform Resource Locator) mà dữ liệu trong biểu mẫu sẽ được đóng gói và gửi tới khi người dùng kích hoạt sự kiện Submit (Gửi).
\begin{itemize}[leftmargin=1.5em, itemsep=2pt, parsep=0pt]
    \item \textbf{Nói cách khác}: Đây là điểm cuối (\term{Endpoint}) tiếp nhận và xử lý dữ liệu biểu mẫu trên máy chủ Backend (thường là một tập tin xử lý như PHP, Python, Node.js, Java hoặc một điểm cuối RESTful API).
    \item \textbf{Cơ chế hoạt động bên dưới}: Khi sự kiện submit xảy ra, trình duyệt web dừng luồng tương tác hiện tại, quét qua tất cả các phần tử nhập liệu hợp lệ trong thẻ \codeinline{<form>}, gom các cặp \codeinline{name=value}, mã hóa dữ liệu theo định dạng \codeinline{enctype} và khởi tạo một HTTP Request chuyển hướng tới địa chỉ URL khai báo trong \codeinline{action}.
\end{itemize}
\end{conceptBox}

% ------------------------------------------------------------------------------
\paragraph{Sơ Đồ Luồng Hoạt Động Của Form Action:}

\begin{center}
\begin{tikzpicture}[node distance=1.6cm, auto,
    block/.style={rectangle, draw=primaryBlue, fill=primaryBlue!6, text width=3.0cm, text centered, rounded corners=3pt, minimum height=1.2cm, font=\sffamily\small},
    cloud/.style={rectangle, draw=codeComment, fill=codeComment!10, text width=3.0cm, text centered, rounded corners=3pt, minimum height=1.2cm, font=\sffamily\small},
    line/.style={draw=primaryBlue, ->, line width=1.2pt}]
    
    \node [block] (input) {\textbf{1. Người Dùng}\\Nhập liệu \& Bấm Submit};
    \node [block, right of=input, node distance=3.7cm] (browser) {\textbf{2. Trình Duyệt}\\Gom cặp \codeinline{name=value}};
    \node [block, right of=browser, node distance=3.7cm] (action) {\textbf{3. HTTP Request}\\Gửi dữ liệu tới \codeinline{action}};
    \node [cloud, right of=action, node distance=3.7cm] (backend) {\textbf{4. Backend Server}\\Tiếp nhận \& Trả phản hồi};
    
    \path [line] (input) -- (browser);
    \path [line] (browser) -- (action);
    \path [line] (action) -- (backend);
\end{tikzpicture}
\end{center}

% ------------------------------------------------------------------------------
\paragraph{Phân Loại Địa Chỉ URL Trong Thuộc Tính action:}

Khi khai báo đường dẫn \codeinline{action}, lập trình viên có thể sử dụng các dạng đường dẫn sau:

\begin{prettyTableBox}[Phân Loại Các Dạng Đường Dẫn Trong Thuộc Tính action]
\rowcolors{2}{white}{primaryBlue!4}
\begin{tabularx}{\linewidth}{K{3.6cm} Z Z}
\rowcolor{primaryBlue}
\thd{Dạng Đường Dẫn} & \thd{Cú Pháp HTML Mẫu} & \thd{Đặc Điểm \& Hành Vi Xử Lý} \\
\textbf{Đường Dẫn Tương Đối}\newline(\term{Relative URL}) & \codeinline{<form action="process.php">}\newline\codeinline{<form action="../save.php">} & Trỏ tới tệp nằm cùng thư mục hoặc thư mục cha/con trên cùng server. Dễ di chuyển mã nguồn giữa các môi trường (Dev/Staging/Prod). \\
\textbf{Đường Dẫn Từ Gốc Root}\newline(\term{Root-Relative URL}) & \codeinline{<form action="/api/v1/users">} & Trỏ tới địa chỉ tính từ tên miền gốc (\term{Domain Root}). Đảm bảo không bị lỗi đường dẫn khi nhúng form ở các trang con có độ sâu cấp thư mục khác nhau. \\
\textbf{Đường Dẫn Tuyệt Đối}\newline(\term{Absolute URL}) & \codeinline{<form action="https://api.domain.com/submit">} & Trỏ trực tiếp tới một tên miền hoặc máy chủ API bên ngoài. Bắt buộc phải khởi chạy qua giao thức an toàn \codeinline{https://}. \\
\textbf{Chuỗi Rỗng Hoặc Bỏ Trống}\newline(\term{Self-Submitting Form}) & \codeinline{<form action="">}\newline\codeinline{<form>} & Trình duyệt mặc định gửi toàn bộ dữ liệu biểu mẫu về chính URL của trang hiện tại. Rất phổ biến khi làm việc với PHP/ASP.NET monolithic hoặc xử lý form 2 giai đoạn (GET hiển thị, POST xử lý). \\
\textbf{Dấu Thăng (\codeinline{action="\#"})}\newline(\term{Hash Navigation}) & \codeinline{<form action="\#">} & Gửi dữ liệu về trang hiện tại nhưng đính kèm thẻ \codeinline{\#} làm trình duyệt tự động cuộn trang lên trên cùng (\term{Top Page}). Không khuyến nghị dùng trong thực tế trừ khi có xử lý JS. \\
\end{tabularx}
\end{prettyTableBox}

% ------------------------------------------------------------------------------
\paragraph{Mối Quan Hệ Cốt Lõi Giữa action, method và enctype:}

Thuộc tính \codeinline{action} không hoạt động độc lập mà luôn kết hợp chặt chẽ với phương thức gửi dữ liệu \codeinline{method} và kiểu mã hóa dữ liệu \codeinline{enctype}:

\begin{enumerate}[leftmargin=1.5em, itemsep=6pt, parsep=0pt]
    \item \textbf{Kết hợp với \codeinline{method="GET"}}:
    Trình duyệt đính kèm toàn bộ dữ liệu vào sau địa chỉ trong \codeinline{action} dưới dạng chuỗi truy vấn (\term{Query String}):
    \begin{center}
    \codeinline{URL\_action?name1=value1\&name2=value2}
    \end{center}
    \textit{Ứng dụng}: Dùng cho các biểu mẫu tìm kiếm, lọc dữ liệu, phân trang mà thông tin không cần bảo mật và cho phép người dùng lưu bookmark URL.

    \item \textbf{Kết hợp với \codeinline{method="POST"}}:
    Trình duyệt đính kèm dữ liệu ẩn bên trong thân của HTTP Request (\term{Request Body}).
    \textit{Ứng dụng}: Dùng cho đăng nhập, đăng ký, thanh toán, cập nhật thông tin cá nhân.

    \item \textbf{Ảnh hưởng của thuộc tính \codeinline{enctype} khi gửi tới action}:
    \begin{itemize}[leftmargin=1.2em, itemsep=2pt]
        \item \codeinline{application/x-www-form-urlencoded} (\textit{Mặc định}): Tất cả khoảng trắng đổi thành \codeinline{+} hoặc \codeinline{\%20}, các ký tự đặc biệt được mã hóa theo dạng Hex.
        \item \codeinline{multipart/form-data}: \textbf{Bắt buộc phải dùng} khi form có ô tải tệp tin (\codeinline{<input type="file">}) để dữ liệu tệp nhị phân được phân chia thành từng khối gửi tới \codeinline{action}.
        \item \codeinline{text/plain}: Gửi dữ liệu dưới dạng văn bản thô không mã hóa (chỉ phục vụ thử nghiệm).
    \end{itemize}
\end{enumerate}

% ------------------------------------------------------------------------------
\paragraph{Ví Dụ Minh Họa Chi Tiết:}

\begin{enumerate}[leftmargin=1.5em, itemsep=6pt, parsep=0pt]
    \item \textbf{Gửi đến file xử lý khác (\codeinline{save\_data.php}):}
    
\begin{codeWindow}[Form gửi dữ liệu tới file xử lý save\_data.php]
\begin{lstlisting}[language=HTML5]
<form action="save_data.php" method="post">
  <input type="text" name="email" placeholder="(*@Nhập email@*)">
  <button type="submit">(*@Đăng ký@*)</button>
</form>
\end{lstlisting}
\end{codeWindow}

\begin{browserWindow}[Kết Quả Trình Duyệt: Gửi Đến File save\_data.php]
\sffamily\small
\tcbox[on line, arc=2pt, colback=white, colframe=primaryBlue!40, boxsep=1pt, left=6pt, right=30pt, top=3pt, bottom=3pt]{\color{gray!80}Nhập email} \;\;
\tcbox[on line, arc=3pt, colback=primaryBlue, colframe=primaryBlue, boxsep=2pt, left=10pt, right=10pt, top=3pt, bottom=3pt]{\color{white}\bfseries\small Đăng ký}
\end{browserWindow}

\textbf{$\rightarrow$} \textit{Kết quả khi submit}: Dữ liệu nhập ở ô email sẽ được gửi đến tập tin \codeinline{save\_data.php}.

    \item \textbf{Gửi đến chính trang hiện tại (\codeinline{action=""}):}

\begin{codeWindow}[Form tìm kiếm gửi dữ liệu về chính trang hiện tại]
\begin{lstlisting}[language=HTML5]
<form action="" method="get">
  <input type="text" name="q" placeholder="(*@Tìm kiếm...@*)">
  <button type="submit">Search</button>
</form>
\end{lstlisting}
\end{codeWindow}

\begin{browserWindow}[Kết Quả Trình Duyệt: Form Tìm Kiếm Gửi Về Trang Hiện Tại]
\sffamily\small
\tcbox[on line, arc=2pt, colback=white, colframe=primaryBlue!40, boxsep=1pt, left=6pt, right=30pt, top=3pt, bottom=3pt]{\color{gray!80}Tìm kiếm...} \;\;
\tcbox[on line, arc=3pt, colback=primaryBlue, colframe=primaryBlue, boxsep=2pt, left=10pt, right=10pt, top=3pt, bottom=3pt]{\color{white}\bfseries\small Search}
\end{browserWindow}

\textbf{$\rightarrow$} \textit{Kết quả khi submit}: Trình duyệt sẽ đổi URL trên thanh địa chỉ thành dạng:
\begin{center}
\codeinline{http://example.com/?q=giatri\_nhap}
\end{center}
\end{enumerate}

% ------------------------------------------------------------------------------
\paragraph{Lưu Ý Quan Trọng \& Cảnh Báo Bảo Mật (Security Best Practices):}

\begin{warnBox}[Các Nguyên Tắc Bảo Mật Cần Nắm Vững Khi Cấu Hình Form Action]
\begin{enumerate}[leftmargin=1.5em, itemsep=3pt, parsep=0pt]
    \item \textbf{Bắt buộc sử dụng giao thức HTTPS trong action}: Không bao giờ trỏ thuộc tính \codeinline{action} tới đường dẫn giao thức không an toàn (\codeinline{http://}) trong sản phẩm thực tế để tránh nguy cơ bị nghe lén dữ liệu trên đường truyền (\term{Man-in-the-Middle Attack}).
    \item \textbf{Phòng chống lỗ hổng Cross-Site Request Forgery (CSRF)}: Khi \codeinline{action} trỏ tới các điểm cuối thay đổi dữ liệu hệ thống (như đổi mật khẩu, chuyển tiền), bắt buộc Backend phải yêu cầu đính kèm \term{CSRF Token} hợp lệ bên trong form.
    \item \textbf{Tránh lỗ hổng Open Redirect}: Không bao giờ cấu hình thuộc tính \codeinline{action} nhận trực tiếp giá trị URL biến động truyền từ tham số chuỗi truy vấn người dùng nhập mà không qua kiểm tra (\term{Whitelist validation}).
    \item \textbf{Phân công vai trò giữa Frontend và Backend}:
    \begin{itemize}[leftmargin=1.2em, itemsep=2pt]
        \item \textbf{Frontend}: Thiết kế giao diện, đảm bảo liên kết thuộc tính \codeinline{action} chính xác, thu thập và kiểm tra sơ bộ tính hợp lệ (\term{Client-side Validation}).
        \item \textbf{Backend}: Xây dựng bộ định tuyến (\term{Router/Endpoint}) tại đúng đường dẫn khai báo trong \codeinline{action}, kiểm tra tính hợp lệ chặt chẽ (\term{Server-side Validation}) và lưu trữ cơ sở dữ liệu.
    \end{itemize}
\end{enumerate}
\end{warnBox}

% ------------------------------------------------------------------------------
\paragraph{Minh Họa Thực Tế \& Kịch Bản Nhập Thử (Form Action Demonstration):}

\begin{codeWindow}[Mã Nguồn Đầy Đủ Minh Họa Form Action Với Kịch Bản Nhập Thử]
\begin{lstlisting}[language=HTML5]
<!-- Form dang ky tai khoan truyen du lieu den server handler -->
<form action="process_register.php" method="POST">
  <label for="user_email">Email dang ky:</label>
  <input type="email" id="user_email" name="email" placeholder="(nhập_email@example.com)" required>
  
  <label for="user_pass">Mat khau:</label>
  <input type="password" id="user_pass" name="password" placeholder="******" required>
  
  <button type="submit">Dang ky tai khoan</button>
</form>
\end{lstlisting}
\end{codeWindow}

\begin{browserWindow}[Kết Quả Hiển Thị Trình Duyệt: Minh Họa Form Action Nhập Thử]
\sffamily\small
\textbf{Form Đăng Ký Tài Khoản (Gửi tới: \codeinline{process\_register.php}):} \\[6pt]
\begin{tabular}{ll}
\textbf{Email đăng ký:} & \tcbox[on line, arc=2pt, colback=white, colframe=primaryBlue!50, boxsep=1pt, left=6pt, right=40pt, top=3pt, bottom=3pt]{\small \textcolor{black}{nguyenvana@gmail.com}} \\
\textbf{Mật khẩu:} & \tcbox[on line, arc=2pt, colback=white, colframe=primaryBlue!50, boxsep=1pt, left=6pt, right=40pt, top=3pt, bottom=3pt]{\small \textcolor{black}{$\bullet\bullet\bullet\bullet\bullet\bullet\bullet\bullet$}} \\
\end{tabular} \\[8pt]
\tcbox[on line, arc=3pt, colback=primaryBlue, colframe=primaryBlue, boxsep=3pt, left=12pt, right=12pt, top=4pt, bottom=4pt]{\color{white}\bfseries\small Đăng ký tài khoản}

\vspace{6pt}
\hrule height 0.4pt \relax
\vspace{6pt}
\textbf{\textcolor{primaryBlue}{Mô Phỏng Trạng Thái Khi Người Dùng Nhấn Nút Submit ("Nhập Thử"):}}
\begin{itemize}[leftmargin=1.2em, itemsep=2pt]
    \item \textbf{Trình duyệt khởi tạo}: HTTP POST Request gửi tới URL \codeinline{process\_register.php}.
    \item \textbf{Dữ liệu đóng gói gửi đi}: \codeinline{email=nguyenvana\%40gmail.com\&password=******}
    \item \textbf{Trạng thái xử lý Backend}: Server xử lý tệp \codeinline{process\_register.php} nhận dữ liệu, tiến hành mã hóa mật khẩu và trả về phản hồi chào mừng.
\end{itemize}
\end{browserWindow}

% ------------------------------------------------------------------------------
% MỤC ĐỌC THÊM (MỞ RỘNG KIẾN THỨC THEO QUY TẮC GIÁO TRÌNH)
% ------------------------------------------------------------------------------
\subsection{Đọc Thêm: Các Kỹ Thuật Nâng Cao Với Form Action Trong Lập Trình Modern Web}

\begin{conceptBox}[Ghi Đè Đường Dẫn Action Bằng Thuộc Tính formaction (HTML5)]
Trong đặc tả HTML5, nếu một form có nhiều nút bấm gửi dữ liệu cho các mục đích khác nhau (ví dụ: Nút "Gửi bài" và Nút "Lưu nháp"), ta có thể dùng thuộc tính \codeinline{formaction} ngay trên nút submit để ghi đè thuộc tính \codeinline{action} chung của form:
\begin{codeWindow}[Ví dụ sử dụng thuộc tính formaction trên từng nút submit]
\begin{lstlisting}[language=HTML5]
<form action="/save-article" method="post">
  (*@<!-- Nut 1: (*@Gửi@*) du lieu den /save-draft -->@*)
  <button type="submit" formaction="/save-draft">Luu Nhap</button>
  
  <!-- Nut 2: Su dung action mac dinh cua form (/save-article) -->
  <button type="submit">Xuat Ban</button>
</form>
\end{lstlisting}
\end{codeWindow}
\end{conceptBox}

\begin{tipBox}[Xử Lý Form Action Trong Các Ứng Dụng Trang Đơn (SPA / AJAX)]
Trong lập trình Frontend hiện đại với JavaScript (React, Vue, Angular, Fetch API/Axios):
\begin{itemize}[leftmargin=1.5em, itemsep=2pt, parsep=0pt]
    \item Lập trình viên thường gọi hàm \codeinline{event.preventDefault()} trong sự kiện \codeinline{onsubmit} để \textbf{chặn hành vi điều hướng trang mặc định} của thuộc tính \codeinline{action}.
    \item Dữ liệu sau đó được đóng gói và gửi qua AJAX/Fetch API một cách bất đồng bộ mà không làm tải lại trang web, giúp tối ưu trải nghiệm người dùng (\term{UX}).
\end{itemize}
\end{tipBox}

\begin{noteBox}[Góc Phỏng Vấn: Các Câu Hỏi Thường Gặp Về Form Action]
\begin{enumerate}[leftmargin=1.5em, itemsep=3pt, parsep=0pt]
    \item \textbf{Phân biệt action="", action="\#" và bỏ qua thuộc tính action?}
    \begin{itemize}[leftmargin=1.2em, itemsep=2pt]
        \item \codeinline{action=""} và bỏ trống \codeinline{action}: Đều gửi form về chính URL trang hiện tại mà không làm đổi thẻ thăng.
        \item \codeinline{action="\#"}: Gửi dữ liệu về trang hiện tại nhưng đính kèm vị trí \codeinline{\#} khiến trình duyệt tự cuộn lên đầu trang.
    \end{itemize}
    \item \textbf{Tại sao thẻ <form> vẫn cần thuộc tính action khi làm việc với React/Vue/Angular?}
    \begin{itemize}[leftmargin=1.2em, itemsep=2pt]
        \item Để đảm bảo tính năng \textbf{Progressive Enhancement} (Cải tiến tiệm tiến) và \textbf{Accessibility (a11y)}: Nếu JavaScript của người dùng bị tắt hoặc gặp sự cố tải, form vẫn có thể gửi dữ liệu lên server backend thông qua đường dẫn \codeinline{action} truyền thống.
    \end{itemize}
\end{enumerate}
\end{noteBox}

% ==============================================================================
% MỤC GIÁO TRÌNH: CHI TIẾT VỀ CHECKBOX trong HTML (<input type="checkbox">)
% ==============================================================================
\section[Thẻ Ô Chọn Checkbox]{Thẻ Ô Chọn Checkbox (<input type="checkbox">)}
\label{sec:html_checkbox}

\begin{conceptBox}[Định Nghĩa Checkbox]
\term{Checkbox} là một loại phần tử nhập liệu (\codeinline{<input>}) dạng ô chọn nhiều (\term{multi-choice}), cho phép người dùng chọn \textbf{0, 1 hoặc nhiều lựa chọn} trong cùng một nhóm.
\begin{itemize}[leftmargin=1.5em, itemsep=2pt, parsep=0pt]
    \item \textbf{Khác biệt với Radio Button}: Radio button chỉ cho phép chọn duy nhất 1 trong nhiều lựa chọn (\term{single-choice}).
    \item \textbf{Cơ chế gửi dữ liệu}: Nếu được tích chọn, checkbox sẽ gửi giá trị tương ứng về máy chủ. Nếu không được chọn, giá trị của nó sẽ không được gửi đi.
    \item \textbf{Ứng dụng phổ biến}: Chọn nhiều tùy chọn (sở thích, kỹ năng, ngôn ngữ lập trình...), xác nhận đồng ý điều khoản, hoặc bật/tắt một tính năng.
\end{itemize}
\end{conceptBox}

% ------------------------------------------------------------------------------
\subsection{Cú Pháp Cơ Bản \& Thuộc Tính Thường Gặp}

\begin{codeWindow}[Cú pháp tạo ô chọn Checkbox đơn giản]
\begin{lstlisting}[language=HTML5]
<input type="checkbox" name="vehicle" value="Bike"> (*@Tôi có xe đạp@*)
<input type="checkbox" name="vehicle" value="Car"> (*@Tôi có ô tô@*)
\end{lstlisting}
\end{codeWindow}

\begin{browserWindow}[Kết Quả Hiển Thị Trình Duyệt: Checkbox Cú Pháp Cơ Bản]
\sffamily\small
\begin{tikzpicture}[baseline=-0.4ex]\draw[primaryBlue!70, line width=1pt, rounded corners=2pt] (0,0) rectangle (12pt,12pt);\end{tikzpicture} \;\textbf{Tôi có xe đạp} \\[5pt]
\begin{tikzpicture}[baseline=-0.4ex]\draw[primaryBlue!70, line width=1pt, rounded corners=2pt] (0,0) rectangle (12pt,12pt);\end{tikzpicture} \;\textbf{Tôi có ô tô}
\end{browserWindow}

Dưới đây là bảng tổng hợp các thuộc tính phổ biến nhất được sử dụng cùng với thẻ Checkbox:

\begin{prettyTableBox}[Bảng Thuộc Tính Checkbox Cần Nắm Vững]
\rowcolors{2}{white}{primaryBlue!4}
\begin{tabularx}{\linewidth}{K{3.2cm}Z Z}
\rowcolor{primaryBlue}
\thd{Thuộc Tính} & \thd{Ý Nghĩa} & \thd{Ví Dụ Đầy Đủ} \\
\codeinline{type="checkbox"} & Xác định phần tử nhập liệu là một checkbox. & \codeinline{<input type="checkbox">} \\
\codeinline{name} & Tên trường dữ liệu, dùng để gom nhóm các checkbox liên quan. & \codeinline{<input type="checkbox" name="hobby">} \\
\codeinline{value} & Giá trị sẽ gửi đi lên máy chủ nếu checkbox được chọn. & \codeinline{<input type="checkbox" value="football">} \\
\codeinline{checked} & Đặt checkbox ở trạng thái mặc định được chọn sẵn khi tải trang. & \codeinline{<input type="checkbox" checked>} \\
\codeinline{required} & Bắt buộc phải chọn ô này trước khi gửi form. & \codeinline{<input type="checkbox" required>} \\
\codeinline{disabled} & Vô hiệu hóa ô chọn, không cho phép người dùng thay đổi trạng thái. & \codeinline{<input type="checkbox" disabled>} \\
\end{tabularx}
\end{prettyTableBox}

% ------------------------------------------------------------------------------
\subsection{Ví Dụ Đầy Đủ \& Cơ Chế Gửi Dữ Liệu Lên Server}

\begin{codeWindow}[Form chọn sở thích sử dụng Checkbox]
\begin{lstlisting}[language=HTML5]
<form action="/submit" method="GET">
  <p>(*@Sở thích của bạn:@*)</p>
  <input type="checkbox" name="hobby" value="music"> (*@Nghe nhạc@*)<br>
  <input type="checkbox" name="hobby" value="sport" checked> (*@Thể thao@*)<br>
  <input type="checkbox" name="hobby" value="travel"> (*@Du lịch@*)<br>
  <br>
  <input type="submit" value="(*@Gửi@*)">
</form>
\end{lstlisting}
\end{codeWindow}

\begin{browserWindow}[Kết Quả Hiển Thị Trình Duyệt: Form Chọn Sở Thích]
\sffamily\small
\textbf{Sở thích của bạn:} \\[5pt]
\begin{tikzpicture}[baseline=-0.4ex]\draw[primaryBlue!70, line width=1pt, rounded corners=2pt] (0,0) rectangle (12pt,12pt);\end{tikzpicture} \; Nghe nhạc \\[4pt]
\begin{tikzpicture}[baseline=-0.4ex]\fill[primaryBlue, rounded corners=2pt] (0,0) rectangle (12pt,12pt); \draw[white, line width=1.2pt] (3pt,6pt) -- (5pt,3pt) -- (9pt,9pt);\end{tikzpicture} \; \textbf{Thể thao} \quad \textcolor{codeComment}{\small\textit{(Checked mặc định)}} \\[4pt]
\begin{tikzpicture}[baseline=-0.4ex]\draw[primaryBlue!70, line width=1pt, rounded corners=2pt] (0,0) rectangle (12pt,12pt);\end{tikzpicture} \; Du lịch \\[8pt]
\tcbox[on line, arc=3pt, colback=primaryBlue, colframe=primaryBlue, boxsep=2pt, left=10pt, right=10pt, top=3pt, bottom=3pt]{\color{white}\bfseries\small Gửi}
\end{browserWindow}

\begin{tipBox}[Kết Quả Truy Vấn HTTP Request]
Nếu người dùng chọn đồng thời hai ô \textbf{Nghe nhạc} và \textbf{Thể thao}, dữ liệu URL Chuỗi truy vấn (Query String) gửi lên máy chủ sẽ là:
\begin{center}
\codeinline{hobby=music\&hobby=sport}
\end{center}
\textit{Giải thích}: Thuộc tính \codeinline{name} đóng vai trò là \textbf{Tên Biến (Key)}, còn thuộc tính \codeinline{value} đóng vai trò là \textbf{Giá Trị (Value)} trong cặp dữ liệu truyền đi \codeinline{key=value}.
\end{tipBox}

\paragraph{Lưu Ý Quan Trọng Khi Gửi Dữ Liệu:}
\begin{enumerate}[leftmargin=1.5em, itemsep=2pt, parsep=0pt]
    \item \textbf{Không chọn (Unchecked)}: Nếu checkbox không được tích chọn, trình duyệt sẽ \textbf{hoàn toàn không gửi bất kỳ dữ liệu nào} tương ứng với \codeinline{name} đó (Đây là cơ chế \term{Successful Control} theo chuẩn đặc tả W3C HTML5).
    \item \textbf{Cùng tên (\codeinline{name})}: Khi nhiều checkbox có cùng \codeinline{name}, khi submit trình duyệt sẽ gửi đi một danh sách các giá trị đã chọn dưới dạng dãy tham số trùng tên.
    \item \textbf{Checkbox độc lập}: Nếu muốn các ô chọn hoạt động hoàn toàn độc lập với nhau (thu nhận từng giá trị riêng biệt), hãy đặt các giá trị \codeinline{name} khác nhau.
\end{enumerate}

% ------------------------------------------------------------------------------
\subsection{So Sánh Checkbox name Giống Nhau (Gom Nhóm) vs name Khác Nhau (Độc Lập)}

Trong thiết kế biểu mẫu HTML, cách đặt tên thuộc tính \codeinline{name} quyết định trực tiếp bản chất cấu trúc dữ liệu mà máy chủ (Backend) thu nhận được.

\begin{prettyTableBox}[Ma Trận So Sánh thuộc tính name Giống Nhau vs Khác Nhau]
\rowcolors{2}{white}{primaryBlue!4}
\begin{tabularx}{\linewidth}{K{3.0cm} Z Z Z}
\rowcolor{primaryBlue}
\thd{Đặc Điểm So Sánh} & \thd{Cú Pháp HTML Mẫu} & \thd{Dữ Liệu Submit (Âm nhạc + Du lịch)} & \thd{Ngữ Cảnh Ứng Dụng} \\
\textbf{\codeinline{name} Giống Nhau}\newline(\term{Gom Nhóm}) & \codeinline{<input name="hobby">}\newline\textcolor{codeComment}{\small\textit{(Dùng chung name)}} & \codeinline{hobby=music\&}\newline\codeinline{hobby=travel}\newline\textcolor{codeComment}{\small\textit{(Gửi mảng 1 key)}} & Nhóm tùy chọn cùng chủ đề (sở thích, kỹ năng). Backend nhận mảng giá trị. \\
\textbf{\codeinline{name} Khác Nhau}\newline(\term{Lựa Chọn Độc Lập}) & \codeinline{<input name="music">}\newline\codeinline{<input name="travel">}\newline\textcolor{codeComment}{\small\textit{(Dùng name riêng)}} & \codeinline{music=yes\&}\newline\codeinline{travel=yes}\newline\textcolor{codeComment}{\small\textit{(Gửi từng key riêng)}} & Các cài đặt/tính năng độc lập (bật/tắt thông báo). Backend nhận từng key riêng. \\
\end{tabularx}
\end{prettyTableBox}

% ------------------------------------------------------------------------------
\paragraph{Trường Hợp 1: Gom Nhóm Checkbox Cùng Name (\codeinline{name="hobby"}):}

Khi các tùy chọn thuộc về cùng một danh mục (ví dụ: Danh sách sở thích), ta đặt chung \codeinline{name="hobby"}. Khi submit, trình duyệt sẽ gom tất cả các ô được chọn thành một danh sách gửi về máy chủ:

\begin{codeWindow}[Form gom nhóm Checkbox sử dụng chung name="hobby"]
\begin{lstlisting}[language=HTML5]
<form action="/submit" method="GET">
  <p>(*@Sở thích của bạn:@*)</p>
  <input type="checkbox" name="hobby" value="music" checked> (*@Âm nhạc@*)
  <input type="checkbox" name="hobby" value="sport"> (*@Thể thao@*)
  <input type="checkbox" name="hobby" value="travel" checked> (*@Du lịch@*)
  <input type="submit" value="(*@Gửi@*)">
</form>
\end{lstlisting}
\end{codeWindow}

\begin{browserWindow}[Mô Phỏng Trình Duyệt: Gom Nhóm Checkbox Chung name="hobby"]
\sffamily\small
\textbf{Sở thích của bạn:} \\[6pt]
\begin{tikzpicture}[baseline=-0.4ex]\fill[primaryBlue, rounded corners=2pt] (0,0) rectangle (12pt,12pt); \draw[white, line width=1.2pt] (3pt,6pt) -- (5pt,3pt) -- (9pt,9pt);\end{tikzpicture} \; \textbf{Âm nhạc} \quad
\begin{tikzpicture}[baseline=-0.4ex]\draw[primaryBlue!70, line width=1pt, rounded corners=2pt] (0,0) rectangle (12pt,12pt);\end{tikzpicture} \; Thể thao \quad
\begin{tikzpicture}[baseline=-0.4ex]\fill[primaryBlue, rounded corners=2pt] (0,0) rectangle (12pt,12pt); \draw[white, line width=1.2pt] (3pt,6pt) -- (5pt,3pt) -- (9pt,9pt);\end{tikzpicture} \; \textbf{Du lịch} \\[8pt]
\tcbox[on line, arc=3pt, colback=primaryBlue, colframe=primaryBlue, boxsep=2pt, left=10pt, right=10pt, top=3pt, bottom=3pt]{\color{white}\bfseries\small Gửi}
\end{browserWindow}

\begin{tipBox}[Kết Quả Chuỗi Truy Vấn Gửi Về Máy Chủ (Query String)]
Khi chọn đồng thời \textbf{Âm nhạc} và \textbf{Du lịch}, dữ liệu nhận được ở Backend là:
\begin{center}
\codeinline{hobby=music\&hobby=travel}
\end{center}
\end{tipBox}

% ------------------------------------------------------------------------------
\paragraph{Trường Hợp 2: Checkbox Độc Lập Khác Name (\codeinline{name} Riêng Biệt):}

Khi mỗi ô chọn đại diện cho một cài đặt/tính năng riêng biệt (như bật/tắt nhận tin tức, khuyến mãi, cập nhật), ta khai báo các thuộc tính \codeinline{name} riêng lẻ:

\begin{codeWindow}[Form tùy chọn Checkbox khai báo name riêng biệt]
\begin{lstlisting}[language=HTML5]
<form action="/submit" method="GET">
  <p>(*@Tùy chọn cài đặt hệ thống:@*)</p>
  <input type="checkbox" name="newsletter" value="yes" checked> (*@Nhận tin tức@*)
  <input type="checkbox" name="offers" value="yes"> (*@Nhận khuyến mãi@*)
  <input type="checkbox" name="updates" value="yes" checked> (*@Nhận cập nhật@*)
  <input type="submit" value="(*@Gửi@*)">
</form>
\end{lstlisting}
\end{codeWindow}

\begin{browserWindow}[Mô Phỏng Trình Duyệt: Tùy Chọn Độc Lập Khác name]
\sffamily\small
\textbf{Tùy chọn cài đặt hệ thống:} \\[6pt]
\begin{tikzpicture}[baseline=-0.4ex]\fill[primaryBlue, rounded corners=2pt] (0,0) rectangle (12pt,12pt); \draw[white, line width=1.2pt] (3pt,6pt) -- (5pt,3pt) -- (9pt,9pt);\end{tikzpicture} \; \textbf{Nhận tin tức} \quad
\begin{tikzpicture}[baseline=-0.4ex]\draw[primaryBlue!70, line width=1pt, rounded corners=2pt] (0,0) rectangle (12pt,12pt);\end{tikzpicture} \; Nhận khuyến mãi \quad
\begin{tikzpicture}[baseline=-0.4ex]\fill[primaryBlue, rounded corners=2pt] (0,0) rectangle (12pt,12pt); \draw[white, line width=1.2pt] (3pt,6pt) -- (5pt,3pt) -- (9pt,9pt);\end{tikzpicture} \; \textbf{Nhận cập nhật} \\[8pt]
\tcbox[on line, arc=3pt, colback=primaryBlue, colframe=primaryBlue, boxsep=2pt, left=10pt, right=10pt, top=3pt, bottom=3pt]{\color{white}\bfseries\small Gửi}
\end{browserWindow}

\begin{tipBox}[Kết Quả Chuỗi Truy Vấn Gửi Về Máy Chủ (Query String)]
Khi chọn đồng thời \textbf{Nhận tin tức} và \textbf{Nhận cập nhật}, dữ liệu nhận được ở Backend là:
\begin{center}
\codeinline{newsletter=yes\&updates=yes}
\end{center}
\textit{Ghi chú}: Tương tự nếu khai báo riêng tên từng sở thích (\codeinline{name="music"}, \codeinline{name="travel"}), dữ liệu gửi về sẽ là: \codeinline{music=yes\&travel=yes}.
\end{tipBox}

% ------------------------------------------------------------------------------
\begin{conceptBox}[Quy Tắc Quyết Định Chọn Thuộc Tính name Chuẩn Senior Frontend Engineer]
\begin{itemize}[leftmargin=1.5em, itemsep=3pt, parsep=0pt]
    \item \textbf{Sử dụng name GIỐNG NHAU} khi các ô chọn thuộc về cùng 1 tập hợp danh mục (sở thích, danh sách kỹ năng, danh mục sản phẩm). Backend sẽ xử lý dữ liệu dưới dạng mảng (\term{Array}).
    \item \textbf{Sử dụng name KHÁC NHAU} khi mỗi ô chọn đại diện cho một cờ bật/tắt (\term{Boolean Flag}) độc lập (như đồng ý điều khoản, đăng ký bản tin, bật thông báo qua Email).
\end{itemize}
\end{conceptBox}

% ------------------------------------------------------------------------------
\subsection{Tối Ưu Trải Nghiệm Người Dùng Với Thẻ <label>}

Nếu không có thẻ \codeinline{<label>}, người dùng vẫn có thể chọn checkbox bằng cách nhấp trực tiếp vào ô vuông nhỏ. Tuy nhiên, việc không sử dụng label sẽ mất đi các lợi ích quan trọng sau:
\begin{itemize}[leftmargin=1.5em, itemsep=2pt, parsep=0pt]
    \item \textbf{Dễ dàng thao tác}: Khi có label, người dùng có thể nhấp vào cả dòng văn bản để bật/tắt checkbox, thay vì phải căn chỉnh con trỏ vào ô vuông nhỏ.
    \item \textbf{Khắc phục khó khăn trên di động}: Giúp tương tác thuận tiện hơn khi kích thước checkbox nhỏ hoặc nằm gần các thành phần khác.
    \item \textbf{Tạo sự liên kết logic}: Khi sử dụng \codeinline{label for="id"}, nhãn được gắn kết trực tiếp với checkbox thông qua thuộc tính \codeinline{id}. Ngược lại nếu không có label, văn bản chỉ là text đơn thuần không có liên kết logic.
\end{itemize}

\begin{codeWindow}[Ví dụ hoàn chỉnh kết hợp Checkbox với thẻ label và thuộc tính action]
\begin{lstlisting}[language=HTML5]
<form action="https://www.w3schools.com/action_page.php" target="_blank">
    <p>(*@Những ngôn ngữ lập trình bạn đã học:@*)</p>
    <input type="checkbox" id="ngonNguC" name="language" value="c">
    <label for="ngonNguC">(*@Ngôn ngữ C@*)</label><br>
    
    <input type="checkbox" id="ngonNguJava" name="language" value="java">
    <label for="ngonNguJava">(*@Ngôn ngữ Java@*)</label><br>
    
    <input type="checkbox" id="ngonNguPhp" name="language" value="php">
    <label for="ngonNguPhp">(*@Ngôn ngữ PHP@*)</label><br>
    
    <input type="submit" value="(*@Gửi@*)">
</form>
\end{lstlisting}
\end{codeWindow}

\begin{browserWindow}[Kết Quả Hiển Thị Trình Duyệt: Tương Tác Chọn Ngôn Ngữ Qua Thẻ Label]
\sffamily\small
\textbf{Những ngôn ngữ lập trình bạn đã học:} \\[5pt]
\begin{tikzpicture}[baseline=-0.4ex]\fill[primaryBlue, rounded corners=2pt] (0,0) rectangle (12pt,12pt); \draw[white, line width=1.2pt] (3pt,6pt) -- (5pt,3pt) -- (9pt,9pt);\end{tikzpicture} \; \textbf{Ngôn ngữ C} \quad \textcolor{primaryBlue}{\small\textbf{$\leftarrow$ Nhấp vào chữ "Ngôn ngữ C" cũng kích hoạt ô chọn}} \\[4pt]
\begin{tikzpicture}[baseline=-0.4ex]\draw[primaryBlue!70, line width=1pt, rounded corners=2pt] (0,0) rectangle (12pt,12pt);\end{tikzpicture} \; Ngôn ngữ Java \\[4pt]
\begin{tikzpicture}[baseline=-0.4ex]\draw[primaryBlue!70, line width=1pt, rounded corners=2pt] (0,0) rectangle (12pt,12pt);\end{tikzpicture} \; Ngôn ngữ PHP \\[8pt]
\tcbox[on line, arc=3pt, colback=primaryBlue, colframe=primaryBlue, boxsep=2pt, left=10pt, right=10pt, top=3pt, bottom=3pt]{\color{white}\bfseries\small Gửi}
\end{browserWindow}

\begin{noteBox}[Giải Thích Cấu Trúc Mã Nguồn]
\begin{itemize}[leftmargin=1.5em, itemsep=2pt, parsep=0pt]
    \item \codeinline{action="https://www.w3schools.com/action\_page.php"}: Khi gửi form, dữ liệu sẽ được chuyển tiếp đến trang này để xử lý.
    \item \codeinline{target="\_blank"}: Kết quả xử lý form sẽ được mở ra trong một tab mới trên trình duyệt.
    \item \codeinline{name="language"}: Gom nhóm các ô chọn ngôn ngữ lập trình lại với nhau.
    \item \codeinline{value="c|java|php"}: Khi tích chọn ô nào, giá trị tương ứng của ô đó sẽ được gửi đi. Ví dụ: nhấp chọn ô Ngôn ngữ C $\rightarrow$ tương ứng với lựa chọn \codeinline{language=c}.
    \item \codeinline{label for="id"}: Liên kết nhãn văn bản với thẻ input tương ứng qua ID.
\end{itemize}
\end{noteBox}

% ------------------------------------------------------------------------------
\subsection{Các Kỹ Thuật Sử Dụng Checkbox Nâng Cao}

\paragraph{Checkbox Bắt Buộc Chọn (\codeinline{required}):}
Thêm thuộc tính \codeinline{required} để bắt buộc người dùng phải tích chọn trước khi cho phép submit form (thường áp dụng cho ô Đồng ý điều khoản):

\begin{codeWindow}[Bắt buộc chọn checkbox điều khoản]
\begin{lstlisting}[language=HTML5]
<form>
    <label>
        <input type="checkbox" name="agree" required> 
        (*@Tôi đồng ý với điều khoản@*)
    </label>
    <input type="submit" value="(*@Gửi@*)">
</form>
\end{lstlisting}
\end{codeWindow}

\begin{browserWindow}[Mô Phỏng Trình Duyệt: Checkbox Bắt Buộc Chọn (Required Validation)]
\sffamily\small
\begin{tikzpicture}[baseline=-0.4ex]\draw[accentOrange, line width=1.2pt, rounded corners=2pt] (0,0) rectangle (12pt,12pt);\end{tikzpicture} \; \textbf{Tôi đồng ý với điều khoản} \quad \textcolor{accentOrange}{\small\textbf{* (Bắt buộc)}} \\[6pt]
\tcbox[on line, arc=3pt, colback=primaryBlue, colframe=primaryBlue, boxsep=2pt, left=10pt, right=10pt, top=3pt, bottom=3pt]{\color{white}\bfseries\small Gửi}
\end{browserWindow}

\paragraph{Checkbox Mặc Định Được Chọn (\codeinline{checked}):}
Sử dụng thuộc tính \codeinline{checked} để ô chọn luôn ở trạng thái đã được tích sẵn khi trang web tải xong:

\begin{codeWindow}[Khởi tạo checkbox mặc định checked]
\begin{lstlisting}[language=HTML5]
<input type="checkbox" name="subscribe" value="yes" checked> (*@Nhận tin tức@*)
\end{lstlisting}
\end{codeWindow}

\paragraph{Kỹ Thuật Checkbox Ẩn (Hidden Checkbox Fallback):}
Mặc định trong chuẩn HTML, nếu checkbox không được tích chọn, trình duyệt sẽ \textbf{hoàn toàn không gửi bất kỳ dữ liệu nào} lên máy chủ. Việc này gây khó khăn cho hệ thống backend khi cần ghi nhận giá trị \codeinline{false} hoặc \codeinline{"no"}.

\textbf{Giải pháp}: Sử dụng một thẻ \codeinline{<input type="hidden">} đặt nằm \textbf{ngay trước} thẻ checkbox có cùng thuộc tính \codeinline{name}.

\begin{codeWindow}[Kỹ thuật gửi giá trị mặc định khi không chọn checkbox]
\begin{lstlisting}[language=HTML5]
<!-- (*@Thẻ Hidden chứa giá trị mặc định khi KHÔNG chọn@*) -->
<input type="hidden" name="subscribe" value="no">

<!-- (*@Thẻ Checkbox chứa giá trị khi ĐƯỢC chọn@*) -->
<input type="checkbox" name="subscribe" value="yes"> (*@Nhận tin tức@*)
\end{lstlisting}
\end{codeWindow}

\begin{conceptBox}[Nguyên Lý Hoạt Động Của Hidden Fallback]
\begin{itemize}[leftmargin=1.5em, itemsep=2pt, parsep=0pt]
    \item \textbf{Khi KHÔNG chọn Checkbox}: Trình duyệt bỏ qua ô checkbox và gửi giá trị của ô Hidden $\rightarrow$ Dữ liệu nhận được: \codeinline{subscribe=no}.
    \item \textbf{Khi CÓ chọn Checkbox}: Cả hai thẻ đều được gửi đi. Do ô Checkbox nằm ở sau, giá trị \codeinline{subscribe=yes} của nó sẽ đè lên (\term{override}) giá trị \codeinline{no} đứng trước $\rightarrow$ Dữ liệu nhận được: \codeinline{subscribe=yes}.
\end{itemize}
\end{conceptBox}

% ------------------------------------------------------------------------------
\subsection{Phân Biệt HTML Attribute checked vs DOM Property}

\begin{prettyTableBox}[So Sánh HTML Attribute vs DOM Property]
\rowcolors{2}{white}{primaryBlue!4}
\begin{tabularx}{\linewidth}{K{3.5cm}Z Z}
\rowcolor{primaryBlue}
\thd{Đặc Điểm} & \thd{HTML Attribute (\codeinline{checked})} & \thd{DOM Property (\codeinline{element.checked})} \\
Bản chất & Xác định trạng thái ban đầu khi tải trang. & Phản ánh trạng thái thời gian thực (Real-time). \\
Tương tác người dùng & \textbf{KHÔNG thay đổi} khi nhấp chọn. & \textbf{TỰ ĐỘNG thay đổi} (\codeinline{true} / \codeinline{false}). \\
Truy cập JavaScript & \codeinline{elem.getAttribute('checked')} & \codeinline{elem.checked} hoặc \codeinline{elem.defaultChecked} \\
Khôi phục Form & Dùng để Reset form về ban đầu. & Dùng để kiểm tra trạng thái trước khi xử lý Logic. \\
\end{tabularx}
\end{prettyTableBox}

% ------------------------------------------------------------------------------
\subsection{Bắt Buộc Chọn Ít Nhất 1 Checkbox Trong Nhóm}

\begin{warnBox}[Hạn Chế Của Thuộc Tính required Mặc Định]
Khi gán \codeinline{required} cho nhiều checkbox có cùng \codeinline{name}, chuẩn HTML5 yêu cầu người dùng phải tích chọn \textbf{TẤT CẢ các ô đó} thay vì chọn ít nhất 1 ô. 
\end{warnBox}

\textbf{Giải pháp}: Sử dụng JavaScript Constraint Validation API (\codeinline{setCustomValidity}) để kiểm tra cả nhóm:

\begin{codeWindow}[Giải pháp Validate bắt buộc chọn ít nhất 1 Checkbox trong nhóm]
\begin{lstlisting}[language=JavaScript]
function validateCheckboxGroup() {
  const checkboxes = document.querySelectorAll('input[name="hobby"]');
  // (*@Kiểm tra xem có ít nhất 1 ô nào được checked hay không@*)
  const isOneChecked = Array.from(checkboxes).some(cb => cb.checked);

  checkboxes.forEach(cb => {
    if (!isOneChecked) {
      // (*@Gán thông báo lỗi nếu chưa chọn ô nào@*)
      cb.setCustomValidity('(*@Vui lòng chọn ít nhất một sở thích!@*)');
    } else {
      // (*@Xóa thông báo lỗi khi đã chọn hợp lệ@*)
      cb.setCustomValidity('');
    }
  });
}
\end{lstlisting}
\end{codeWindow}

% ------------------------------------------------------------------------------
\subsection{Thu Thập Dữ Liệu Trong JavaScript Modern \& React}

\paragraph{Lấy Toàn Bộ Mảng Checkbox Được Chọn Với FormData:}
API \codeinline{FormData.getAll()} trong ES6+ giúp thu thập danh sách cực kỳ gọn gàng:

\begin{codeWindow}[Thu thập mảng dữ liệu Checkbox với FormData]
\begin{lstlisting}[language=JavaScript]
const form = document.querySelector('form');
form.addEventListener('submit', (e) => {
  e.preventDefault();
  const formData = new FormData(form);
  
  // (*@Trả về mảng chứa tất cả các value được chọn của 'hobby'@*)
  const selectedHobbies = formData.getAll('hobby'); 
  console.log('(*@Các sở thích đã chọn:@*)', selectedHobbies);
});
\end{lstlisting}
\end{codeWindow}

\paragraph{Quản Lý State Checkbox Trong React (Controlled Component):}
Trong các Framework như React, trạng thái của danh sách Checkbox thường được quản lý dưới dạng mảng State:

\begin{codeWindow}[Xử lý nhóm Checkbox trong React Component]
\begin{lstlisting}[language=JavaScript]
import React, { useState } from 'react';

function HobbySelector() {
  const [selectedHobbies, setSelectedHobbies] = useState([]);

  const handleCheckboxChange = (event) => {
    const { value, checked } = event.target;
    if (checked) {
      // (*@Thêm giá trị vào mảng state nếu được chọn@*)
      setSelectedHobbies([...selectedHobbies, value]);
    } else {
      // (*@Lọc bỏ giá trị khỏi mảng state nếu bỏ chọn@*)
      setSelectedHobbies(selectedHobbies.filter(item => item !== value;
    }
  };

  return (
    <div>
      <label>
        <input type="checkbox" value="music" onChange={handleCheckboxChange} /> (*@Nghe nhạc@*)
      </label>
      <label>
        <input type="checkbox" value="sport" onChange={handleCheckboxChange} /> (*@Thể thao@*)
      </label>
    </div>
  );
}
\end{lstlisting}
\end{codeWindow}

% ------------------------------------------------------------------------------
\subsection{Cảnh Báo Bảo Mật (Security Notes)}

\begin{warnBox}[Cảnh Báo Bảo Mật Dữ Liệu Checkbox]
\begin{itemize}[leftmargin=1.5em, itemsep=2pt, parsep=0pt]
    \item \textbf{Không bao giờ tin tưởng giá trị \codeinline{value} gửi từ Client}: Người dùng có thể dễ dàng mở DevTools trên trình duyệt để sửa đổi thuộc tính \codeinline{value="free"} thành \codeinline{value="admin"} trước khi bấm Submit.
    \item \textbf{Khử trùng dữ liệu (Sanitization \& Validation ở Backend)}: Máy chủ (Backend) luôn phải kiểm tra danh sách giá trị nhận được có nằm trong tập danh mục cho phép (\term{Whitelist}) hay không trước khi lưu vào cơ sở dữ liệu.
\end{itemize}
\end{warnBox}

% ------------------------------------------------------------------------------
\subsection{Minh Họa Sơ Đồ Luồng Xử Lý Dữ Liệu Checkbox Khi Submit}

\begin{figure}[H]
\centering
\begin{tikzpicture}[
    auto, 
    >=stealth, 
    semithick,
    node distance=2.2cm,
    startstop/.style={rectangle, rounded corners=5pt, minimum width=3.8cm, minimum height=1cm, text centered, draw=primaryBlue, fill=primaryBlue!10, font=\sffamily\bfseries\small},
    decision/.style={diamond, aspect=2.2, minimum width=4cm, minimum height=1.1cm, text centered, draw=accentOrange, fill=accentOrange!10, font=\sffamily\bfseries\small},
    process/.style={rectangle, rounded corners=4pt, minimum width=3.8cm, minimum height=1cm, text centered, draw=secondaryTeal, fill=secondaryTeal!10, font=\sffamily\small}
]
    \node (start) [startstop] {1. Bấm Submit Form};
    \node (scan) [process, below of=start, node distance=1.8cm] {2. Quét ô \codeinline{<input type="checkbox">}};
    \node (check) [decision, below of=scan, node distance=2.2cm] {Tích chọn? (\codeinline{checked})};
    
    \node (send) [process, right=2.8cm of check] {Đóng gói: \codeinline{name=value}};
    \node (ignore) [process, left=2.8cm of check] {Bỏ qua, không gửi};
    
    \node (end) [startstop, below of=check, node distance=2.6cm] {3. Gửi HTTP Request};

    \draw [->] (start) -- (scan);
    \draw [->] (scan) -- (check);
    \draw [->] (check) -- node [above=2pt, font=\sffamily\bfseries\scriptsize, fill=white, inner sep=2pt] {\color{primaryBlue}CÓ (true)} (send);
    \draw [->] (check) -- node [above=2pt, font=\sffamily\bfseries\scriptsize, fill=white, inner sep=2pt] {\color{accentOrange}KHÔNG (false)} (ignore);
    \draw [->] (send) |- (end);
    \draw [->] (ignore) |- (end);
\end{tikzpicture}
\caption{Sơ đồ quy trình đóng gói dữ liệu Checkbox của trình duyệt khi submit Form}
\label{fig:checkbox_flowchart}
\end{figure}

% ------------------------------------------------------------------------------
\subsection{Bảng Tổng Kết \& Ghi Nhớ Cốt Lõi}

\begin{prettyTableBox}[Tóm Tắt Các Thuộc Tính Và Kỹ Thuật Checkbox]
\rowcolors{2}{white}{primaryBlue!4}
\begin{tabularx}{\linewidth}{K{2.8cm}Z Z}
\rowcolor{primaryBlue}
\thd{Thuộc Tính / Thẻ} & \thd{Mô Tả} & \thd{Ghi Chú Thực Hành} \\
\codeinline{checked} & Mặc định được chọn khi tải trang. & Giúp tăng tỷ lệ chuyển đổi cho lựa chọn ưu tiên. \\
\codeinline{required} & Bắt buộc phải chọn ô này mới được submit. & Dùng khi yêu cầu đồng ý điều khoản dịch vụ. \\
\codeinline{name} & Tên của trường dữ liệu gửi đi. & Thuộc tính quan trọng nhất để server nhận diện. \\
\codeinline{value} & Giá trị gửi lên máy chủ khi ô được chọn. & Nếu không khai báo value, trình duyệt mặc định gửi "on". \\
\codeinline{hidden input} & Thẻ ẩn gửi giá trị mặc định khi không chọn. & Tránh bị mất dữ liệu hay lỗi khuyết key ở Backend. \\
\end{tabularx}
\end{prettyTableBox}

\begin{tipBox}[Những Điểm Bắt Buộc Phải Ghi Nhớ]
\begin{enumerate}[leftmargin=1.5em, itemsep=2pt, parsep=0pt]
    \item Checkbox \textbf{không chọn thì không gửi dữ liệu} $\rightarrow$ Sử dụng \codeinline{<input type="hidden">} để khắc phục.
    \item Các checkbox có \textbf{cùng name} khi submit sẽ gửi dữ liệu dưới dạng một danh sách giá trị (dãy các tham số trùng tên).
    \item Luôn phải khai báo thuộc tính \codeinline{name} để máy chủ thu nhận được dữ liệu.
    \item Nên luôn bọc hoặc kết hợp với thẻ \codeinline{<label>} để tăng tính thân thiện và trải nghiệm người dùng.
\end{enumerate}
\end{tipBox}

% ------------------------------------------------------------------------------
\subsection{Đọc Thêm (Mở Rộng Kiến Thức Chuyên Sâu)}
\label{subsec:checkbox_advanced_readmore}

Để nâng cao năng lực thiết kế và phát triển ứng dụng web thực tế, lập trình viên Frontend cần nắm vững các khía cạnh chuyên sâu sau đây liên quan đến Checkbox:

\paragraph{a. Trạng Thái Không Xác Định (Indeterminate State) Trong JavaScript:}
Ngoài hai trạng thái hiển thị cơ bản là \codeinline{checked = true} và \codeinline{checked = false}, HTML Checkbox còn sở hữu trạng thái thứ 3 gọi là \term{Indeterminate} (Lơ lửng / Không xác định). Trạng thái này thường xuất hiện trong giao diện danh sách cây (Tree-view) hoặc ô "Chọn tất cả" (\term{Select All}) khi chỉ có một số ô con được chọn.

\begin{codeWindow}[Thiết lập trạng thái Indeterminate bằng JavaScript]
\begin{lstlisting}[language=JavaScript]
const parentCheckbox = document.getElementById('selectAll');
// (*@Trạng thái này KHÔNG THỂ thiết lập trực tiếp bằng thuộc tính HTML@*)
parentCheckbox.indeterminate = true; // (*@Hiển thị dấu gạch ngang (-) trên ô chọn@*)
\end{lstlisting}
\end{codeWindow}

\begin{browserWindow}[Mô Phỏng Trình Duyệt: Trạng Thái Indeterminate (Lơ lửng)]
\sffamily\small
\begin{tikzpicture}[baseline=-0.4ex]\fill[primaryBlue, rounded corners=2pt] (0,0) rectangle (12pt,12pt); \draw[white, line width=1.5pt] (3pt,6pt) -- (9pt,6pt);\end{tikzpicture} \; \textbf{Chọn tất cả danh mục} \quad \textcolor{accentOrange}{\small\textbf{$\leftarrow$ Trạng thái Indeterminate (Dấu gạch ngang)}} \\[5pt]
\quad \begin{tikzpicture}[baseline=-0.4ex]\fill[primaryBlue, rounded corners=2pt] (0,0) rectangle (12pt,12pt); \draw[white, line width=1.2pt] (3pt,6pt) -- (5pt,3pt) -- (9pt,9pt);\end{tikzpicture} \; Danh mục 1 (Đã chọn) \\[4pt]
\quad \begin{tikzpicture}[baseline=-0.4ex]\draw[primaryBlue!70, line width=1pt, rounded corners=2pt] (0,0) rectangle (12pt,12pt);\end{tikzpicture} \; Danh mục 2 (Chưa chọn)
\end{browserWindow}

\paragraph{b. Gom Nhóm Truy Cập Chuẩn Accessibility Với <fieldset> Và <legend>:}
Khi có một nhóm nhiều Checkbox liên quan đến cùng một câu hỏi (như danh sách sở thích), việc gom nhóm bằng các thẻ chuẩn Accessibility là yêu cầu bắt buộc đối với Trình đọc màn hình (Screen Reader):

\begin{codeWindow}[Gom nhóm Checkbox hợp chuẩn WCAG/a11y]
\begin{lstlisting}[language=HTML5]
<fieldset>
  <legend>(*@Chọn các kỹ năng Frontend bạn làm chủ:@*)</legend>
  <label><input type="checkbox" name="skill" value="html"> HTML5</label>
  <label><input type="checkbox" name="skill" value="css"> CSS3</label>
  <label><input type="checkbox" name="skill" value="js"> JavaScript</label>
</fieldset>
\end{lstlisting}
\end{codeWindow}

\begin{browserWindow}[Mô Phỏng Trình Duyệt: Khung Gom Nhóm Thẻ <fieldset> \& <legend>]
\sffamily\small
\textbf{Chọn các kỹ năng Frontend bạn làm chủ:} \textcolor{codeComment}{\small\textit{(Tiêu đề <legend>)}} \\[6pt]
\begin{tikzpicture}[baseline=-0.4ex]\draw[primaryBlue!70, line width=1pt, rounded corners=2pt] (0,0) rectangle (12pt,12pt);\end{tikzpicture} \; HTML5 \qquad
\begin{tikzpicture}[baseline=-0.4ex]\fill[primaryBlue, rounded corners=2pt] (0,0) rectangle (12pt,12pt); \draw[white, line width=1.2pt] (3pt,6pt) -- (5pt,3pt) -- (9pt,9pt);\end{tikzpicture} \; \textbf{CSS3} \qquad
\begin{tikzpicture}[baseline=-0.4ex]\fill[primaryBlue, rounded corners=2pt] (0,0) rectangle (12pt,12pt); \draw[white, line width=1.2pt] (3pt,6pt) -- (5pt,3pt) -- (9pt,9pt);\end{tikzpicture} \; \textbf{JavaScript}
\end{browserWindow}

\paragraph{c. Tùy Biến Giao Diện Checkbox Bằng CSS Modern (\codeinline{appearance: none}):}
Các ô chọn mặc định của trình duyệt thường có giao diện thô cứng và không đồng nhất giữa các hệ điều hành (Windows, macOS, iOS). Kỹ thuật CSS hiện đại cho phép tái thiết kế giao diện Checkbox hoàn toàn theo thiết kế Figma:

\begin{codeWindow}[Custom giao diện Checkbox đẹp mắt bằng CSS3]
\begin{lstlisting}[language=HTML5]
input[type="checkbox"] {
  appearance: none; /* (*@Bỏ giao diện mặc định của hệ điều hành@*) */
  width: 20px;
  height: 20px;
  border: 2px solid #0056b3;
  border-radius: 4px;
  cursor: pointer;
  transition: all 0.2s ease-in-out;
}

input[type="checkbox"]:checked {
  background-color: #0056b3;
  position: relative;
}

/* (*@Tạo dấu tích chọn (Checkmark) bằng Pseudo-element ::after@*) */
input[type="checkbox"]:checked::after {
  content: '';
  position: absolute;
  left: 6px; top: 2px;
  width: 5px; height: 10px;
  border: solid white;
  border-width: 0 2px 2px 0;
  transform: rotate(45deg);
}
\end{lstlisting}
\end{codeWindow}

\begin{browserWindow}[Mô Phỏng Trình Duyệt: Giao Diện Custom Checkbox Bằng CSS Modern]
\sffamily\small
\begin{tikzpicture}[baseline=-0.4ex]\fill[primaryBlue, rounded corners=3pt] (0,0) rectangle (14pt,14pt); \draw[white, line width=1.5pt] (3.5pt,7pt) -- (6pt,4.5pt) -- (10.5pt,10.5pt);\end{tikzpicture} \;\textbf{Checkbox Đã Được Styling Với Bo Góc 4pt \& Màu Brand Primary Blue} \\[5pt]
\begin{tikzpicture}[baseline=-0.4ex]\draw[primaryBlue, line width=1.2pt, rounded corners=3pt] (0,0) rectangle (14pt,14pt);\end{tikzpicture} \; Checkbox Chưa Chọn Với Viền Xanh Đẹp Mắt
\end{browserWindow}

\paragraph{d. Xử Lý Mảng Checkbox Ở Backend (PHP, Node.js, Express):}
Trong các ngôn ngữ server-side như PHP, nếu muốn tự động nhận nhiều checkbox gửi lên dưới dạng một mảng (Array) chính thức, cú pháp đặt tên \codeinline{name} cần có thêm ký tự ngoặc vuông \codeinline{[]}:

\begin{codeWindow}[Đặt name dạng mảng cho PHP Backend]
\begin{lstlisting}[language=HTML5]
<input type="checkbox" name="hobby[]" value="music"> (*@Nghe nhạc@*)
<input type="checkbox" name="hobby[]" value="sport"> (*@Thể thao@*)
<!-- (*@Trên PHP Backend sẽ thu nhận được@*) $_POST['hobby'] (*@dưới dạng Array: ['music', 'sport']@*) -->
\end{lstlisting}
\end{codeWindow}

% ------------------------------------------------------------------------------
\subsection{Bộ Câu Hỏi Phỏng Vấn Frontend Thực Chiến (Interview Q\&A)}

\begin{noteBox}[Câu Hỏi Phỏng Vấn 1: Phân Biệt Checkbox Và Radio Button]
\textbf{Hỏi}: Khi nào nên sử dụng Checkbox và khi nào nên sử dụng Radio Button trong thiết kế giao diện Form?
\par\vspace{3pt}
\textbf{Trả lời}:
\begin{itemize}[leftmargin=1.5em, itemsep=2pt, parsep=0pt]
    \item \textbf{Checkbox}: Dùng cho bài toán chọn nhiều lựa chọn độc lập (\term{Multi-choice}, chọn 0, 1 hoặc nhiều) hoặc trường hợp bật/tắt một tính năng đơn lẻ (\term{Toggle switch}).
    \item \textbf{Radio Button}: Dùng cho bài toán bắt buộc chọn duy nhất 1 trong danh sách các lựa chọn loại trừ lẫn nhau (\term{Single-choice}).
\end{itemize}
\end{noteBox}

\begin{noteBox}[Câu Hỏi Phỏng Vấn 2: Xử Lý Vấn Đề Khuyết Dữ Liệu Unchecked]
\textbf{Hỏi}: Tại sao ô Checkbox không tích chọn lại không gửi dữ liệu về máy chủ? Làm sao để khắc phục ở Frontend?
\par\vspace{3pt}
\textbf{Trả lời}: Theo đặc tả kỹ thuật HTML, trình duyệt chỉ đóng gói các phần tử ở trạng thái \term{successful control} (có \codeinline{name} và được \codeinline{checked}). Để khắc phục, ta đặt một thẻ \codeinline{<input type="hidden">} có cùng \codeinline{name} với giá trị mặc định (ví dụ \codeinline{value="no"}) ngay trước thẻ checkbox.
\end{noteBox}

\begin{noteBox}[Câu Hỏi Phỏng Vấn 3: Thiết Lập Trạng Thái Indeterminate]
\textbf{Hỏi}: Trạng thái lơ lửng (\term{Indeterminate}) của Checkbox có thể khai báo trực tiếp bằng thuộc tính HTML được không?
\par\vspace{3pt}
\textbf{Trả lời}: KHÔNG. Trạng thái \term{Indeterminate} không có thuộc tính HTML tương ứng mà bắt buộc phải thiết lập thông qua thuộc tính DOM JavaScript: \codeinline{checkboxElement.indeterminate = true}.
\end{noteBox}

% ------------------------------------------------------------------------------
\subsection{Bảng Độ Tương Thích Trình Duyệt \& Best Practices Checklist}

\begin{prettyTableBox}[Bảng Độ Tương Thích Trình Duyệt Của Các Tính Năng Checkbox]
\rowcolors{2}{white}{primaryBlue!4}
\begin{tabularx}{\linewidth}{K{3.2cm}C{2.2cm}C{2.2cm}C{2.2cm}Z}
\rowcolor{primaryBlue}
\thd{Tính Năng} & \thd{Chrome / Edge} & \thd{Firefox} & \thd{Safari (iOS)} & \thd{Ghi Chú Hỗ Trợ} \\
\codeinline{<input type="checkbox">} & Full (v1+) & Full (v1+) & Full (v1+) & Hỗ trợ 100\% trên tất cả trình duyệt. \\
\codeinline{elem.indeterminate} & Full (v1+) & Full (v1+) & Full (v1+) & Thuộc tính DOM JS chuẩn W3C. \\
CSS \codeinline{appearance: none} & Full (v84+) & Full (v80+) & Full (v15.4+) & Tùy biến giao diện không cần prefix. \\
\codeinline{FormData.getAll()} & Full (v50+) & Full (v44+) & Full (v11.1+) & Gom mảng value checkbox được chọn. \\
\end{tabularx}
\end{prettyTableBox}

\begin{tipBox}[Checklist 5 Quy Tắc Vàng Khi Lập Trình Checkbox]
\begin{enumerate}[leftmargin=1.5em, itemsep=2pt, parsep=0pt]
    \item \textbf{Luôn kết hợp với thẻ \codeinline{<label>}}: Tăng diện tích tương tác click cho người dùng trên thiết bị di động.
    \item \textbf{Khai báo thuộc tính \codeinline{name} chính xác}: Bắt buộc phải có \codeinline{name} để máy chủ nhận dạng được trường dữ liệu.
    \item \textbf{Cung cấp \codeinline{value} có nghĩa}: Đặt giá trị \codeinline{value} rõ ràng (ví dụ \codeinline{value="email"} thay vì bỏ trống để mặc định gửi \codeinline{"on"}).
    \item \textbf{Gom nhóm với \codeinline{<fieldset>} \& \codeinline{<legend>}}: Đảm bảo tính khả truy cập (Accessibility / a11y) đạt chuẩn WCAG 2.1 cho Trình đọc màn hình.
    \item \textbf{Luôn Validate \& Sanitization ở Backend}: Không tin tưởng giá trị \codeinline{value} gửi từ Client, kiểm tra Whitelist trên Server.
\end{enumerate}
\end{tipBox}

% ==============================================================================
% MỤC GIÁO TRÌNH: CHI TIẾT VỀ RADIO BUTTON trong HTML (<input type="radio">)
% ==============================================================================
\section[Thẻ Nút Chọn Đơn Radio Button]{Thẻ Nút Chọn Đơn Radio Button (<input type="radio">)}
\label{sec:html_radio_button}

\begin{conceptBox}[Định Nghĩa Radio Button]
\term{Radio Button} là một loại phần tử nhập liệu (\codeinline{<input>}) dạng ô chọn đơn (\term{single-choice}), cho phép người dùng chọn \textbf{duy nhất 1 lựa chọn} trong một nhóm các tùy chọn loại trừ lẫn nhau (\term{mutually exclusive}).
\begin{itemize}[leftmargin=1.5em, itemsep=2pt, parsep=0pt]
    \item \textbf{Cùng \codeinline{name} gom nhóm}: Trong một nhóm các nút Radio có cùng thuộc tính \codeinline{name}, khi người dùng tích chọn một ô mới, trình duyệt sẽ tự động bỏ chọn (\term{uncheck}) ô đã chọn trước đó.
    \item \textbf{Khác biệt với Checkbox}: Checkbox cho phép chọn nhiều lựa chọn cùng lúc (\term{multi-choice}), trong khi Radio Button bắt buộc chọn duy nhất một giá trị.
    \item \textbf{Ứng dụng thực tế}: Chọn giới tính (Nam/Nữ/Khác), Phương thức thanh toán (Thẻ tín dụng/PayPal/COD), Hình thức vận chuyển (Giao hàng nhanh/Tiết kiệm), Mức độ hài lòng,...
\end{itemize}
\end{conceptBox}

% ------------------------------------------------------------------------------
\subsection{Cú Pháp Cơ Bản \& Bảng Thuộc Tính Thường Gặp}

\begin{codeWindow}[Cú pháp tạo Radio Button cơ bản]
\begin{lstlisting}[language=HTML5]
<input type="radio" name="gender" value="male"> Nam
<input type="radio" name="gender" value="female"> (*@Nữ@*)
\end{lstlisting}
\end{codeWindow}

\begin{browserWindow}[Kết Quả Hiển Thị Trình Duyệt: Radio Button Cơ Bản]
\sffamily\small
\begin{tikzpicture}[baseline=-0.4ex]\fill[primaryBlue] (0,0) circle (6pt); \fill[white] (0,0) circle (2.2pt);\end{tikzpicture} \;\textbf{Nam} \quad \textcolor{codeComment}{\small\textit{(Được chọn mặc định)}} \qquad
\begin{tikzpicture}[baseline=-0.4ex]\draw[primaryBlue!70, line width=1pt] (0,0) circle (6pt);\end{tikzpicture} \; Nữ
\end{browserWindow}

\begin{tipBox}[Giải Thích Cú Pháp Gom Nhóm]
Vì cả hai ô chọn trên đều có cùng thuộc tính \codeinline{name="gender"}, trình duyệt hiểu rằng chúng thuộc về cùng một nhóm câu hỏi $\rightarrow$ Người dùng chỉ có thể nhấp chọn \textbf{duy nhất 1 trong 2} tùy chọn.
\end{tipBox}

Dưới đây là bảng tổng hợp các thuộc tính phổ biến nhất được sử dụng với Radio Button:

\begin{prettyTableBox}[Bảng Thuộc Tính Radio Button Cần Nắm Vững]
\rowcolors{2}{white}{primaryBlue!4}
\begin{tabularx}{\linewidth}{K{3.0cm} Z Z}
\rowcolor{primaryBlue}
\thd{Thuộc Tính} & \thd{Ý Nghĩa Kỹ Thuật} & \thd{Ví Dụ Đầy Đủ} \\
\codeinline{type="radio"} & Xác định phần tử nhập liệu là một nút Radio Button. & \codeinline{<input type="radio">} \\
\codeinline{name} & Gom nhóm các nút Radio với nhau (các ô cùng nhóm bắt buộc có \codeinline{name} giống hệt nhau). & \codeinline{<input type="radio" name="gender">} \\
\codeinline{value} & Giá trị dữ liệu được gửi về máy chủ khi nút Radio đó được tích chọn. & \codeinline{<input type="radio" value="male">} \\
\codeinline{checked} & Thiết lập nút Radio ở trạng thái mặc định được chọn sẵn khi tải trang. & \codeinline{<input type="radio" checked>} \\
\codeinline{required} & Bắt buộc người dùng phải chọn 1 nút Radio trong nhóm trước khi submit form. & \codeinline{<input type="radio" name="gender" required>} \\
\codeinline{disabled} & Vô hiệu hóa tùy chọn, không cho phép người dùng tương tác hay thay đổi. & \codeinline{<input type="radio" disabled>} \\
\end{tabularx}
\end{prettyTableBox}

% ------------------------------------------------------------------------------
\subsection{Ví Dụ Đầy Đủ \& Cơ Chế Gửi Dữ Liệu Lên Server}

\begin{codeWindow}[Form chọn giới tính hoàn chỉnh sử dụng Radio Button]
\begin{lstlisting}[language=HTML5]
<form action="/submit" method="POST">
  <p>(*@Chọn giới tính:@*)</p>
  <input type="radio" name="gender" value="male" checked> Nam<br>
  <input type="radio" name="gender" value="female"> (*@Nữ@*)<br>
  <input type="radio" name="gender" value="other"> (*@Khác@*)<br>
  <br>
  <button type="submit">(*@Gửi@*)</button>
</form>
\end{lstlisting}
\end{codeWindow}

\begin{browserWindow}[Kết Quả Hiển Thị Trình Duyệt: Form Chọn Giới Tính]
\sffamily\small
\textbf{Chọn giới tính:} \\[6pt]
\begin{tikzpicture}[baseline=-0.4ex]\fill[primaryBlue] (0,0) circle (6pt); \fill[white] (0,0) circle (2.2pt);\end{tikzpicture} \; \textbf{Nam} \quad \textcolor{codeComment}{\small\textit{(Checked mặc định)}} \\[4pt]
\begin{tikzpicture}[baseline=-0.4ex]\draw[primaryBlue!70, line width=1pt] (0,0) circle (6pt);\end{tikzpicture} \; Nữ \\[4pt]
\begin{tikzpicture}[baseline=-0.4ex]\draw[primaryBlue!70, line width=1pt] (0,0) circle (6pt);\end{tikzpicture} \; Khác \\[8pt]
\tcbox[on line, arc=3pt, colback=primaryBlue, colframe=primaryBlue, boxsep=2pt, left=10pt, right=10pt, top=3pt, bottom=3pt]{\color{white}\bfseries\small Gửi}
\end{browserWindow}

\begin{tipBox}[Kết Quả Dữ Liệu Gửi Về Máy Chủ (HTTP Request)]
Khi form được submit:
\begin{itemize}[leftmargin=1.5em, itemsep=2pt, parsep=0pt]
    \item Nếu chọn \textbf{Nữ}, dữ liệu gửi đi là: \codeinline{gender=female}.
    \item Nếu chọn \textbf{Nam}, dữ liệu gửi đi là: \codeinline{gender=male}.
\end{itemize}
\textit{Nguyên lý}: Chỉ duy nhất giá trị \codeinline{value} của nút Radio được tích chọn mới được đóng gói gửi đi lên máy chủ.
\end{tipBox}

% ------------------------------------------------------------------------------
\subsection{Tối Ưu Trải Nghiệm Tương Tác Với Thẻ <label>}

Nếu không có thẻ \codeinline{<label>}, người dùng bắt buộc phải căn chỉnh con trỏ chuột chính xác vào hình tròn nhỏ để chọn. Việc kết hợp với thẻ \codeinline{<label>} giúp nhấp vào chữ để chọn và tăng trải nghiệm sử dụng (UX).

\paragraph{Cách 1: Liên kết qua thuộc tính \codeinline{id} và \codeinline{for}:}
\begin{codeWindow}[Liên kết Radio Button với thẻ label qua ID]
\begin{lstlisting}[language=HTML5]
<form>
    <input type="radio" id="male" name="gender" value="male">
    <label for="male">Nam</label><br>
    
    <input type="radio" id="female" name="gender" value="female">
    <label for="female">(*@Nữ@*)</label><br>
    
    <input type="radio" id="other" name="gender" value="other">
    <label for="other">(*@Khác@*)</label><br>
</form>
\end{lstlisting}
\end{codeWindow}

\begin{browserWindow}[Mô Phỏng Trình Duyệt: Tương Tác Chọn Radio Qua Thẻ Label (Liên kết for=id)]
\sffamily\small
\begin{tikzpicture}[baseline=-0.4ex]\fill[primaryBlue] (0,0) circle (6pt); \fill[white] (0,0) circle (2.2pt);\end{tikzpicture} \; \textbf{Nam} \quad \textcolor{primaryBlue}{\small\textbf{$\leftarrow$ Nhấp chữ "Nam" (liên kết for="male") để chọn}} \\[4pt]
\begin{tikzpicture}[baseline=-0.4ex]\draw[primaryBlue!70, line width=1pt] (0,0) circle (6pt);\end{tikzpicture} \; Nữ \\[4pt]
\begin{tikzpicture}[baseline=-0.4ex]\draw[primaryBlue!70, line width=1pt] (0,0) circle (6pt);\end{tikzpicture} \; Khác
\end{browserWindow}

\paragraph{Cách 2: Bọc trực tiếp thẻ input bên trong thẻ label:}
\begin{codeWindow}[Bọc trực tiếp Radio Button trong thẻ label]
\begin{lstlisting}[language=HTML5]
<form>
  <p>(*@Phương thức thanh toán:@*)</p>
  <label>
    <input type="radio" name="payment" value="credit"> (*@Thẻ tín dụng@*)
  </label><br>
  <label>
    <input type="radio" name="payment" value="paypal"> PayPal
  </label><br>
</form>
\end{lstlisting}
\end{codeWindow}

\begin{browserWindow}[Kết Quả Hiển Thị Trình Duyệt: Tương Tác Chọn Phương Thức Thanh Toán]
\sffamily\small
\textbf{Phương thức thanh toán:} \\[6pt]
\begin{tikzpicture}[baseline=-0.4ex]\draw[primaryBlue!70, line width=1pt] (0,0) circle (6pt);\end{tikzpicture} \; Thẻ tín dụng \\[4pt]
\begin{tikzpicture}[baseline=-0.4ex]\fill[primaryBlue] (0,0) circle (6pt); \fill[white] (0,0) circle (2.2pt);\end{tikzpicture} \; \textbf{PayPal} \quad \textcolor{primaryBlue}{\small\textbf{$\leftarrow$ Nhấp vào chữ "PayPal" cũng kích hoạt chọn}}
\end{browserWindow}

% ------------------------------------------------------------------------------
\subsection{Các Lưu Ý Quan Trọng \& Lỗi Thường Gặp}

\begin{warnBox}[4 Lưu Ý Cốt Lõi Khi Sử Dụng Radio Button]
\begin{enumerate}[leftmargin=1.5em, itemsep=3pt, parsep=0pt]
    \item \textbf{Cùng \codeinline{name} để gom nhóm}: Các nút Radio phải có cùng thuộc tính \codeinline{name} thì mới tạo thành một nhóm lựa chọn duy nhất.
    \item \textbf{Lỗi quên \codeinline{name} hoặc khác \codeinline{name}}: Nếu để thuộc tính \codeinline{name} khác nhau (hoặc bỏ trống \codeinline{name}), người dùng có thể chọn nhiều nút Radio cùng một lúc $\rightarrow$ Mất đi bản chất chọn đơn của Radio Button.
    \item \textbf{Xử lý khi không chọn}: Nếu người dùng không chọn nút Radio nào và form không có thuộc tính \codeinline{required}, khi submit trình duyệt sẽ \textbf{không gửi bất kỳ giá trị nào} đại diện cho nhóm đó.
    \item \textbf{Luôn dùng \codeinline{<label>}}: Giúp tăng diện tích nhấp trên thiết bị di động, tạo sự thân thiện cho người dùng.
\end{enumerate}
\end{warnBox}

\paragraph{Minh Hoạ Lỗi Thường Gặp (Khác \codeinline{name} gây chọn nhiều):}

\begin{codeWindow}[CẢNH BÁO LỖI: Radio Button không có cùng name]
\begin{lstlisting}[language=HTML5]
<!-- (*@LỖI: Hai nút Radio có name/id riêng lẻ, không có name chung gom nhóm@*) -->
<input type="radio" id="opt1" value="A"> (*@Tùy chọn A@*)  
<input type="radio" id="opt2" value="B"> (*@Tùy chọn B@*)  
<!-- (*@HẬU QUẢ: Người dùng có thể chọn cả A và B cùng lúc (Sai logic Radio!)@*) -->
\end{lstlisting}
\end{codeWindow}

% ------------------------------------------------------------------------------
\subsection{So Sánh Chi Tiết Checkbox vs Radio Button}

Dưới đây là bảng ma trận so sánh toàn diện giúp lập trình viên đưa ra quyết định thiết kế chính xác giữa Checkbox và Radio Button:

\begin{prettyTableBox}[Ma Trận So Sánh Checkbox vs Radio Button]
\rowcolors{2}{white}{primaryBlue!4}
\begin{tabularx}{\linewidth}{K{2.8cm} Z Z}
\rowcolor{primaryBlue}
\thd{Đặc Điểm So Sánh} & \thd{Checkbox (<input type="checkbox">)} & \thd{Radio Button (<input type="radio">)} \\
Số lượng lựa chọn & Chọn \textbf{nhiều lựa chọn} cùng lúc (0, 1 hoặc nhiều). & Chọn \textbf{duy nhất 1 giá trị} trong nhóm. \\
Tác dụng của \codeinline{name} & Có thể giống nhau (gom mảng) hoặc khác nhau (độc lập). & \textbf{Bắt buộc phải giống nhau} để gom nhóm chọn duy nhất 1. \\
Cơ chế gửi dữ liệu & Gửi nhiều cặp tham số (\codeinline{hobby=music\&hobby=sport}). & Gửi \textbf{duy nhất 1 giá trị} của ô được chọn (\codeinline{gender=male}). \\
Mặc định chọn trước & Thuộc tính \codeinline{checked} có thể đặt trên \textbf{nhiều ô}. & Thuộc tính \codeinline{checked} chỉ nên đặt trên \textbf{1 ô duy nhất}. \\
Ngữ cảnh sử dụng & Chọn sở thích, danh mục kỹ năng, ngôn ngữ lập trình... & Chọn giới tính, phương thức thanh toán, hình thức vận chuyển... \\
\end{tabularx}
\end{prettyTableBox}

\paragraph{Ví Dụ Mã Nguồn So Sánh Trực Quan:}

\begin{codeWindow}[Ví dụ đối chiếu cú pháp Radio vs Checkbox]
\begin{lstlisting}[language=HTML5]
<!-- (*@Radio Button: Bắt buộc chọn duy nhất 1 giới tính@*) -->
<h3>(*@Radio (Chỉ chọn 1)@*)</h3>
<input type="radio" name="gender" value="male"> Nam  
<input type="radio" name="gender" value="female"> (*@Nữ@*)  

<!-- (*@Checkbox: Có thể chọn nhiều sở thích cùng lúc@*) -->
<h3>(*@Checkbox (Có thể chọn nhiều)@*)</h3>
<input type="checkbox" name="hobby" value="reading"> (*@Đọc sách@*)  
<input type="checkbox" name="hobby" value="travel"> (*@Du lịch@*)  
\end{lstlisting}
\end{codeWindow}

\begin{browserWindow}[Mô Phỏng Trình Duyệt: Giao Diện So Sánh Radio vs Checkbox]
\sffamily\small
\textbf{Radio Button (Chỉ chọn 1):} \\[4pt]
\begin{tikzpicture}[baseline=-0.4ex]\fill[primaryBlue] (0,0) circle (6pt); \fill[white] (0,0) circle (2.2pt);\end{tikzpicture} \; Nam \qquad
\begin{tikzpicture}[baseline=-0.4ex]\draw[primaryBlue!70, line width=1pt] (0,0) circle (6pt);\end{tikzpicture} \; Nữ \\[10pt]
\textbf{Checkbox (Có thể chọn nhiều):} \\[4pt]
\begin{tikzpicture}[baseline=-0.4ex]\fill[primaryBlue, rounded corners=2pt] (0,0) rectangle (12pt,12pt); \draw[white, line width=1.2pt] (3pt,6pt) -- (5pt,3pt) -- (9pt,9pt);\end{tikzpicture} \; \textbf{Đọc sách} \qquad
\begin{tikzpicture}[baseline=-0.4ex]\fill[primaryBlue, rounded corners=2pt] (0,0) rectangle (12pt,12pt); \draw[white, line width=1.2pt] (3pt,6pt) -- (5pt,3pt) -- (9pt,9pt);\end{tikzpicture} \; \textbf{Du lịch}
\end{browserWindow}

% ------------------------------------------------------------------------------
\subsection{Minh Họa Sơ Đồ Luồng Xử Lý Radio Button}

\begin{figure}[H]
\centering
\begin{tikzpicture}[
    auto, 
    >=stealth, 
    semithick,
    node distance=2.2cm,
    startstop/.style={rectangle, rounded corners=5pt, minimum width=3.8cm, minimum height=1cm, text centered, draw=primaryBlue, fill=primaryBlue!10, font=\sffamily\bfseries\small},
    decision/.style={diamond, aspect=2.2, minimum width=4cm, minimum height=1.1cm, text centered, draw=accentOrange, fill=accentOrange!10, font=\sffamily\bfseries\small},
    process/.style={rectangle, rounded corners=4pt, minimum width=3.8cm, minimum height=1cm, text centered, draw=secondaryTeal, fill=secondaryTeal!10, font=\sffamily\small}
]
    \node (start) [startstop] {1. Nhấp chọn Nút Radio B};
    \node (scan) [process, below of=start, node distance=1.8cm] {2. Trình duyệt quét nhóm cùng \codeinline{name}};
    \node (check) [process, below of=scan, node distance=1.8cm] {3. Bỏ tích chọn nút A cũ (\codeinline{checked = false})};
    \node (activate) [process, below of=check, node distance=1.8cm] {4. Bật tích chọn nút B mới (\codeinline{checked = true})};
    \node (end) [startstop, below of=activate, node distance=1.8cm] {5. Đóng gói 1 giá trị: \codeinline{name=valueB}};

    \draw [->] (start) -- (scan);
    \draw [->] (scan) -- (check);
    \draw [->] (check) -- (activate);
    \draw [->] (activate) -- (end);
\end{tikzpicture}
\caption{Sơ đồ quy trình loại trừ lẫn nhau (Mutually Exclusive) của nhóm Radio Button}
\label{fig:radio_button_flowchart}
\end{figure}

% ------------------------------------------------------------------------------
\subsection{Đọc Thêm (Mở Rộng Kiến Thức Chuyên Sâu)}
\label{subsec:radio_advanced_readmore}

\paragraph{a. Điều Hướng Bàn Phím Chuẩn Accessibility (a11y Keyboard Navigation):}
Radio Button sở hữu cơ chế điều hướng bàn phím rất đặc thù theo tiêu chuẩn WCAG:
\begin{itemize}[leftmargin=1.5em, itemsep=2pt, parsep=0pt]
    \item \textbf{Phím Tab}: Nhảy vào nhóm Radio (chỉ dừng lại ở nút Radio đang được tích chọn trong nhóm).
    \item \textbf{Phím mũi tên ($\leftarrow \rightarrow \uparrow \downarrow$)}: Di chuyển và tự động chuyển đổi tích chọn giữa các nút Radio trong cùng nhóm.
\end{itemize}

\paragraph{b. Tùy Biến Giao Diện Radio Button Bằng CSS Modern (\codeinline{appearance: none}):}
Tương tự Checkbox, ta có thể dùng CSS3 để tạo chấm tròn Radio đẹp mắt đồng bộ thiết kế thương hiệu:

\begin{codeWindow}[Custom giao diện Radio Button đẹp mắt với CSS3]
\begin{lstlisting}[language=HTML5]
input[type="radio"] {
  appearance: none; /* (*@Xóa giao diện mặc định hệ điều hành@*) */
  width: 20px; height: 20px;
  border: 2px solid #0056b3;
  border-radius: 50%; /* (*@Tạo hình tròn chuẩn Radio@*) */
  outline: none; cursor: pointer;
}

input[type="radio"]:checked {
  background-color: white;
  border-color: #0056b3;
  box-shadow: inset 0 0 0 4px #0056b3; /* (*@Tạo chấm tròn nhỏ bên trong@*) */
}
\end{lstlisting}
\end{codeWindow}

\begin{browserWindow}[Mô Phỏng Trình Duyệt: Custom Radio Button CSS Modern]
\sffamily\small
\begin{tikzpicture}[baseline=-0.4ex]\fill[primaryBlue] (0,0) circle (7pt); \fill[white] (0,0) circle (3pt);\end{tikzpicture} \;\textbf{Custom Radio Button Đã Chọn Với Viền \& Chấm Tròn Thương Hiệu} \\[6pt]
\begin{tikzpicture}[baseline=-0.4ex]\draw[primaryBlue, line width=1.2pt] (0,0) circle (7pt);\end{tikzpicture} \; Custom Radio Button Chưa Chọn Viền Xanh Đẹp Mắt
\end{browserWindow}

\paragraph{c. Quản Lý State Radio Button Trong React (Controlled Component):}
Trong React, nhóm Radio Button thường được liên kết với 1 biến State chuỗi duy nhất:

\begin{codeWindow}[Xử lý nhóm Radio Button trong React]
\begin{lstlisting}[language=JavaScript]
import React, { useState } from 'react';

function GenderForm() {
  const [gender, setGender] = useState('male'); // (*@State lưu 1 giá trị duy nhất@*)

  return (
    <div>
      <label>
        <input 
          type="radio" 
          name="gender" 
          value="male" 
          checked={gender === 'male'} 
          onChange={(e) => setGender(e.target.value)} 
        /> Nam
      </label>
      <label>
        <input 
          type="radio" 
          name="gender" 
          value="female" 
          checked={gender === 'female'} 
          onChange={(e) => setGender(e.target.value)} 
        /> (*@Nữ@*)
      </label>
    </div>
  );
}
\end{lstlisting}
\end{codeWindow}

% ------------------------------------------------------------------------------
\subsection{Bộ Câu Hỏi Phỏng Vấn Frontend Thực Chiến (Interview Q\&A)}

\begin{noteBox}[Câu Hỏi Phỏng Vấn 1: Điều Kiện Để Radio Button Gom Nhóm]
\textbf{Hỏi}: Điều kiện bắt buộc để nhiều nút Radio Button hoạt động loại trừ lẫn nhau (chỉ chọn được 1) là gì? Điều gì xảy ra nếu quên khai báo thuộc tính này?
\par\vspace{3pt}
\textbf{Trả lời}: Điều kiện bắt buộc là các nút Radio đó phải có cùng giá trị thuộc tính \codeinline{name}. Nếu quên khai báo hoặc đặt tên \codeinline{name} khác nhau, trình duyệt sẽ coi chúng là các nút riêng biệt và cho phép người dùng chọn tất cả cùng lúc (sai nguyên tắc UX).
\end{noteBox}

\begin{noteBox}[Câu Hỏi Phỏng Vấn 2: Dữ Liệu Radio Gửi Đi Khi Submit]
\textbf{Hỏi}: Khi Submit Form, dữ liệu của nhóm Radio Button được trình duyệt đóng gói gửi đi như thế nào?
\par\vspace{3pt}
\textbf{Trả lời}: Trình duyệt chỉ đóng gói duy nhất 1 cặp \codeinline{name=value} tương ứng với nút Radio đang được tích chọn (\codeinline{checked = true}). Các nút không được chọn trong cùng nhóm sẽ hoàn toàn bị bỏ qua.
\end{noteBox}

% ------------------------------------------------------------------------------
\subsection{Bảng Độ Tương Thích Trình Duyệt \& Best Practices Checklist}

\begin{prettyTableBox}[Bảng Độ Tương Thích Trình Duyệt Của Radio Button]
\rowcolors{2}{white}{primaryBlue!4}
\begin{tabularx}{\linewidth}{K{3.2cm}C{2.2cm}C{2.2cm}C{2.2cm}Z}
\rowcolor{primaryBlue}
\thd{Tính Năng} & \thd{Chrome / Edge} & \thd{Firefox} & \thd{Safari (iOS)} & \thd{Ghi Chú Hỗ Trợ} \\
\codeinline{<input type="radio">} & Full (v1+) & Full (v1+) & Full (v1+) & Hỗ trợ 100\% trên mọi thiết bị. \\
Phím mũi tên (Arrow Keys) & Full (v1+) & Full (v1+) & Full (v1+) & Chuẩn WCAG a11y keyboard. \\
CSS \codeinline{appearance: none} & Full (v84+) & Full (v80+) & Full (v15.4+) & Tùy biến hình tròn Radio đẹp mắt. \\
\end{tabularx}
\end{prettyTableBox}

\begin{tipBox}[Checklist 5 Quy Tắc Vàng Khi Sử Dụng Radio Button]
\begin{enumerate}[leftmargin=1.5em, itemsep=2pt, parsep=0pt]
    \item \textbf{Kiểm tra thuộc tính \codeinline{name} trùng khớp}: Luôn đảm bảo các nút trong cùng nhóm có tên \codeinline{name} giống hệt nhau.
    \item \textbf{Khai báo giá trị \codeinline{value} rõ ràng}: Đặt \codeinline{value} cụ thể cho từng nút (ví dụ \codeinline{value="credit"} thay vì để trống).
    \item \textbf{Nên chọn mặc định 1 ô (\codeinline{checked})}: Đặt \codeinline{checked} trên tùy chọn phổ biến nhất để giảm bớt 1 thao tác cho người dùng.
    \item \textbf{Luôn kết hợp với thẻ \codeinline{<label>}}: Giúp người dùng nhấp vào văn bản để chọn radio dễ dàng trên điện thoại.
    \item \textbf{Sử dụng đúng mục đích UX}: Chỉ dùng Radio Button khi số lượng lựa chọn ít (từ 2 - 5 lựa chọn). Nếu danh sách quá dài (> 5 lựa chọn), hãy sử dụng thẻ chọn \codeinline{<select>} (Dropdown Menu).
\end{enumerate}
\end{tipBox}

% ==============================================================================
% MỤC GIÁO TRÌNH: CHI TIẾT VỀ CÁC THẺ NÚT BẤM TRONG FORM (submit, button, <button>)
% ==============================================================================
\section[Các Thẻ Nút Bấm Trong Form]{Các Thẻ Nút Bấm Trong Form (<input type="submit">, <input type="button">, <button>)}
\label{sec:html_form_buttons}

\begin{conceptBox}[Tổng Quan Về Nút Bấm Trong Form]
Nút bấm (\term{Buttons}) là thành phần kích hoạt tương tác cốt lõi trong biểu mẫu HTML. Trình duyệt cung cấp hai cách chính để khai báo nút bấm: thông qua thẻ nhập liệu tự đóng \codeinline{<input>} hoặc thẻ mở/đóng linh hoạt \codeinline{<button>}.
\begin{itemize}[leftmargin=1.5em, itemsep=2pt, parsep=0pt]
    \item \textbf{Thẻ \codeinline{<input type="submit">}}: Nút gửi form truyền thống, chỉ chứa văn bản đơn thuần.
    \item \textbf{Thẻ \codeinline{<input type="button">}}: Nút bấm đơn giản dùng để kích hoạt các kịch bản JavaScript.
    \item \textbf{Thẻ \codeinline{<button>}}: Nút bấm thế hệ mới cực kỳ linh hoạt, cho phép chứa văn bản, icon, thẻ \codeinline{<span>}, hình ảnh và các phần tử HTML bên trong.
\end{itemize}
\end{conceptBox}

% ------------------------------------------------------------------------------
\subsection{Thẻ Nút Gửi Form (<input type="submit">)}
\label{subsec:input_submit}

\paragraph{Định Nghĩa \& Bản Chất:}
\codeinline{<input type="submit">} là loại nút dùng để gửi toàn bộ dữ liệu nhập trong form lên URL được chỉ định trong thuộc tính \codeinline{action} thông qua phương thức \codeinline{method} (\codeinline{GET} hoặc \codeinline{POST}). Đây là loại nút \textbf{mặc định} trong thẻ \codeinline{<form>} nếu một ô nút bấm không được khai báo rõ thuộc tính \codeinline{type}.

\begin{codeWindow}[Cú pháp cơ bản của nút input type="submit"]
\begin{lstlisting}[language=HTML5]
<form action="/submit" method="post">
  <input type="text" name="username" placeholder="(*@Tên đăng nhập@*)" required>
  <input type="submit" value="(*@Đăng nhập@*)">
</form>
\end{lstlisting}
\end{codeWindow}

\begin{browserWindow}[Kết Quả Hiển Thị Trình Duyệt: Nút Input Submit]
\sffamily\small
\tcbox[on line, arc=3pt, colback=codeBg, colframe=codeBorder, boxsep=1pt, left=6pt, right=6pt, top=2pt, bottom=2pt]{\color{darkText}\small Tên đăng nhập...} \quad
\tcbox[on line, arc=3pt, colback=primaryBlue, colframe=primaryBlue, boxsep=2pt, left=10pt, right=10pt, top=3pt, bottom=3pt]{\color{white}\bfseries\small Đăng nhập}
\end{browserWindow}

\begin{tipBox}[Giải Thích Luồng Gửi Dữ Liệu]
Khi người dùng nhấp vào nút \textbf{Đăng nhập}, trình duyệt sẽ gom tất cả dữ liệu có thuộc tính \codeinline{name} trong form (ví dụ \codeinline{username}) và đóng gói gửi tới đường dẫn \codeinline{/submit}.
\end{tipBox}

\paragraph{Các Thuộc Tính Thường Gặp Của \codeinline{<input type="submit">}:}

\begin{prettyTableBox}[Bảng Thuộc Tính Hỗ Trợ Nút Submit]
\rowcolors{2}{white}{primaryBlue!4}
\begin{tabularx}{\linewidth}{K{3.0cm} Z Z}
\rowcolor{primaryBlue}
\thd{Thuộc Tính} & \thd{Ý Nghĩa Kỹ Thuật} & \thd{Ví Dụ Cú Pháp} \\
\codeinline{value} & Nội dung văn bản hiển thị trên nút (nếu không khai báo, mặc định trình duyệt hiển thị "Submit"). & \codeinline{<input type="submit" value="Gửi">} \\
\codeinline{formaction} & Ghi đè đường dẫn \codeinline{action} của \codeinline{<form>} (chỉ áp dụng khi bấm nút này). & \codeinline{<input type="submit" formaction="/login">} \\
\codeinline{formmethod} & Ghi đè phương thức \codeinline{method} của form (\codeinline{GET} hoặc \codeinline{POST}). & \codeinline{<input type="submit" formmethod="get">} \\
\codeinline{formtarget} & Quy định nơi mở trang kết quả (\codeinline{\_self}, \codeinline{\_blank}, \codeinline{\_parent}, \codeinline{\_top}). & \codeinline{<input type="submit" formtarget="\_blank">} \\
\codeinline{formnovalidate} & Bỏ qua tất cả các kiểm tra hợp lệ (\term{Validation}) của HTML5 khi bấm nút này. & \codeinline{<input type="submit" formnovalidate>} \\
\codeinline{disabled} & Vô hiệu hóa nút gửi, không cho phép người dùng nhấp chọn. & \codeinline{<input type="submit" disabled>} \\
\end{tabularx}
\end{prettyTableBox}

\paragraph{Ví Dụ Nâng Cao (Ghi Đè Thuộc Tính Form Với Nút Submit):}

\begin{codeWindow}[Ví dụ sử dụng nhiều nút submit với các thuộc tính ghi đè formaction và formnovalidate]
\begin{lstlisting}[language=HTML5]
<form action="/register" method="post">
  <input type="email" name="email" required>
  
  <!-- (*@Nút gửi bình thường: Gửi đến /register và kiểm tra required@*) -->
  <input type="submit" value="(*@Đăng ký@*)">
  
  <!-- (*@Nút này gửi nhưng BỎ QUA kiểm tra required@*) -->
  <input type="submit" value="(*@Đăng ký nhanh@*)" formnovalidate>
  
  <!-- (*@Nút này GHI ĐÈ chuyển tiếp gửi dữ liệu đến trang /login@*) -->
  <input type="submit" value="(*@Đăng nhập@*)" formaction="/login">
</form>
\end{lstlisting}
\end{codeWindow}

\begin{browserWindow}[Mô Phỏng Trình Duyệt: Nhóm Nút Submit Ghi Đè Action]
\sffamily\small
\tcbox[on line, arc=3pt, colback=codeBg, colframe=codeBorder, boxsep=1pt, left=6pt, right=6pt, top=2pt, bottom=2pt]{\color{darkText}\small user@example.com} \\[6pt]
\tcbox[on line, arc=3pt, colback=primaryBlue, colframe=primaryBlue, boxsep=2pt, left=8pt, right=8pt, top=3pt, bottom=3pt]{\color{white}\bfseries\small Đăng ký} \quad
\tcbox[on line, arc=3pt, colback=secondaryTeal, colframe=secondaryTeal, boxsep=2pt, left=8pt, right=8pt, top=3pt, bottom=3pt]{\color{white}\bfseries\small Đăng ký nhanh} \quad
\tcbox[on line, arc=3pt, colback=codeComment, colframe=codeComment, boxsep=2pt, left=8pt, right=8pt, top=3pt, bottom=3pt]{\color{white}\bfseries\small Đăng nhập}
\end{browserWindow}

\begin{warnBox}[3 Lưu Ý Quan Trọng Khi Dùng Nút Submit]
\begin{enumerate}[leftmargin=1.5em, itemsep=3pt, parsep=0pt]
    \item \textbf{Nhiều nút submit trong cùng form}: Khi bấm nút nào, trình duyệt sẽ áp dụng các thuộc tính riêng (\codeinline{formaction}, \codeinline{formnovalidate}...) của chính nút đó để xử lý việc gửi dữ liệu.
    \item \textbf{Mặc định hiển thị chữ "Submit"}: Nếu bỏ trống thuộc tính \codeinline{value}, văn bản hiển thị trên nút sẽ mặc định là "Submit" (hoặc từ tương đương tùy theo ngôn ngữ trình duyệt người dùng).
    \item \textbf{Đặt nút ngoài form bằng thuộc tính \codeinline{form="id"}}: Nút submit có thể nằm bất kỳ đâu ngoài thẻ \codeinline{<form>} miễn là khai báo thuộc tính \codeinline{form="id\_form"}.
\end{enumerate}
\end{warnBox}

\begin{codeWindow}[Ví dụ đặt nút submit nằm bên ngoài thẻ form]
\begin{lstlisting}[language=HTML5]
<!-- (*@Form được gắn ID riêng@*) -->
<form id="myForm" action="/save">
  <input type="text" name="data">
</form>

<!-- (*@Nút Submit nằm BÊN NGOÀI form nhưng vẫn kích hoạt gửi myForm@*) -->
<input type="submit" form="myForm" value="(*@Lưu dữ liệu@*)">
\end{lstlisting}
\end{codeWindow}

\begin{browserWindow}[Mô Phỏng Trình Duyệt: Nút Submit Nằm Bên Ngoài Form Kết Nối Qua form="id"]
\sffamily\small
\textbf{Khung Form (id="myForm"):} \tcbox[on line, arc=3pt, colback=codeBg, colframe=codeBorder, boxsep=1pt, left=6pt, right=6pt, top=2pt, bottom=2pt]{\color{darkText}\small Dữ liệu nhập...} \\[6pt]
\textcolor{codeComment}{\small\textit{Nút submit nằm ngoài DOM Form nhưng khai báo form="myForm":}} \quad
\tcbox[on line, arc=3pt, colback=primaryBlue, colframe=primaryBlue, boxsep=2pt, left=10pt, right=10pt, top=3pt, bottom=3pt]{\color{white}\bfseries\small Lưu dữ liệu}
\end{browserWindow}

% ------------------------------------------------------------------------------
\subsection{Thẻ Nút Bấm JavaScript (<input type="button">)}
\label{subsec:input_button}

\paragraph{Định Nghĩa \& Đặc Điểm:}
\codeinline{<input type="button">} tạo ra một nút bấm bình thường. Nút này \textbf{không có bất kỳ hành động mặc định nào} (không gửi form, không reset form). Nó được thiết kế chuyên biệt để kết hợp với ngôn ngữ JavaScript nhằm thực thi các sự kiện tương tác (\codeinline{onclick}, \codeinline{addEventListener}).

\begin{codeWindow}[Cú pháp nút input type="button" kích hoạt sự kiện onclick]
\begin{lstlisting}[language=HTML5]
<input type="button" value="(*@Bấm tôi@*)" onclick="alert('(*@Xin chào!@*)')">
\end{lstlisting}
\end{codeWindow}

\begin{browserWindow}[Kết Quả Hiển Thị Trình Duyệt: Nút Input Button Kích Hoạt Alert]
\sffamily\small
Khi người dùng nhấp vào nút:
\tcbox[on line, arc=3pt, colback=codeBg, colframe=codeBorder, boxsep=2pt, left=10pt, right=10pt, top=3pt, bottom=3pt]{\color{darkText}\bfseries\small Bấm tôi} $\rightarrow$ Trình duyệt hiển thị hộp thoại thông báo: \textcolor{accentOrange}{\textbf{"Xin chào!"}} (Form không bị reload hay gửi đi).
\end{browserWindow}

\paragraph{Phân Biệt \codeinline{<input type="button">} vs \codeinline{<input type="submit">} Trong Form:}

\begin{codeWindow}[Form kết hợp cả input type="button" và input type="submit"]
\begin{lstlisting}[language=HTML5]
<form action="/submit" method="POST">
  <input type="text" placeholder="(*@Nhập họ tên@*)">
  
  <!-- (*@Nút này chỉ chạy JS kiểm tra, KHÔNG gửi form@*) -->
  <input type="button" value="(*@Kiểm tra@*)" onclick="alert('(*@Nút này không gửi form!@*)')">
  
  <!-- (*@Nút này GỬI FORM lên máy chủ@*) -->
  <input type="submit" value="(*@Gửi form@*)">
</form>
\end{lstlisting}
\end{codeWindow}

\begin{browserWindow}[Mô Phỏng Trình Duyệt: Phân Biệt Nút JS Button Và Nút Submit Form]
\sffamily\small
\tcbox[on line, arc=3pt, colback=codeBg, colframe=codeBorder, boxsep=1pt, left=6pt, right=6pt, top=2pt, bottom=2pt]{\color{codeComment}\small Nhập họ tên...} \quad
\tcbox[on line, arc=3pt, colback=secondaryTeal, colframe=secondaryTeal, boxsep=2pt, left=8pt, right=8pt, top=3pt, bottom=3pt]{\color{white}\bfseries\small Kiểm tra (Chạy JS)} \quad
\tcbox[on line, arc=3pt, colback=primaryBlue, colframe=primaryBlue, boxsep=2pt, left=8pt, right=8pt, top=3pt, bottom=3pt]{\color{white}\bfseries\small Gửi form (Submit)}
\end{browserWindow}

\begin{prettyTableBox}[Ma Trận So Sánh \codeinline{<input type="button">} vs \codeinline{<button type="button">}]
\rowcolors{2}{white}{primaryBlue!4}
\begin{tabularx}{\linewidth}{K{2.8cm} Z Z}
\rowcolor{primaryBlue}
\thd{Tiêu Chí} & \thd{Thẻ \codeinline{<input type="button">}} & \thd{Thẻ \codeinline{<button type="button">}} \\
Cú pháp mã HTML & \codeinline{<input type="button" value="Click">} & \codeinline{<button type="button">Click</button>} \\
Nội dung hiển thị & \textbf{Chỉ chứa văn bản thuần} lấy từ thuộc tính \codeinline{value}. & \textbf{Linh hoạt chứa HTML} (văn bản, icon SVG, hình ảnh, thẻ \codeinline{<span>}...). \\
Mức độ tùy biến UI & Hạn chế hơn khi thiết kế giao diện phức tạp. & Rất linh hoạt, dễ dàng định dạng CSS Modern. \\
Thực tế sử dụng & Ít được sử dụng hơn trong các dự án hiện đại. & \textbf{Được ưu tiên lựa chọn} trong hầu hết ứng dụng web. \\
\end{tabularx}
\end{prettyTableBox}

\begin{tipBox}[Lưu Ý Về Nút Khôi Phục Dữ Liệu \codeinline{<input type="reset">}]
Bên cạnh \codeinline{<input type="submit">} và \codeinline{<input type="button">}, HTML5 còn cung cấp nút khôi phục dữ liệu \codeinline{<input type="reset" value="Nhập lại">}. Nút này có tác dụng đưa toàn bộ các phần tử nhập liệu trong form trở về trạng thái giá trị mặc định ban đầu.
\par\vspace{3pt}
\textbf{Lời khuyên UX/UI}: Tránh lạm dụng nút Reset trong thiết kế giao diện ứng dụng web hiện đại vì người dùng có thể nhấp nhầm làm mất toàn bộ thông tin vừa điền!
\end{tipBox}

% ------------------------------------------------------------------------------
\subsection{Thẻ Nút Bấm Đa Năng Linh Hoạt (<button>)}
\label{subsec:button_tag}

\begin{conceptBox}[Định Nghĩa \& Bản Chất Vượt Trội Của Thẻ <button>]
Thẻ \codeinline{<button>} tạo ra một nút bấm tương tác đa năng trong biểu mẫu HTML. Điểm cải tiến vượt trội cốt lõi của thẻ \codeinline{<button>} so me với thẻ tự đóng \codeinline{<input type="submit">} là nó được thiết kế dưới dạng \term{Container Element} (thẻ cặp có thẻ mở \codeinline{<button>} và thẻ đóng \codeinline{</button>}).
\begin{itemize}[leftmargin=1.5em, itemsep=3pt, parsep=0pt]
    \item \textbf{Khả năng chứa nội dung HTML phong phú}: Bạn có thể lồng bất kỳ văn bản, thẻ định dạng (\codeinline{<span>}, \codeinline{<strong>}), biểu tượng SVG, hình ảnh (\codeinline{<img>}), hoặc thậm chí cả badge hiệu ứng Spinner Loading vào bên trong nút.
    \item \textbf{Linh hoạt tùy biến giao diện (CSS UI/UX)}: Dễ dàng áp dụng kỹ thuật Flexbox/Grid, hiệu ứng chuyển cảnh (\term{Transition}), Hover, Active, và Ripple effect chuẩn Modern Web.
    \item \textbf{Tương thích Accessibility (a11y)}: Hỗ trợ hoàn hảo cho trình đọc màn hình (Screen Reader) và dễ dàng kết hợp các thuộc tính ARIA (\codeinline{aria-expanded}, \codeinline{aria-controls}, \codeinline{aria-label}).
\end{itemize}
\end{conceptBox}

\paragraph{Ví Dụ Đơn Giản: Nút Bấm Chứa Icon SVG \& Thẻ HTML Con:}

\begin{codeWindow}[Cú pháp thẻ button chứa icon SVG và văn bản HTML]
\begin{lstlisting}[language=HTML5]
<form action="/submit" method="post">
  <button type="submit" class="btn-send">
    <!-- (*@Hình ảnh hoặc biểu tượng SVG lồng bên trong nút@*) -->
    <svg width="16" height="16" viewBox="0 0 24 24" fill="currentColor">
      <path d="M2.01 21L23 12 2.01 3 2 10l15 2-15 2z"/>
    </svg>
    <span>(*@Gửi Thông Tin Đăng Ký@*)</span>
  </button>
</form>
\end{lstlisting}
\end{codeWindow}

\begin{browserWindow}[Mô Phỏng Trình Duyệt: Giao Diện Thẻ Button Lồng Icon SVG \& Text]
\sffamily\small
\tcbox[on line, arc=4pt, colback=primaryBlue, colframe=primaryBlue, boxsep=2pt, left=12pt, right=12pt, top=4pt, bottom=4pt]{\color{white}\bfseries\small $\rightarrow$ \quad Gửi Thông Tin Đăng Ký}
\end{browserWindow}

% ------------------------------------------------------------------------------
\paragraph{3 Giá Trị Thuộc Tính \codeinline{type} Trong Thẻ \codeinline{<button>}:}

Khác với thẻ \codeinline{<input>}, thuộc tính \codeinline{type} trong thẻ \codeinline{<button>} quy định trực tiếp hành động mặc định của nút khi người dùng nhấp chuột:

\begin{prettyTableBox}[Bảng Phân Loại Giá Trị Thuộc Tính type Trong Thẻ <button>]
\rowcolors{2}{white}{primaryBlue!4}
\begin{tabularx}{\linewidth}{K{2.8cm} Z Z Z}
\rowcolor{primaryBlue}
\thd{Giá Trị \codeinline{type}} & \thd{Hành Động Mặc Định} & \thd{Ngữ Cảnh Áp Dụng Thực Tế} & \thd{Ví Dụ Cú Pháp} \\
\codeinline{type="submit"}\newline(\textbf{Mặc định}) & Tự động gom dữ liệu và gửi form lên Server (giống \codeinline{<input type="submit">}). & Dùng cho nút gửi form đăng ký, đăng nhập, thanh toán. & \codeinline{<button type="submit">Gửi</button>} \\
\codeinline{type="reset"} & Tự động xóa và khôi phục tất cả ô nhập về trạng thái ban đầu. & Dùng cho nút xóa dữ liệu vừa nhập để điền lại từ đầu. & \codeinline{<button type="reset">Xóa làm lại</button>} \\
\codeinline{type="button"} & \textbf{Vô hướng (Dead-end)}, không có bất kỳ hành động mặc định nào. & Dùng để kích hoạt kịch bản JavaScript (Modal, Dropdown, AJAX). & \codeinline{<button type="button" onclick="...">Bấm</button>} \\
\end{tabularx}
\end{prettyTableBox}

\begin{warnBox}[CẢNH BÁO NGUY HIỂM 1: Bẫy Mặc Định Khi Bỏ Quên Thuộc Tính `type`]
Nếu bạn khai báo thẻ \codeinline{<button>} nằm bên trong thẻ \codeinline{<form>} mà \textbf{bỏ quên thuộc tính \codeinline{type}}, trình duyệt sẽ tự động coi đó là \codeinline{type="submit"}. 
\par\vspace{3pt}
$\rightarrow$ \textbf{Hậu quả nghiêm trọng}: Khi người dùng nhấp vào một nút bấm vốn chỉ dành để mở Popup/Modal hoặc chạy đoạn mã JavaScript, trình duyệt sẽ tự động kích hoạt gửi form và reload lại toàn bộ trang ngoài ý muốn!
\par\vspace{3pt}
$\rightarrow$ \textbf{Quy Tắc Vàng Senior}: Luôn luôn khai báo rõ ràng \codeinline{type="button"} cho tất cả các nút bấm xử lý JavaScript.
\end{warnBox}

\begin{warnBox}[CẢNH BÁO NGUY HIỂM 2: Lỗi Thẻ HTML Không Hợp Lệ (Invalid HTML Nesting)]
Theo chuẩn đặc tả W3C HTML5 Specification, thẻ \codeinline{<button>} thuộc nhóm \term{Interactive Content}. Do đó:
\begin{itemize}[leftmargin=1.5em, itemsep=2pt, parsep=0pt]
    \item \textbf{Tuyệt đối KHÔNG lồng thẻ \codeinline{<a>} (chứa \codeinline{href}) vào trong \codeinline{<button>}}.
    \item \textbf{Tuyệt đối KHÔNG lồng một thẻ \codeinline{<button>} khác hoặc thẻ \codeinline{<input>} vào trong \codeinline{<button>}}.
\end{itemize}
Việc lồng sai này sẽ khiến trình duyệt phá vỡ cây DOM (\term{DOM Parsing Error}) và làm hỏng trải nghiệm người dùng.
\end{warnBox}

% ------------------------------------------------------------------------------
\paragraph{Cơ Chế Đặc Biệt: Gửi Dữ Liệu Với Thuộc Tính \codeinline{name} \& \codeinline{value} Trên Nút Submit:}

Một tính năng cực kỳ mạnh mẽ của thẻ \codeinline{<button type="submit">} là nó cho phép gán thuộc tính \codeinline{name} và \codeinline{value}. Khi người dùng nhấp vào một nút Submit cụ thể trong Form, \textbf{chỉ duy nhất \codeinline{name} và \codeinline{value} của nút được nhấp mới được đóng gói gửi đi cùng dữ liệu Form}!

\begin{tipBox}[Mẹo Senior: Xử Lý Nhiều Hành Động Submit Trong 1 Form Mà Không Cần JavaScript]
Bạn có thể tạo ra các nút Submit khác nhau trong cùng 1 form (ví dụ nút "Lưu nháp" và nút "Xuất bản bài viết"). Máy chủ Backend dựa vào giá trị \codeinline{value} nhận được để phân loại hành động xử lý:
\end{tipBox}

\begin{codeWindow}[Ví dụ xử lý nhiều nút Submit với name và value khác nhau]
\begin{lstlisting}[language=HTML5]
<form action="/article/save" method="post">
  <input type="text" name="title" placeholder="(*@Tiêu đề bài viết@*)" required>
  
  <!-- (*@Nút 1: Gửi name="action" và value="draft"@*) -->
  <button type="submit" name="action" value="draft">(*@Lưu Nháp@*)</button>
  
  <!-- (*@Nút 2: Gửi name="action" và value="publish"@*) -->
  <button type="submit" name="action" value="publish">(*@Xuất Bản Ngay@*)</button>
</form>
\end{lstlisting}
\end{codeWindow}

\begin{browserWindow}[Mô Phỏng Trình Duyệt: Form 2 Nút Submit Phân Loại Hành Động]
\sffamily\small
\tcbox[on line, arc=3pt, colback=codeBg, colframe=codeBorder, boxsep=1pt, left=6pt, right=6pt, top=2pt, bottom=2pt]{\color{darkText}\small Tiêu đề bài viết mới...} \\[6pt]
\tcbox[on line, arc=3pt, colback=codeComment, colframe=codeComment, boxsep=2pt, left=8pt, right=8pt, top=3pt, bottom=3pt]{\color{white}\bfseries\small Lưu Nháp (action=draft)} \quad
\tcbox[on line, arc=3pt, colback=primaryBlue, colframe=primaryBlue, boxsep=2pt, left=8pt, right=8pt, top=3pt, bottom=3pt]{\color{white}\bfseries\small Xuất Bản Ngay (action=publish)}
\end{browserWindow}

% ------------------------------------------------------------------------------
\paragraph{Bảng Thuộc Tính Mở Rộng Đầy Đủ Của Thẻ \codeinline{<button>}:}

\begin{prettyTableBox}[Bảng Thuộc Tính Mở Rộng Hỗ Trợ Bởi Thẻ <button>]
\rowcolors{2}{white}{primaryBlue!4}
\begin{tabularx}{\linewidth}{K{3.0cm} Z Z}
\rowcolor{primaryBlue}
\thd{Thuộc Tính} & \thd{Ý Nghĩa Chức Năng Kỹ Thuật} & \thd{Ví Dụ Cú Pháp} \\
\codeinline{type} & Quy định loại nút (\codeinline{submit}, \codeinline{reset}, \codeinline{button}). & \codeinline{<button type="button">Click</button>} \\
\codeinline{name} & Tên trường đại diện cho nút bấm khi submit form. & \codeinline{<button type="submit" name="btn">} \\
\codeinline{value} & Giá trị gửi về máy chủ khi nút đó được nhấp. & \codeinline{<button type="submit" value="delete">} \\
\codeinline{form} & Liên kết nút bấm với ID của một thẻ \codeinline{<form>} nằm bất kỳ đâu ngoài DOM. & \codeinline{<button type="submit" form="mainForm">} \\
\codeinline{formaction} & Ghi đè đường dẫn \codeinline{action} của Form khi gửi bằng nút này. & \codeinline{<button type="submit" formaction="/save">} \\
\codeinline{formmethod} & Ghi đè phương thức gửi dữ liệu (\codeinline{GET} hoặc \codeinline{POST}). & \codeinline{<button type="submit" formmethod="get">} \\
\codeinline{formtarget} & Quy định nơi mở cửa sổ kết quả (\codeinline{\_blank}, \codeinline{\_self}...). & \codeinline{<button type="submit" formtarget="\_blank">} \\
\codeinline{formnovalidate} & Gửi form mà không chạy kiểm tra Validation HTML5. & \codeinline{<button type="submit" formnovalidate>} \\
\codeinline{disabled} & Vô hiệu hóa nút, đổi giao diện xám và không thể nhấp chọn. & \codeinline{<button type="submit" disabled>} \\
\codeinline{autofocus} & Tự động trỏ con trỏ/focus vào nút ngay khi trang tải xong. & \codeinline{<button type="button" autofocus>} \\
\end{tabularx}
\end{prettyTableBox}

% ------------------------------------------------------------------------------
\paragraph{Ví Dụ Tổng Hợp Đầy Đủ Tất Cả Các Loại Nút Bấm Trong Form:}

\begin{codeWindow}[Ví dụ tổng hợp đầy đủ các loại nút submit, reset, JS button và nút ngoài form]
\begin{lstlisting}[language=HTML5]
<form id="regForm" action="/register" method="post">
  <input type="text" name="user" required placeholder="(*@Tên đăng nhập@*)">
  
  <!-- (*@1. Nút Submit chính: Gửi dữ liệu form@*) -->
  <button type="submit" class="btn-primary">(*@Đăng ký tài khoản@*)</button>
  
  <!-- (*@2. Nút Reset: Khôi phục ô nhập về giá trị ban đầu@*) -->
  <button type="reset" class="btn-secondary">(*@Xóa dữ liệu@*)</button>
  
  <!-- (*@3. Nút JS Button: Kích hoạt hàm kiểm tra JavaScript@*) -->
  <button type="button" onclick="alert('(*@Đã bấm nút JS!@*)')">(*@Nhấn thử JS@*)</button>
</form>

<!-- (*@4. Nút Submit nằm BÊN NGOÀI form nhưng liên kết qua form="regForm"@*) -->
<button type="submit" form="regForm" formaction="/login" class="btn-outer">
  (*@Đăng nhập (Nút ngoài Form)@*)
</button>
\end{lstlisting}
\end{codeWindow}

\begin{browserWindow}[Mô Phỏng Trình Duyệt: Form Đầy Đủ Tất Cả Các Loại Nút Bấm]
\sffamily\small
\tcbox[on line, arc=3pt, colback=codeBg, colframe=codeBorder, boxsep=1pt, left=6pt, right=6pt, top=2pt, bottom=2pt]{\color{darkText}\small Tên đăng nhập...} \\[6pt]
\tcbox[on line, arc=3pt, colback=primaryBlue, colframe=primaryBlue, boxsep=2pt, left=8pt, right=8pt, top=3pt, bottom=3pt]{\color{white}\bfseries\small Đăng ký tài khoản} \quad
\tcbox[on line, arc=3pt, colback=accentOrange, colframe=accentOrange, boxsep=2pt, left=8pt, right=8pt, top=3pt, bottom=3pt]{\color{white}\bfseries\small Xóa dữ liệu} \quad
\tcbox[on line, arc=3pt, colback=secondaryTeal, colframe=secondaryTeal, boxsep=2pt, left=8pt, right=8pt, top=3pt, bottom=3pt]{\color{white}\bfseries\small Nhấn thử JS} \\[6pt]
\textcolor{codeComment}{\small\textit{Nút bấm nằm bên ngoài form:}} \quad
\tcbox[on line, arc=3pt, colback=codeComment, colframe=codeComment, boxsep=2pt, left=8pt, right=8pt, top=3pt, bottom=3pt]{\color{white}\bfseries\small Đăng nhập (formaction="/login")}
\end{browserWindow}

% ------------------------------------------------------------------------------
\subsection{Thẻ Nút Khôi Phục Dữ Liệu Reset Form (<input type="reset"> \& <button type="reset">)}
\label{subsec:form_reset_button}

\begin{conceptBox}[Định Nghĩa \& Bản Chất Của Nút Reset Form]
Nút \codeinline{type="reset"} (khai báo qua \codeinline{<input type="reset">} hoặc \codeinline{<button type="reset">}) là loại nút bấm chức năng đặc biệt dùng để \textbf{xóa toàn bộ thông tin người dùng vừa nhập} trong biểu mẫu và khôi phục tất cả các ô nhập liệu về \textbf{trạng thái / giá trị khởi tạo ban đầu} (\term{initial default state}).
\begin{itemize}[leftmargin=1.5em, itemsep=3pt, parsep=0pt]
    \item \textbf{Phạm vi hoạt động}: Chỉ có tác dụng khi được đặt bên trong thẻ \codeinline{<form>} (hoặc liên kết qua thuộc tính \codeinline{form="id"}).
    \item \textbf{Cơ chế đối với Server}: \textbf{KHÔNG gửi dữ liệu đi}, không phát sinh bất kỳ HTTP Request nào tới máy chủ Backend và không làm reload lại trang.
\end{itemize}
\end{conceptBox}

\paragraph{Cú Pháp Cơ Bản \& Ví Dụ Minh Họa Khôi Phục Giá Trị Mặc Định:}

Dưới đây là ví dụ minh họa chi tiết sự khác biệt giữa các ô nhập có khai báo thuộc tính \codeinline{value} mặc định và không có thuộc tính \codeinline{value}:

\begin{codeWindow}[Cú pháp sử dụng nút reset với ô nhập có giá trị mặc định]
\begin{lstlisting}[language=HTML5]
<form action="/update-profile" method="post">
  <!-- (*@Ô 1: Có giá trị mặc định value="Nguyễn Văn A"@*) -->
  <label>(*@Họ và tên:@*)</label>
  <input type="text" name="fullname" value="(*@Nguyễn Văn A@*)">

  <!-- (*@Ô 2: Không có giá trị mặc định (rỗng)@*) -->
  <label>(*@Email mới:@*)</label>
  <input type="email" name="email" placeholder="(*@nhập email@*)">

  <!-- (*@Nút Reset dạng input và button@*) -->
  <input type="reset" value="(*@Xóa làm lại@*)">
  <button type="reset">Reset Form</button>
  <button type="submit">(*@Gửi Form@*)</button>
</form>
\end{lstlisting}
\end{codeWindow}

\begin{browserWindow}[Mô Phỏng Trình Duyệt: Giao Diện Form \& Tương Tác Khi Bấm Nút Reset Form]
\sffamily\small
\textbf{Giao diện Form hiển thị trên trình duyệt:} \\[4pt]
Họ và tên: \tcbox[on line, arc=3pt, colback=codeBg, colframe=codeBorder, boxsep=1pt, left=6pt, right=6pt, top=2pt, bottom=2pt]{\color{darkText}\small Nguyễn Văn A} \quad
Email mới: \tcbox[on line, arc=3pt, colback=codeBg, colframe=codeBorder, boxsep=1pt, left=6pt, right=6pt, top=2pt, bottom=2pt]{\color{codeComment}\small nhập email...} \\[8pt]
\tcbox[on line, arc=3pt, colback=codeComment, colframe=codeComment, boxsep=2pt, left=8pt, right=8pt, top=3pt, bottom=3pt]{\color{white}\bfseries\small Xóa làm lại} \quad
\tcbox[on line, arc=3pt, colback=accentOrange, colframe=accentOrange, boxsep=2pt, left=8pt, right=8pt, top=3pt, bottom=3pt]{\color{white}\bfseries\small Reset Form} \quad
\tcbox[on line, arc=3pt, colback=primaryBlue, colframe=primaryBlue, boxsep=2pt, left=8pt, right=8pt, top=3pt, bottom=3pt]{\color{white}\bfseries\small Gửi Form} \\[10pt]
{\color{codeBorder}\hrule height 0.5pt}\vspace{6pt}
\textcolor{accentOrange}{\textbf{Cơ chế khi nhấp chọn nút Reset Form (\codeinline{<button type="reset">}):}} \\
$\rightarrow$ Ô \textbf{Họ và tên} khôi phục về giá trị mặc định: \textbf{"Nguyễn Văn A"} (lấy từ thuộc tính \codeinline{value}). \\
$\rightarrow$ Ô \textbf{Email mới} khôi phục về trạng thái rỗng: \textbf{""} (do không có thuộc tính \codeinline{value}).
\end{browserWindow}

% ------------------------------------------------------------------------------
\paragraph{So Sánh Chi Tiết Hành Động Của 3 Loại Nút Bấm Trong Form:}

\begin{prettyTableBox}[Ma Trận So Sánh Hành Động Mặc Định Của Các Loại Nút Bấm]
\rowcolors{2}{white}{primaryBlue!4}
\begin{tabularx}{\linewidth}{K{2.8cm} Z Z Z}
\rowcolor{primaryBlue}
\thd{Loại Nút (\codeinline{type})} & \thd{Hành Động Mặc Định} & \thd{Tác Động Dữ Liệu} & \thd{Phát Request Server?} \\
\codeinline{type="submit"} & Gửi dữ liệu form tới đường dẫn \codeinline{action}. & Đóng gói toàn bộ \codeinline{name=value}. & \textbf{CÓ} (Tải lại trang hoặc AJAX). \\
\codeinline{type="reset"} & Xóa/khôi phục ô nhập về giá trị ban đầu. & Trả các input về \codeinline{initial value}. & \textbf{KHÔNG} (Chỉ xử lý ở Client). \\
\codeinline{type="button"} & \textbf{Vô hướng}, không làm gì mặc định. & Không tự động thay đổi dữ liệu. & \textbf{KHÔNG} (Chờ gắn sự kiện JS). \\
\end{tabularx}
\end{prettyTableBox}

% ------------------------------------------------------------------------------
\paragraph{3 Lưu Ý Kỹ Thuật Cốt Lõi Khi Sử Dụng Nút Reset:}

\begin{warnBox}[LƯU Ý QUAN TRỌNG: 3 Quy Tắc Cần Nhớ Về Nút Reset]
\begin{enumerate}[leftmargin=1.5em, itemsep=3pt, parsep=0pt]
    \item \textbf{Reset đưa về giá trị ban đầu, KHÔNG PHẢI xóa rỗng hoàn toàn}:
    \begin{itemize}[leftmargin=1.2em, itemsep=1pt]
        \item Nếu ô nhập có sẵn \codeinline{<input value="abc">} $\rightarrow$ Khi bấm Reset, giá trị quay về \textbf{"abc"}.
        \item Nếu ô nhập không có \codeinline{value} $\rightarrow$ Khi bấm Reset, giá trị quay về rỗng \textbf{""}.
    \end{itemize}
    \item \textbf{Phạm vi tác động gói gọn trong Form}: Nút Reset chỉ xóa các ô nhập liệu thuộc cùng thẻ \codeinline{<form>} chứa nó (hoặc liên kết qua thuộc tính \codeinline{form="id"}). Các ô nằm ngoài form sẽ không bị ảnh hưởng.
    \item \textbf{Xóa sạch 100\% về rỗng phải dùng JavaScript}: Nếu muốn xóa sạch tất cả các ô về rỗng (bất chấp ô đó có \codeinline{value} mặc định hay không), bạn bắt buộc phải dùng đoạn mã JavaScript để can thiệp vào cây DOM.
\end{enumerate}
\end{warnBox}

\begin{warnBox}[CẢNH BÁO NGUY HIỂM UX: Tại Sao Nên Hạn Chế Sử Dụng Nút Reset Trên Form Hiện Đại?]
Trong thiết kế trải nghiệm người dùng (UX/UI Design) hiện đại, các chuyên gia hàng đầu \textbf{KHUYÊN NÊN TRÁNH DÙNG NÚT RESET} trên các biểu mẫu giao diện:
\par\vspace{3pt}
$\rightarrow$ \textbf{Rủi ro nghiêm trọng}: Người dùng thường vô tình nhấp nhầm vào nút "Reset" (xóa) đặt ngay cạnh nút "Submit" (gửi), dẫn đến việc toàn bộ thông tin họ vừa tốn nhiều thời gian điền bị xóa sạch!
\end{warnBox}

% ------------------------------------------------------------------------------
\subsection{So Sánh Chi Tiết \codeinline{<input type="submit">} vs \codeinline{<button type="submit">}}

Dưới đây là ma trận so sánh toàn diện giúp bạn đưa ra lựa chọn chính xác khi thiết kế biểu mẫu HTML5:

\begin{prettyTableBox}[Ma Trận So Sánh \codeinline{<input type="submit">} vs \codeinline{<button type="submit">}]
\rowcolors{2}{white}{primaryBlue!4}
\begin{tabularx}{\linewidth}{K{3.0cm} Z Z}
\rowcolor{primaryBlue}
\thd{Tiêu Chí So Sánh} & \thd{Thẻ \codeinline{<input type="submit">}} & \thd{Thẻ \codeinline{<button type="submit">}} \\
Cấu trúc phần tử HTML & Thẻ tự đóng (\term{Void Element}), không có phần thân chứa thẻ con. & Thẻ cặp (\term{Container Element}), chứa được văn bản, icon, thẻ HTML con. \\
Nội dung hiển thị & \textbf{Chỉ chứa văn bản thuần} lấy từ thuộc tính \codeinline{value}. & \textbf{Cực kỳ linh hoạt}: Văn bản, icon SVG, hình ảnh, spinner badge,... \\
Tùy biến giao diện (UI) & Hạn chế, khó tạo các hiệu ứng nút bấm hiện đại. & \textbf{Rất dễ dàng tùy biến CSS Modern} (Flexbox, Pseudo-elements, Ripple). \\
Hỗ trợ thuộc tính & Hỗ trợ các thuộc tính ghi đè \codeinline{formaction}, \codeinline{formnovalidate}... & Hỗ trợ đầy đủ các thuộc tính tương tự. \\
Khuyên dùng thực tế & Dùng cho các biểu mẫu đơn giản, không yêu cầu thiết kế icon. & \textbf{Khuyên dùng cho 100\% dự án ứng dụng Web hiện đại}. \\
\end{tabularx}
\end{prettyTableBox}

\begin{conceptBox}[Tóm Tắt Quy Tắc Lựa Chọn Thẻ Nút Bấm Chuẩn Senior Engineer]
\begin{itemize}[leftmargin=1.5em, itemsep=3pt, parsep=0pt]
    \item \textbf{Ưu tiên 100\% sử dụng thẻ \codeinline{<button>}} khi thiết kế giao diện ứng dụng Web hiện đại để dễ dàng làm chủ CSS, thêm icon SVG và cải thiện trải nghiệm người dùng (UX).
    \item \textbf{Luôn luôn khai báo thuộc tính \codeinline{type}} rõ ràng (\codeinline{type="submit"}, \codeinline{type="button"}, \codeinline{type="reset"}) để tránh bẫy tự động submit form ngoài ý muốn.
\end{itemize}
\end{conceptBox}

% ------------------------------------------------------------------------------
\subsection{Minh Họa Sơ Đồ Luồng Xử Lý Submit Form Khi Nhấp Nút}

\begin{figure}[H]
\centering
\begin{tikzpicture}[
    auto, 
    >=stealth, 
    semithick,
    node distance=2.2cm,
    startstop/.style={rectangle, rounded corners=5pt, minimum width=3.8cm, minimum height=1cm, text centered, draw=primaryBlue, fill=primaryBlue!10, font=\sffamily\bfseries\small},
    decision/.style={diamond, aspect=2.2, minimum width=4cm, minimum height=1.1cm, text centered, draw=accentOrange, fill=accentOrange!10, font=\sffamily\bfseries\small},
    process/.style={rectangle, rounded corners=4pt, minimum width=3.8cm, minimum height=1cm, text centered, draw=secondaryTeal, fill=secondaryTeal!10, font=\sffamily\small}
]
    \node (start) [startstop] {1. Bấm Nút Submit Form};
    \node (check) [decision, below of=start, node distance=1.8cm] {Có \codeinline{formnovalidate}?};
    \node (val) [process, right=2.5cm of check] {2. Kiểm Tra Validation HTML5};
    \node (skip) [process, left=2.5cm of check] {Bỏ Qua Validation};
    \node (override) [decision, below of=check, node distance=2.4cm] {Có \codeinline{formaction}?};
    \node (custom) [process, left=2.5cm of override] {Gửi Đến URL Nút Submit};
    \node (default) [process, right=2.5cm of override] {Gửi Đến URL Action Form};
    \node (end) [startstop, below of=override, node distance=2.4cm] {3. Phát HTTP Request};

    \draw [->] (start) -- (check);
    \draw [->] (check) -- node [above=2pt, font=\sffamily\bfseries\scriptsize, fill=white, inner sep=1.5pt] {\color{accentOrange}KHÔNG} (val);
    \draw [->] (check) -- node [above=2pt, font=\sffamily\bfseries\scriptsize, fill=white, inner sep=1.5pt] {\color{primaryBlue}CÓ} (skip);
    \draw [->] (val) |- ([yshift=12pt]override.north) -- (override.north);
    \draw [->] (skip) |- ([yshift=12pt]override.north) -- (override.north);
    \draw [->] (override) -- node [above=2pt, font=\sffamily\bfseries\scriptsize, fill=white, inner sep=1.5pt] {\color{primaryBlue}CÓ} (custom);
    \draw [->] (override) -- node [above=2pt, font=\sffamily\bfseries\scriptsize, fill=white, inner sep=1.5pt] {\color{accentOrange}KHÔNG} (default);
    \draw [->] (custom) |- (end);
    \draw [->] (default) |- (end);
\end{tikzpicture}
\caption{Sơ đồ luồng quyết định gửi dữ liệu và xử lý thuộc tính ghi đè (Override Attributes) của nút Submit}
\label{fig:button_submit_flowchart}
\end{figure}

% ------------------------------------------------------------------------------
\subsection{Đọc Thêm (Mở Rộng Kiến Thức Chuyên Sâu)}
\label{subsec:buttons_advanced_readmore}

\paragraph{a. Trải Nghiệm Phím Enter Trong Form (Enter Key Form Submit):}
Mặc định trên các trình duyệt, khi người dùng đang nhập dữ liệu trong một ô \codeinline{<input type="text">} và nhấn phím \textbf{Enter}, trình duyệt sẽ tự động tìm kiếm nút \codeinline{type="submit"} đầu tiên trong form và kích hoạt sự kiện Submit.

\paragraph{b. Styling Nút Bấm Đẹp Với CSS Modern (Flexbox, Icon SVG \& CSS Pseudo-classes):}
Kỹ thuật thiết kế nút bấm chuẩn UI/UX hiện đại với hiệu ứng Hover, Active và Disabled:

\begin{codeWindow}[Mẫu CSS nút bấm chuẩn UI/UX hiện đại]
\begin{lstlisting}[language=HTML5]
.btn-primary {
  display: inline-flex;
  align-items: center;
  gap: 8px;
  padding: 10px 20px;
  background-color: #0f4c81;
  color: white;
  border: none;
  border-radius: 6px;
  font-weight: 600;
  cursor: pointer;
  transition: all 0.2s ease-in-out;
}

.btn-primary:hover {
  background-color: #0a355c;
  transform: translateY(-1px);
}

.btn-primary:active {
  transform: translateY(0);
}

.btn-primary:disabled {
  opacity: 0.6;
  cursor: not-allowed;
}
\end{lstlisting}
\end{codeWindow}

\begin{browserWindow}[Mô Phỏng Trình Duyệt: Custom Modern CSS Button]
\sffamily\small
\tcbox[on line, arc=4pt, colback=primaryBlue, colframe=primaryBlue, boxsep=2pt, left=12pt, right=12pt, top=4pt, bottom=4pt]{\color{white}\bfseries\small [Register] Đăng Ký Tài Khoản (Hover State)} \qquad
\tcbox[on line, arc=4pt, colback=codeComment!50, colframe=codeComment!50, boxsep=2pt, left=12pt, right=12pt, top=4pt, bottom=4pt]{\color{white}\bfseries\small Đang Xử Lý... (Disabled State)}
\end{browserWindow}

\paragraph{c. Xử Lý Submit Form Trong React (e.preventDefault()):}
Trong các ứng dụng Single Page Application (SPA) như React, ta sử dụng \codeinline{e.preventDefault()} để ngăn trình duyệt reload lại trang khi nhấp nút Submit:

\begin{codeWindow}[Xử lý sự kiện Submit Form trong React]
\begin{lstlisting}[language=JavaScript]
import React from 'react';

function LoginForm() {
  const handleSubmit = (event) => {
    event.preventDefault(); // (*@Ngăn trình duyệt reload lại trang@*)
    console.log('(*@Đã gửi dữ liệu qua API Ajax/Fetch@*)');
  };

  return (
    <form onSubmit={handleSubmit}>
      <input type="text" name="username" />
      {/* (*@Bắt buộc khai báo type="submit" cho nút gửi form@*) */}
      <button type="submit">(*@Đăng nhập@*)</button>
    </form>
  );
}
\end{lstlisting}
\end{codeWindow}

% ------------------------------------------------------------------------------
\subsection{Bộ Câu Hỏi Phỏng Vấn Frontend Thực Chiến (Interview Q\&A)}

\begin{noteBox}[Câu Hỏi Phỏng Vấn 1: Sự Khác Biệt Giữa \codeinline{<input type="submit">} và \codeinline{<button type="submit">}]
\textbf{Hỏi}: Tại sao trong các dự án ứng dụng web hiện đại, các lập trình viên Frontend thường ưu tiên sử dụng thẻ \codeinline{<button type="submit">} hơn thẻ \codeinline{<input type="submit">}?
\par\vspace{3pt}
\textbf{Trả lời}: Bởi vì thẻ \codeinline{<button>} cho phép lồng các phần tử HTML phức tạp bên trong (như icon SVG, spinner loading, hình ảnh, thẻ \codeinline{<span>} định dạng chữ), giúp việc tùy biến giao diện (UI/UX) và tạo hiệu ứng linh hoạt hơn rất nhiều so với thẻ \codeinline{<input>} vốn chỉ nhận chuỗi văn bản thuần qua thuộc tính \codeinline{value}.
\end{noteBox}

\begin{noteBox}[Câu Hỏi Phỏng Vấn 2: Sự Cố Bỏ Quên \codeinline{type} Trong Thẻ \codeinline{<button>}]
\textbf{Hỏi}: Điều gì xảy ra nếu bạn đặt một thẻ \codeinline{<button>} xử lý sự kiện nhấp chuột bằng JavaScript trong form nhưng quên khai báo thuộc tính \codeinline{type="button"}?
\par\vspace{3pt}
\textbf{Trả lời}: Mặc định theo đặc tả HTML5, thẻ \codeinline{<button>} không có thuộc tính \codeinline{type} sẽ tự động mặc định là \codeinline{type="submit"}. Do đó khi người dùng nhấp vào, ngoài việc chạy đoạn mã JavaScript, trình duyệt sẽ tự động kích hoạt gửi form và reload lại trang, gây ra lỗi ngoài ý muốn.
\end{noteBox}

\begin{noteBox}[Câu Hỏi Phỏng Vấn 3: Cơ Chế Khôi Phục Dữ Liệu Của Nút Reset Form]
\textbf{Hỏi}: Tại sao khi nhấp vào nút \codeinline{<button type="reset">}, một số ô nhập liệu lại không bị xóa rỗng mà lại hiển thị lại nội dung cũ? Muốn xóa sạch 100\% về rỗng thì làm thế nào?
\par\vspace{3pt}
\textbf{Trả lời}: Nút Reset trong HTML5 không có chức năng "xóa sạch về rỗng" mà là "khôi phục về giá trị mặc định khởi tạo ban đầu (\term{initial HTML value attribute})". Nếu ô nhập có sẵn \codeinline{value="abc"}, nút Reset sẽ đưa ô đó về lại chữ \codeinline{"abc"}. Muốn xóa sạch 100\% tất cả các ô về rỗng (bất kể có \codeinline{value} hay không), lập trình viên phải dùng đoạn mã JavaScript can thiệp đặt \codeinline{input.value = ''} hoặc lặp qua các phần tử điều khiển của Form.
\end{noteBox}

% ------------------------------------------------------------------------------
\subsection{Bảng Độ Tương Thích Trình Duyệt \& Best Practices Checklist}

\begin{prettyTableBox}[Bảng Độ Tương Thích Trình Duyệt Của Thẻ Nút Bấm]
\rowcolors{2}{white}{primaryBlue!4}
\begin{tabularx}{\linewidth}{K{3.2cm}C{2.2cm}C{2.2cm}C{2.2cm}Z}
\rowcolor{primaryBlue}
\thd{Tính Năng} & \thd{Chrome / Edge} & \thd{Firefox} & \thd{Safari (iOS)} & \thd{Ghi Chú Hỗ Trợ} \\
\codeinline{<input type="submit">} & Full (v1+) & Full (v1+) & Full (v1+) & Hỗ trợ 100\% trên mọi trình duyệt. \\
Thẻ \codeinline{<button>} & Full (v1+) & Full (v1+) & Full (v1+) & Hỗ trợ tốt các phần tử con HTML. \\
\codeinline{formaction / form...} & Full (v9+) & Full (v4+) & Full (v5.1+) & Ghi đè thuộc tính form ở từng nút. \\
Thuộc tính \codeinline{form="id"} & Full (v10+) & Full (v4+) & Full (v5.1+) & Đặt nút ngoài thẻ form. \\
\end{tabularx}
\end{prettyTableBox}

\begin{tipBox}[Checklist 5 Quy Tắc Vàng Khi Làm Việc Với Nút Bấm Form]
\begin{enumerate}[leftmargin=1.5em, itemsep=2pt, parsep=0pt]
    \item \textbf{Luôn khai báo thuộc tính \codeinline{type} rõ ràng}: Luôn viết rõ \codeinline{type="submit"}, \codeinline{type="button"} hoặc \codeinline{type="reset"} cho tất cả các thẻ \codeinline{<button>}.
    \item \textbf{Ưu tiên dùng \codeinline{<button>} cho giao diện hiện đại}: Giúp dễ dàng thêm icon, văn bản đa dạng và làm chủ thiết kế CSS.
    \item \textbf{Sử dụng thuộc tính \codeinline{disabled} đúng lúc}: Chuyển nút sang trạng thái \codeinline{disabled} khi đang gửi dữ liệu (loading state) để tránh người dùng nhấp trùng lặp (Double submit).
    \item \textbf{Tận dụng thuộc tính \codeinline{formaction}}: Khi form có nhiều hành động gửi khác nhau (ví dụ: Lưu nháp, Đăng ký, Đăng nhập).
    \item \textbf{Đảm bảo Accessibility (a11y)}: Nút bấm phải hiển thị rõ trạng thái Focus khi điều hướng bằng phím Tab và có văn bản mô tả mục đích hành động rõ ràng.
\end{enumerate}
\end{tipBox}

% ------------------------------------------------------------------------------
\section[Thẻ Ô Nhập Số]{Thẻ Ô Nhập Số (<input type="number">)}
\label{sec:html_number}

\begin{conceptBox}[Định Nghĩa Thẻ Ô Nhập Số input type="number"]
Thẻ \codeinline{<input type="number">} là phần tử nhập liệu chuyên biệt cho phép người dùng nhập giá trị số (\term{Integer} hoặc \term{Float}). 
\begin{itemize}[leftmargin=1.5em, itemsep=2pt, parsep=0pt]
    \item \textbf{Giao diện trình duyệt}: Trình duyệt thường hiển thị thêm hai nút mũi tên tăng/giảm (\term{Spinner buttons}) ở cạnh phải để người dùng tăng hoặc giảm số nhanh chóng bằng chuột.
    \item \textbf{Bàn phím di động}: Khi truy cập trên điện thoại (iOS, Android), màn hình sẽ tự động bật bàn phím số (\term{Numeric Keypad}).
    \item \textbf{Giá trị mặc định}: \codeinline{min} và \codeinline{max} không giới hạn, \codeinline{step = 1}, \codeinline{value} mặc định là rỗng (\codeinline{""}).
\end{itemize}
\end{conceptBox}

% ------------------------------------------------------------------------------
\subsection{Cú Pháp Cơ Bản \& Bảng Thuộc Tính Thường Gặp}

\begin{codeWindow}[Cú pháp khởi tạo ô nhập số đơn giản]
\begin{lstlisting}[language=HTML5]
<input type="number">
\end{lstlisting}
\end{codeWindow}

\begin{browserWindow}[Mô Phỏng Trình Duyệt: Ô Nhập Số Cú Pháp Mặc Định]
\sffamily\small
\textbf{Ô nhập số mặc định:} \quad
\tcbox[on line, arc=4pt, colback=white, colframe=primaryBlue!50, boxsep=0pt, left=10pt, right=0pt, top=4pt, bottom=4pt, baseline=-0.6ex]{%
  \textcolor{codeComment}{\small Nhập số...} \hspace{24pt}
  \begin{tikzpicture}[baseline=-0.6ex]
    \fill[primaryBlue!12] (0,-4pt) rectangle (14pt,12pt);
    \draw[primaryBlue!30, line width=0.5pt] (0,-4pt) -- (0,12pt);
    \draw[primaryBlue!30, line width=0.5pt] (0,4pt) -- (14pt,4pt);
    \fill[primaryBlue] (7pt,8.5pt) -- (4pt,5.5pt) -- (10pt,5.5pt) -- cycle;
    \fill[primaryBlue] (7pt,-0.5pt) -- (4pt,2.5pt) -- (10pt,2.5pt) -- cycle;
  \end{tikzpicture}%
} \quad \textcolor{codeComment}{\small\textit{(Nút Spinner tăng/giảm ở cạnh phải)}}
\end{browserWindow}

Dưới đây là bảng tổng hợp đầy đủ các thuộc tính thường dùng cùng thẻ \codeinline{<input type="number">}:

\begin{prettyTableBox}[Bảng Thuộc Tính Thẻ input type="number" Cần Nắm Vững]
\rowcolors{2}{white}{primaryBlue!4}
\begin{tabularx}{\linewidth}{K{2.8cm} Z K{2.5cm}}
\rowcolor{primaryBlue}
\thd{Thuộc Tính} & \thd{Ý Nghĩa} & \thd{Giá Trị Mặc Định} \\
\codeinline{min} & Giá trị nhỏ nhất được phép nhập vào form. & Không giới hạn \\
\codeinline{max} & Giá trị lớn nhất được phép nhập vào form. & Không giới hạn \\
\codeinline{value} & Giá trị mặc định hiển thị khi tải trang. & Rỗng (\codeinline{""}) \\
\codeinline{step} & Bước nhảy tăng/giảm (ví dụ \codeinline{step="2"} chọn số chẵn, \codeinline{step="0.1"} nhập số thập phân). & \codeinline{1} \\
\codeinline{name} & Tên trường dữ liệu gửi đi khi Submit form. & Không có \\
\codeinline{required} & Bắt buộc người dùng phải nhập số mới cho submit. & \codeinline{false} \\
\codeinline{readonly} / \codeinline{disabled} & Không cho sửa hoặc vô hiệu hóa hoàn toàn ô nhập. & \codeinline{false} \\
\end{tabularx}
\end{prettyTableBox}

% ------------------------------------------------------------------------------
\subsection{Ví Dụ Minh Họa Chi Tiết}

\paragraph{a. Ví Dụ Cơ Bản Giới Hạn Khoảng Giá Trị (1--100):}
Ví dụ biểu mẫu yêu cầu người dùng nhập độ tuổi hợp lệ trong khoảng từ 1 đến 100 với giá trị mặc định là 18:

\begin{codeWindow}[Ví dụ giới hạn độ tuổi với min, max và value]
\begin{lstlisting}[language=HTML5]
<form>
  <label for="userAge">(*@Nhập tuổi (1-100):@*)</label>
  <input type="number" id="userAge" name="age" min="1" max="100" value="18">
  <button type="submit">(*@Gửi Đính Kèm@*)</button>
</form>
\end{lstlisting}
\end{codeWindow}

\begin{browserWindow}[Mô Phỏng Trình Duyệt: Biểu Mẫu Nhập Độ Tuổi Hợp Lệ]
\sffamily\small
\textbf{Nhập tuổi (1-100):} \quad
\tcbox[on line, arc=4pt, colback=white, colframe=primaryBlue!60, boxsep=0pt, left=12pt, right=0pt, top=4.5pt, bottom=4.5pt, baseline=-0.6ex]{%
  \textbf{18} \hspace{30pt}
  \begin{tikzpicture}[baseline=-0.6ex]
    \fill[primaryBlue!12] (0,-4pt) rectangle (14pt,12pt);
    \draw[primaryBlue!40, line width=0.5pt] (0,-4pt) -- (0,12pt);
    \draw[primaryBlue!40, line width=0.5pt] (0,4pt) -- (14pt,4pt);
    \fill[primaryBlue] (7pt,8.5pt) -- (4pt,5.5pt) -- (10pt,5.5pt) -- cycle;
    \fill[primaryBlue] (7pt,-0.5pt) -- (4pt,2.5pt) -- (10pt,2.5pt) -- cycle;
  \end{tikzpicture}%
} \qquad
\tcbox[on line, arc=5pt, colback=primaryBlue, colframe=primaryBlue, boxsep=0pt, left=14pt, right=14pt, top=5pt, bottom=5pt, baseline=-0.6ex]{\color{white}\bfseries\small Gửi Dữ Liệu}
\end{browserWindow}

\paragraph{b. Ví Dụ Nâng Cao Nhập Số Thập Phân Với Thuộc Tính step:}
Mặc định \codeinline{step="1"} chỉ cho phép nhập số nguyên. Khi cần nhập số thập phân (như điểm số thang điểm 10.0), ta khai báo \codeinline{step="0.1"}:

\begin{codeWindow}[Ví dụ cho phép nhập số thập phân với step="0.1"]
\begin{lstlisting}[language=HTML5]
<label for="score">(*@Điểm số (thang 0.0 - 10.0):@*)</label>
<input type="number" id="score" name="score" min="0" max="10" step="0.1" value="7.5">
\end{lstlisting}
\end{codeWindow}

\begin{browserWindow}[Mô Phỏng Trình Duyệt: Ô Nhập Điểm Số Thập Phân (step="0.1")]
\sffamily\small
\textbf{Điểm số (thang 0.0 - 10.0):} \quad
\tcbox[on line, arc=4pt, colback=white, colframe=primaryBlue!60, boxsep=0pt, left=12pt, right=0pt, top=4.5pt, bottom=4.5pt, baseline=-0.6ex]{%
  \textbf{7.5} \hspace{30pt}
  \begin{tikzpicture}[baseline=-0.6ex]
    \fill[primaryBlue!12] (0,-4pt) rectangle (14pt,12pt);
    \draw[primaryBlue!40, line width=0.5pt] (0,-4pt) -- (0,12pt);
    \draw[primaryBlue!40, line width=0.5pt] (0,4pt) -- (14pt,4pt);
    \fill[primaryBlue] (7pt,8.5pt) -- (4pt,5.5pt) -- (10pt,5.5pt) -- cycle;
    \fill[primaryBlue] (7pt,-0.5pt) -- (4pt,2.5pt) -- (10pt,2.5pt) -- cycle;
  \end{tikzpicture}%
} \quad \textcolor{secondaryTeal}{\small\textbf{$\leftarrow$ Cho phép chọn 7.5, 7.6...}}
\end{browserWindow}

\paragraph{c. Ví Dụ Custom Giao Diện Nút Tăng/Giảm Số Lượng Chuẩn E-Commerce:}
Trong các website bán hàng, người ta thường ẩn nút spinner mặc định và tạo 2 nút tăng/giảm đẹp mắt bằng HTML/CSS \& JS:

\begin{codeWindow}[Ví dụ custom nút tăng/giảm số lượng sản phẩm chuẩn UX E-Commerce]
\begin{lstlisting}[language=HTML5]
<div class="quantity-picker">
  <button type="button" onclick="qty.value--">-</button>
  <input type="number" id="qty" name="quantity" min="1" max="99" value="18">
  <button type="button" onclick="qty.value++">+</button>
</div>
\end{lstlisting}
\end{codeWindow}

\begin{browserWindow}[Mô Phỏng Trình Duyệt: Custom UI Ô Chọn Số Lượng E-Commerce]
\sffamily\small
\textbf{Số lượng đặt mua:} \quad
\tcbox[on line, arc=6pt, colback=primaryBlue!4, colframe=primaryBlue!30, boxsep=0pt, left=4pt, right=4pt, top=3pt, bottom=3pt, baseline=-0.6ex]{%
  \tcbox[on line, arc=4pt, colback=white, colframe=primaryBlue!40, boxsep=0pt, left=8pt, right=8pt, top=3pt, bottom=3pt, baseline=-0.6ex]{\color{primaryBlue}\bfseries\small $-$} \hspace{10pt}
  \textbf{18} \hspace{10pt}
  \tcbox[on line, arc=4pt, colback=primaryBlue, colframe=primaryBlue, boxsep=0pt, left=8pt, right=8pt, top=3pt, bottom=3pt, baseline=-0.6ex]{\color{white}\bfseries\small $+$}
} \qquad \textcolor{codeComment}{\small\textit{(Thiết kế nút tăng/giảm tùy biến chuẩn E-Commerce UI/UX)}}
\end{browserWindow}

% ------------------------------------------------------------------------------
\subsection{Các Lưu Ý Quan Trọng \& Bẫy Thường Gặp}

\begin{warnBox}[Cảnh Báo Lỗi Kiểm Lỗi Dữ Liệu (Validation Errors)]
\begin{itemize}[leftmargin=1.5em, itemsep=2pt, parsep=0pt]
    \item \textbf{Vượt quá khoảng min/max}: Nếu người dùng cố tình nhập số nhỏ hơn \codeinline{min} hoặc lớn hơn \codeinline{max}, trình duyệt sẽ ngăn Submit form và hiển thị thông báo lỗi HTML5 Validation (trừ khi dùng thuộc tính \codeinline{novalidate}).
    \item \textbf{Lỗi không chia hết cho step}: Nếu giá trị nhập không chia hết cho bội số của \codeinline{step} tính từ mốc \codeinline{min}, form cũng bị báo lỗi không hợp lệ (Ví dụ: \codeinline{step="2"}, người dùng nhập \codeinline{value="3"} $\rightarrow$ Lỗi!).
\end{itemize}
\end{warnBox}

\begin{warnBox}[Bẫy UX Quan Trọng: Khi Nào KHÔNG NÊN Dùng input type="number"?]
\begin{itemize}[leftmargin=1.5em, itemsep=2pt, parsep=0pt]
    \item \textbf{Số điện thoại, Số thẻ ngân hàng, Mã bưu chính (Zip code)}: KHÔNG NÊN dùng \codeinline{type="number"}.
    \item \textbf{Lý do}:
      \begin{enumerate}[leftmargin=1.2em, itemsep=1pt]
        \item Số điện thoại bắt đầu bằng số \codeinline{0} (ví dụ \codeinline{0912345678}) khi ép về kiểu số sẽ bị \textbf{mất số 0 ở đầu} (thành \codeinline{912345678}).
        \item Nút mũi tên tăng/giảm (\term{Spinner}) hoặc thao tác cuộn chuột vô tình làm thay đổi số thẻ tín dụng/số điện thoại của người dùng mà họ không hề hay biết!
      \end{enumerate}
    \item \textbf{Giải pháp Senior}: Sử dụng \codeinline{<input type="tel">} hoặc \codeinline{<input type="text" inputmode="numeric" pattern="[0-9]*">}.
\end{itemize}
\end{warnBox}

\begin{tipBox}[Lưu Ý Thực Chiến Khi Dùng Ô Nhập Số]
\begin{enumerate}[leftmargin=1.5em, itemsep=2pt, parsep=0pt]
    \item Người dùng vẫn có thể gõ các ký tự chữ từ bàn phím trên một số trình duyệt cũ, nhưng khi Submit form, trình duyệt sẽ tự động kích hoạt HTML5 Validation để chặn.
    \item Để cho phép nhập bất kỳ số thập phân nào mà không bị gò bó bởi bội số, bạn có thể gán \codeinline{step="any"}.
    \item Nếu cần kiểm soát logic số phức tạp hơn (ví dụ tính tổng tiền real-time), hãy luôn kết hợp kiểm tra bằng mã JavaScript.
\end{enumerate}
\end{tipBox}

% ------------------------------------------------------------------------------
\subsection{Bảng Tổng Kết \& Ghi Nhớ Cốt Lõi}

\begin{prettyTableBox}[Tóm Tắt Kỹ Thuật Ô Nhập Số]
\rowcolors{2}{white}{primaryBlue!4}
\begin{tabularx}{\linewidth}{K{2.8cm}Z Z}
\rowcolor{primaryBlue}
\thd{Thuộc Tính / Kỹ Thuật} & \thd{Mô Tả} & \thd{Ghi Chú Thực Hành} \\
\codeinline{min / max} & Giới hạn vùng giá trị hợp lệ. & Tránh người dùng nhập số âm hoặc số quá lớn. \\
\codeinline{step} & Quy định khoảng tăng giảm & Dùng \codeinline{0.1} cho số thập phân, \codeinline{any} cho tự do. \\
\codeinline{value} & Giá trị khởi tạo ban đầu. & Nên khai báo sẵn nếu có giá trị gợi ý cho người dùng. \\
\end{tabularx}
\end{prettyTableBox}

% ------------------------------------------------------------------------------
\subsection{Bộ Câu Hỏi Phỏng Vấn Frontend Thực Chiến (Interview Q\&A)}

\begin{noteBox}[Câu Hỏi Phỏng Vấn 1: Nhập Số Thập Phân Trong input type="number"]
\textbf{Hỏi}: Mặc định ô \codeinline{<input type="number">} báo lỗi khi nhập số có dấu phẩy thập phân (như 3.14). Làm sao để cho phép nhập số thập phân?
\par\vspace{3pt}
\textbf{Trả lời}: Do mặc định thuộc tính \codeinline{step} nhận giá trị là \codeinline{1} (chỉ cho nhập số nguyên). Để nhập số thập phân, ta khai báo thuộc tính \codeinline{step="0.1"} (hoặc số lẻ nhỏ hơn tùy độ chính xác) hoặc dùng \codeinline{step="any"} để chấp nhận mọi số thực.
\end{noteBox}

\begin{noteBox}[Câu Hỏi Phỏng Vấn 2: Ẩn Nút Tăng Giảm Spinner Bằng CSS]
\textbf{Hỏi}: Làm thế nào để ẩn hai nút mũi tên tăng/giảm (\term{Spinner buttons}) mặc định của trình duyệt trên ô nhập số?
\par\vspace{3pt}
\textbf{Trả lời}: Ta sử dụng thuộc tính \codeinline{-webkit-appearance: none} cho các trình duyệt WebKit (Chrome, Safari, Edge) và \codeinline{-moz-appearance: textfield} cho Firefox để loại bỏ hoàn toàn nút spinner:
\end{noteBox}

\begin{codeWindow}[Đoạn mã CSS ẩn nút spinner tăng/giảm trên mọi trình duyệt]
\begin{lstlisting}[language=HTML5]
/* (*@Ẩn nút spinner trên Chrome, Safari, Edge@*) */
input[type="number"]::-webkit-outer-spin-button,
input[type="number"]::-webkit-inner-spin-button {
  -webkit-appearance: none;
  margin: 0;
}

/* (*@Ẩn nút spinner trên Firefox@*) */
input[type="number"] {
  -moz-appearance: textfield;
}
\end{lstlisting}
\end{codeWindow}

% ------------------------------------------------------------------------------
\subsection{Đọc Thêm (Mở Rộng Kiến Thức Chuyên Sâu)}
\label{subsec:number_advanced_readmore}

\paragraph{a. Thuộc Tính inputmode="numeric" vs pattern="[0-9]*":}
Khi thiết kế biểu mẫu nhập số điện thoại hoặc số tài khoản ngân hàng, thay vì dùng \codeinline{type="number"} (vốn bị lỗi mất số \codeinline{0} ở đầu), lập trình viên Frontend chuẩn Senior sẽ kết hợp:
\begin{codeWindow}[Tối ưu bàn phím di động chuẩn UX không bị mất số 0]
\begin{lstlisting}[language=HTML5]
<input type="text" inputmode="numeric" pattern="[0-9]*" placeholder="(*@Nhập số tài khoản...@*)">
\end{lstlisting}
\end{codeWindow}

\paragraph{b. Chặn Sự Kiện Wheel (Cuộn Chuột) Làm Đổi Giá Trị Ngoài Ý Muốn:}
Mặc định khi người dùng di chuột qua ô \codeinline{<input type="number">} và cuộn bánh xe chuột (\term{Mouse wheel}), trình duyệt sẽ tự động tăng hoặc giảm số. Để tránh người dùng bị nhảy số vô tình khi lướt trang, ta dùng đoạn mã JavaScript:

\begin{codeWindow}[Chặn cuộn chuột đổi số trên input type="number"]
\begin{lstlisting}[language=JavaScript]
document.querySelectorAll('input[type="number"]').forEach(input => {
  input.addEventListener('wheel', (e) => e.preventDefault());
});
\end{lstlisting}
\end{codeWindow}

% ==============================================================================
% MỤC GIÁO TRÌNH: CHI TIẾT VỀ INPUT TYPE="RANGE" (SLIDER)
% ==============================================================================
\section[Thẻ Ô Chọn Khoảng Giá Trị]{Thẻ Ô Chọn Khoảng Giá Trị (<input type="range">)}
\label{sec:html_range}

\begin{conceptBox}[Định Nghĩa Thẻ Range Input (Thanh Trượt Slider)]
Thẻ \codeinline{<input type="range">} là phần tử biểu mẫu hiển thị dưới dạng một thanh trượt (\term{Slider}), cho phép người dùng lựa chọn một giá trị số nằm trong một khoảng xác định (\term{Range}).
\begin{itemize}[leftmargin=1.5em, itemsep=2pt, parsep=0pt]
    \item \textbf{Cơ chế hiển thị}: Trình duyệt hiển thị một thanh ngang cùng một nút kéo (\term{Thumb}) có thể di chuyển bằng chuột hoặc thao tác vuốt cảm ứng.
    \item \textbf{Bản chất dữ liệu}: Giá trị chọn luôn là một chuỗi số.
    \item \textbf{Giá trị mặc định chuẩn HTML5}:
      \begin{itemize}
        \item \codeinline{min = 0}
        \item \codeinline{max = 100}
        \item \codeinline{value = 50}
      \end{itemize}
    \item \textbf{Công thức tự động tính value mặc định}: Nếu khai báo \codeinline{min} và \codeinline{max} riêng biệt (ví dụ \codeinline{min="20"} và \codeinline{max="200"}) mà không khai báo thuộc tính \codeinline{value}, trình duyệt sẽ tự động tính giá trị mặc định theo công thức trung bình cộng:
    \[
    \text{value}_{\text{mặc định}} = \frac{\text{min} + \text{max}}{2} = \frac{20 + 200}{2} = 110
    \]
\end{itemize}
\end{conceptBox}

% ------------------------------------------------------------------------------
\subsection{Cú Pháp Cơ Bản \& Bảng Thuộc Tính Thường Gặp}

\begin{codeWindow}[Cú pháp khởi tạo thanh trượt Slider đơn giản]
\begin{lstlisting}[language=HTML5]
<input type="range">
\end{lstlisting}
\end{codeWindow}

\begin{browserWindow}[Mô Phỏng Trình Duyệt: Thanh Trượt Slider Mặc Định (Vị trí 50\%)]
\sffamily\small
\textbf{Thanh trượt mặc định (0--100):} \quad
\begin{tikzpicture}[baseline=-0.5ex]
  % Thanh truot nen
  \fill[primaryBlue!20, rounded corners=2pt] (0,-2pt) rectangle (140pt,2pt);
  % Doan da chon 50%
  \fill[primaryBlue, rounded corners=2pt] (0,-2pt) rectangle (70pt,2pt);
  % Nut keo (Thumb)
  \fill[primaryBlue, draw=white, line width=1.5pt] (70pt,0) circle (6pt);
\end{tikzpicture} \quad \textcolor{codeComment}{\small\textit{(Vị trí mặc định: 50)}}
\end{browserWindow}

Dưới đây là bảng tổng hợp các thuộc tính thường dùng cho thanh trượt \codeinline{<input type="range">}:

\begin{prettyTableBox}[Bảng Thuộc Tính Thẻ input type="range" Cần Nắm Vững]
\rowcolors{2}{white}{primaryBlue!4}
\begin{tabularx}{\linewidth}{K{2.8cm} Z K{3.0cm}}
\rowcolor{primaryBlue}
\thd{Thuộc Tính} & \thd{Ý Nghĩa} & \thd{Giá Trị Mặc Định} \\
\codeinline{min} & Giá trị nhỏ nhất ở đầu thanh trượt bên trái. & \codeinline{0} \\
\codeinline{max} & Giá trị lớn nhất ở đầu thanh trượt bên phải. & \codeinline{100} \\
\codeinline{value} & Giá trị khởi tạo của con trỏ Slider. & \codeinline{50} hoặc $\frac{\text{min} + \text{max}}{2}$ \\
\codeinline{step} & Bước nhảy khoảng cách giữa các điểm dừng hợp lệ khi kéo. & \codeinline{1} \\
\codeinline{name} & Tên tham số dữ liệu gửi lên máy chủ khi submit form. & Không có \\
\end{tabularx}
\end{prettyTableBox}

% ------------------------------------------------------------------------------
\subsection{Ví Dụ Minh Họa Chi Tiết}

\paragraph{a. Ví Dụ Biểu Mẫu Điều Chỉnh Âm Lượng (0--100):}
Ví dụ tạo thanh trượt điều chỉnh âm lượng bắt đầu ở mức 30:

\begin{codeWindow}[Ví dụ thanh trượt điều chỉnh âm lượng mức 30]
\begin{lstlisting}[language=HTML5]
<form>
  <label for="vol">(*@Âm lượng (0-100):@*)</label>
  <input type="range" id="vol" name="volume" min="0" max="100" value="30">
</form>
\end{lstlisting}
\end{codeWindow}

\begin{browserWindow}[Mô Phỏng Trình Duyệt: Thanh Trượt Âm Lượng Bắt Đầu Ở Mức 30]
\sffamily\small
\textbf{Âm lượng (0-100):} \quad
\begin{tikzpicture}[baseline=-0.5ex]
  \fill[primaryBlue!20, rounded corners=2pt] (0,-2pt) rectangle (140pt,2pt);
  \fill[primaryBlue, rounded corners=2pt] (0,-2pt) rectangle (42pt,2pt);
  \fill[primaryBlue, draw=white, line width=1.5pt] (42pt,0) circle (6pt);
\end{tikzpicture} \quad \textbf{30}
\end{browserWindow}

\paragraph{b. Ví Dụ Nâng Cao Hiển Thị Giá Trị Thời Gian Thực Với Thẻ <output>:}
Mặc định thanh trượt Slider không tự hiển thị con số bên cạnh. Ta kết hợp sự kiện JavaScript \codeinline{oninput} và thẻ \codeinline{<output>} để hiển thị giá trị thời gian thực khi người dùng kéo con trỏ:

\begin{codeWindow}[Ví dụ hiển thị giá trị kéo Slider thời gian thực với thẻ output]
\begin{lstlisting}[language=HTML5]
<label for="brightness">(*@Độ sáng:@*)</label>
<input type="range" id="brightness" name="brightness" min="0" max="200" step="10" value="100" 
       oninput="output.value = this.value">
<output id="output">100</output>
\end{lstlisting}
\end{codeWindow}

\begin{browserWindow}[Mô Phỏng Trình Duyệt: Thanh Trượt Độ Sáng Cập Nhật Con Số Real-time]
\sffamily\small
\textbf{Độ sáng:} \quad
\begin{tikzpicture}[baseline=-0.5ex]
  \fill[primaryBlue!20, rounded corners=2pt] (0,-2pt) rectangle (140pt,2pt);
  \fill[primaryBlue, rounded corners=2pt] (0,-2pt) rectangle (70pt,2pt);
  \fill[primaryBlue, draw=white, line width=1.5pt] (70pt,0) circle (6pt);
\end{tikzpicture} \quad
\tcbox[on line, arc=3pt, colback=primaryBlue!10, colframe=primaryBlue, boxsep=1pt, left=6pt, right=6pt, top=2pt, bottom=2pt]{\color{primaryBlue}\bfseries\small 100}
\end{browserWindow}

\paragraph{c. Kỹ Thuật Nâng Cao: Tạo Các Nấc Điểm Dừng (Tick Marks) Với Thẻ <datalist>:}
Chuẩn HTML5 cho phép kết nối thanh trượt Slider với thẻ \codeinline{<datalist>} để hiển thị các vạch nấc dừng (\term{Tick marks / Snap points}) trực quan:

\begin{codeWindow}[Ví dụ tạo các nấc điểm dừng trên Slider với datalist]
\begin{lstlisting}[language=HTML5]
<label for="temp">(*@Nhiệt độ phòng (°C):@*)</label>
<input type="range" id="temp" name="temp" min="0" max="40" step="10" list="ticks">

<!-- (*@Khai báo danh sách các vạch mốc nấc dừng@*) -->
<datalist id="ticks">
  <option value="0" label="0 deg C"></option>
  <option value="10"></option>
  <option value="20" label="20 deg C"></option>
  <option value="30"></option>
  <option value="40" label="40 deg C"></option>
</datalist>
\end{lstlisting}
\end{codeWindow}

\begin{browserWindow}[Mô Phỏng Trình Duyệt: Slider Có Các Vạch Nấc Dừng (Tick Marks)]
\sffamily\small
\textbf{Nhiệt độ (°C):} \quad
\begin{tikzpicture}[baseline=-0.5ex]
  \fill[primaryBlue!20, rounded corners=2pt] (0,-2pt) rectangle (140pt,2pt);
  \fill[primaryBlue, rounded corners=2pt] (0,-2pt) rectangle (70pt,2pt);
  \fill[primaryBlue, draw=white, line width=1.5pt] (70pt,0) circle (6pt);
  % Vach nac dung (Ticks)
  \foreach \x in {0, 35, 70, 105, 140} {
    \draw[primaryBlue!60, line width=1pt] (\x pt, -6pt) -- (\x pt, -3pt);
  }
\end{tikzpicture} \quad \textcolor{codeComment}{\small\textit{(Các vạch mốc: 0, 10, 20, 30, 40)}}
\end{browserWindow}

% ------------------------------------------------------------------------------
\subsection{Các Lưu Ý Quan Trọng \& Bẫy Thường Gặp}

\begin{tipBox}[4 Quy Tắc Cốt Lõi Cần Nhớ Về Thẻ Range]
\begin{enumerate}[leftmargin=1.5em, itemsep=2pt, parsep=0pt]
    \item \textbf{Tự động sửa giá trị (Clamping)}: Nếu gán thuộc tính \codeinline{value} nhỏ hơn \codeinline{min} hoặc lớn hơn \codeinline{max}, trình duyệt sẽ tự động điều chỉnh con trỏ về mốc giới hạn gần nhất.
    \item \textbf{Tác dụng của step}: Thuộc tính \codeinline{step} giúp hạn chế các mốc số dư (ví dụ \codeinline{step="10"} chỉ cho phép dừng tại 0, 10, 20, 30...).
    \item \textbf{Dữ liệu gửi trong Form}: Khi Submit form, trình duyệt chỉ gửi \textbf{giá trị số duy nhất} (ví dụ \codeinline{brightness=100}), không gửi tọa độ con trỏ hay tỷ lệ phần trăm.
    \item \textbf{Trải nghiệm người dùng (UX)}: Nên luôn đi kèm thẻ \codeinline{<output>} hoặc văn bản hiển thị con số để người dùng biết chính xác giá trị họ đang chọn.
\end{enumerate}
\end{tipBox}

% ------------------------------------------------------------------------------
\subsection{Bảng Tổng Kết \& Ghi Nhớ Cốt Lõi}

\begin{prettyTableBox}[Tóm Tắt Thuộc Tính \& Tính Toán Mặc Định Thẻ Range]
\rowcolors{2}{white}{primaryBlue!4}
\begin{tabularx}{\linewidth}{K{3.0cm}Z Z}
\rowcolor{primaryBlue}
\thd{Trường Hợp} & \thd{Giá Trị Mặc Định} & \thd{Ví Dụ Minh Họa} \\
Không khai báo gì & \codeinline{min=0}, \codeinline{max=100}, \codeinline{value=50} & \codeinline{<input type="range">} $\rightarrow$ Vị trí 50. \\
Khai báo min, max & \codeinline{value} = \codeinline{(min + max) / 2} & \codeinline{min="20" max="200"} $\rightarrow$ \codeinline{value = 110}. \\
\end{tabularx}
\end{prettyTableBox}

% ------------------------------------------------------------------------------
\subsection{Bộ Câu Hỏi Phỏng Vấn Frontend Thực Chiến (Interview Q\&A)}

\begin{noteBox}[Câu Hỏi Phỏng Vấn 1: Cập Nhật Con Số Real-time Khi Kéo Slider]
\textbf{Hỏi}: Làm sao để hiển thị con số thay đổi liên tục theo thời gian thực khi người dùng đang kéo thanh trượt Slider?
\par\vspace{3pt}
\textbf{Trả lời}: Ta lắng nghe sự kiện \codeinline{oninput} (hoặc sự kiện \codeinline{input} trong JavaScript) trên thẻ range và gán giá trị \codeinline{this.value} vào một thẻ \codeinline{<output>} hoặc phần tử hiển thị văn bản. Không nên dùng sự kiện \codeinline{change} vì \codeinline{change} chỉ nhả ra kết quả khi người dùng đã thả tay khỏi chuột.
\end{noteBox}

\begin{noteBox}[Câu Hỏi Phỏng Vấn 2: Tùy Biến Giao Diện Thanh Trượt Bằng CSS]
\textbf{Hỏi}: Giao diện thanh trượt Slider mặc định trên các hệ điều hành trông rất thô. Làm sao để custom màu sắc và nút kéo bằng CSS?
\par\vspace{3pt}
\textbf{Trả lời}: Ta sử dụng thuộc tính \codeinline{appearance: none} để loại bỏ giao diện mặc định, sau đó tùy biến đường rãnh (\term{Track}) và nút kéo (\term{Thumb}) thông qua các Pseudo-elements của từng trình duyệt:
\begin{itemize}[leftmargin=1.5em, itemsep=2pt, parsep=0pt]
    \item Chrome / Safari / Edge: \codeinline{::-webkit-slider-runnable-track} và \codeinline{::-webkit-slider-thumb}
    \item Firefox: \codeinline{::-moz-range-track} và \codeinline{::-moz-range-thumb}
\end{itemize}
\end{noteBox}

% ------------------------------------------------------------------------------
\subsection{Đọc Thêm (Mở Rộng Kiến Thức Chuyên Sâu)}
\label{subsec:range_advanced_readmore}

\paragraph{a. Thiết Lập Thanh Trượt Dọc (Vertical Range Slider):}
Mặc định thanh trượt Slider hiển thị theo chiều ngang. Khi thiết kế equalizer âm thanh hoặc cột đo chiều cao, ta có thể xoay thanh trượt thành chiều dọc bằng thuộc tính CSS \codeinline{writing-mode}:

\begin{codeWindow}[Thiết lập thanh trượt Slider xoay dọc]
\begin{lstlisting}[language=HTML5]
<input type="range" id="volVertical" min="0" max="100" value="70" 
       style="writing-mode: vertical-lr; direction: rtl; width: 16px; height: 120px;">
\end{lstlisting}
\end{codeWindow}

\paragraph{b. Kỹ Thuật Đổi Màu Tiến Trình Kéo Slider (Dynamic Gradient Range Fill):}
Để làm nổi bật phần thanh trượt đã kéo qua (màu xanh brand) khác với phần chưa kéo (màu xám nhạt), ta lắng nghe sự kiện \codeinline{input} trong JavaScript để cập nhật dải dốc màu nền \codeinline{linear-gradient} tương ứng với tỷ lệ phần trăm:

\begin{codeWindow}[Custom hiệu ứng màu tiến trình kéo Slider như Spotify]
\begin{lstlisting}[language=JavaScript]
const rangeInput = document.querySelector('input[type="range"]');
rangeInput.addEventListener('input', function() {
  const percent = this.value - this.min) / (this.max - this.min * 100;
  // (*@Cập nhật nền gradient động theo tỷ lệ phần trăm@*)
  this.style.background = `linear-gradient(to right, #0056b3 ${percent}%, #e0e0e0 ${percent}%)`;
});
\end{lstlisting}
\end{codeWindow}

% ==============================================================================
% MỤC GIÁO TRÌNH: CHI TIẾT VỀ THẺ TEXTAREA trong HTML (<textarea>)
% ==============================================================================
\section[Thẻ Ô Nhập Văn Bản Nhiều Dòng Textarea]{Thẻ Ô Nhập Văn Bản Nhiều Dòng Textarea (<textarea>)}
\label{sec:html_textarea}

\begin{conceptBox}[Định Nghĩa Thẻ Textarea]
\term{Textarea} là thẻ HTML chuyên dụng dùng để tạo một ô nhập văn bản nhiều dòng (\term{Multi-line text input field}). Khác với thẻ \codeinline{<input type="text">} vốn chỉ cho phép nhập văn bản trên một dòng đơn lẻ, \codeinline{<textarea>} cho phép người dùng xuống dòng và nhập các đoạn văn bản dài.
\begin{itemize}[leftmargin=1.5em, itemsep=2pt, parsep=0pt]
    \item \textbf{Ứng dụng phổ biến}: Nhập ý kiến phản hồi (\term{Feedback}), bình luận bài viết (\term{Comment}), mô tả sản phẩm (\term{Description}), ghi chú, hoặc nhập địa chỉ giao hàng (\term{Address}).
    \item \textbf{Khi nào nên dùng}: Sử dụng khi cần thu thập dữ liệu dạng văn bản dài hoặc dữ liệu có cấu trúc nhiều dòng thay vì chuỗi ký tự ngắn.
\end{itemize}
\end{conceptBox}

% ------------------------------------------------------------------------------
\subsection{Cú Pháp Cơ Bản \& Cấu Trúc Khai Báo}

\begin{codeWindow}[Cú pháp cơ bản tạo ô nhập liệu Textarea]
\begin{lstlisting}[language=HTML5]
<label for="msg">Tin nhan:</label>
<textarea id="msg" name="message"></textarea>
\end{lstlisting}
\end{codeWindow}

\begin{browserWindow}[Kết Quả Hiển Thị Trình Duyệt: Textarea Cú Pháp Cơ Bản]
\sffamily\small
\textbf{Tin nhắn:} \\[4pt]
\tcbox[arc=2pt, colback=white, colframe=primaryBlue!40, boxsep=2pt, width=0.6\linewidth, height=45pt]{
  \color{gray!60}\small\textit{(Người dùng có thể gõ nhiều dòng văn bản vào đây...)}
}
\end{browserWindow}

\begin{tipBox}[Giải Thích Thành Phần Cấu Trúc]
\begin{itemize}[leftmargin=1.2em, itemsep=2pt, parsep=0pt]
    \item \codeinline{name="message"}: Tên trường dữ liệu đóng gói gửi về máy chủ khi submit form. Nếu không khai báo \codeinline{name}, dữ liệu trong ô sẽ không được gửi đi.
    \item \codeinline{id="msg"}: Định danh duy nhất kết nối logic với thuộc tính \codeinline{for="msg"} của thẻ \codeinline{<label>}.
\end{itemize}
\end{tipBox}

% ------------------------------------------------------------------------------
\subsection{Các Thuộc Tính Quan Trọng Của Thẻ <textarea>}

Dưới đây là bảng tổng hợp chi tiết tất cả các thuộc tính quan trọng thường được sử dụng cùng thẻ \codeinline{<textarea>}:

\begin{prettyTableBox}[Bảng Tổng Hợp Thuộc Tính Của Thẻ Textarea]
\rowcolors{2}{white}{primaryBlue!4}
\begin{tabularx}{\linewidth}{K{2.8cm} Z K{2.2cm} Z}
\rowcolor{primaryBlue}
\thd{Thuộc Tính} & \thd{Ý Nghĩa \& Chức Năng} & \thd{Giá Trị Mặc Định} & \thd{Ví Dụ Cú Pháp} \\
\codeinline{name} & Tên trường dữ liệu gửi lên máy chủ khi submit form. & Không có & \codeinline{name="comment"} \\
\codeinline{id} & Định danh duy nhất để liên kết với thẻ \codeinline{<label>}. & Không có & \codeinline{id="commentMsg"} \\
\codeinline{rows} & Số lượng hàng (chiều cao) hiển thị mặc định của ô nhập. & 2 dòng & \codeinline{rows="5"} \\
\codeinline{cols} & Số lượng cột (chiều rộng ký tự) hiển thị mặc định trên mỗi dòng. & 20 ký tự & \codeinline{cols="40"} \\
\codeinline{placeholder} & Đoạn văn bản gợi ý nội dung ẩn hiển thị khi ô nhập rỗng. & Không có & \codeinline{placeholder="Nhập nội dung..."} \\
\codeinline{maxlength} & Giới hạn số lượng ký tự tối đa cho phép nhập vào ô. & Không giới hạn & \codeinline{maxlength="200"} \\
\codeinline{readonly} & Thiết lập ô chỉ đọc, người dùng không thể chỉnh sửa nội dung. & \codeinline{false} & \codeinline{readonly} \\
\codeinline{disabled} & Vô hiệu hóa ô nhập, không thể tương tác và không gửi dữ liệu. & \codeinline{false} & \codeinline{disabled} \\
\codeinline{required} & Bắt buộc người dùng phải điền dữ liệu trước khi submit form. & \codeinline{false} & \codeinline{required} \\
\codeinline{wrap} & Quy định cơ chế xử lý xuống dòng khi gửi form (\codeinline{soft}, \codeinline{hard}, \codeinline{off}). & \codeinline{soft} & \codeinline{wrap="soft"} \\
\end{tabularx}
\end{prettyTableBox}

% ------------------------------------------------------------------------------
\subsection{Các Giá Trị Mặc Định Của <textarea> Trong HTML}

Khi không chỉ định bất kỳ thuộc tính nào, trình duyệt web sẽ áp dụng bộ giá trị mặc định tiêu chuẩn sau cho thẻ \codeinline{<textarea>}:

\begin{prettyTableBox}[Bảng Chi Tiết Giá Trị Mặc Định Của Thẻ Textarea]
\rowcolors{2}{white}{primaryBlue!4}
\begin{tabularx}{\linewidth}{K{2.8cm} K{3.0cm} Z}
\rowcolor{primaryBlue}
\thd{Thuộc Tính} & \thd{Giá Trị Mặc Định} & \thd{Giải Thích Hành Vi Trình Duyệt} \\
Nội dung bên trong & Rỗng (\codeinline{""}) & Thẻ \codeinline{<textarea></textarea>} không chứa bất kỳ văn bản nào. \\
\codeinline{rows} & \codeinline{2} & Trình duyệt mặc định hiển thị chiều cao tương đương 2 dòng văn bản. \\
\codeinline{cols} & \codeinline{20} & Trình duyệt mặc định hiển thị chiều rộng tương đương 20 ký tự trên mỗi dòng. \\
\codeinline{wrap} & \codeinline{soft} & Văn bản tự động hiển thị xuống dòng trên giao diện, nhưng khi gửi về server chỉ có ký tự xuống dòng do người dùng nhấn Enter thực sự mới được đính kèm. \\
\codeinline{spellcheck} & \codeinline{true} & Trình duyệt tự động bật tính năng kiểm tra chính tả (nếu trình duyệt có hỗ trợ ngôn ngữ). \\
\codeinline{resize} & \codeinline{both} (\textit{CSS}) & Cho phép người dùng kéo thả chuột để thay đổi kích thước ô nhập theo cả chiều ngang và chiều dọc. \\
\codeinline{maxlength} & Không giới hạn & Người dùng có thể nhập văn bản với độ dài tùy ý. \\
\codeinline{placeholder} & Không có & Không hiển thị văn bản mờ gợi ý khi ô nhập trống. \\
\codeinline{readonly} & \codeinline{false} & Người dùng có thể thoải mái nhập và chỉnh sửa nội dung. \\
\codeinline{disabled} & \codeinline{false} & Ô nhập hoạt động bình thường và sẵn sàng thu nhận tương tác. \\
\codeinline{required} & \codeinline{false} & Biểu mẫu không bắt buộc phải điền ô này mới cho phép submit. \\
\codeinline{autofocus} & \codeinline{false} & Không tự động đặt con trỏ chuột vào ô nhập khi tải xong trang web. \\
\end{tabularx}
\end{prettyTableBox}

% ------------------------------------------------------------------------------
\subsection{So Sánh Textarea Mặc Định vs Textarea Cấu Hình Thuộc Tính}

Để thấy rõ sự khác biệt giữa cài đặt mặc định và khi khai báo thuộc tính, hãy xem ví dụ so sánh dưới đây:

\begin{codeWindow}[Mã nguồn so sánh Textarea mặc định và có thuộc tính]
\begin{lstlisting}[language=HTML5]
<!-- 1. Textarea mac dinh -->
<label>Textarea mac dinh:</label>
<textarea></textarea>

<!-- 2. Textarea co thuoc tinh -->
<label>Textarea co thuoc tinh:</label>
<textarea rows="5" cols="40" placeholder="(*@Nhập ghi chú tại đây...@*)"></textarea>
\end{lstlisting}
\end{codeWindow}

\begin{browserWindow}[Mô Phỏng Trình Duyệt: So Sánh Mặc Định vs Có Thuộc Tính]
\sffamily\small
\textbf{1. Textarea mặc định (2 dòng $\times$ 20 ký tự):} \par\vspace{4pt}
\tcbox[arc=2pt, colback=white, colframe=primaryBlue!40, boxsep=2pt, width=0.45\linewidth, height=30pt]{} \par\vspace{8pt}

\textbf{2. Textarea có thuộc tính (5 dòng $\times$ 40 ký tự + Placeholder):} \par\vspace{4pt}
\tcbox[arc=2pt, colback=white, colframe=primaryBlue!40, boxsep=2pt, width=0.75\linewidth, height=65pt]{
  \color{gray!70}\small\textit{Nhập ghi chú tại đây...}
}
\end{browserWindow}

% ------------------------------------------------------------------------------
\subsection{Các Ví Dụ Minh Họa Chi Tiết Thực Tế}

\paragraph{Ví Dụ 1: Ô nhập phản hồi có giới hạn số dòng, cột và tối đa ký tự:}

\begin{codeWindow}[Form nhập phản hồi giới hạn 4 dòng, 50 cột và tối đa 200 ký tự]
\begin{lstlisting}[language=HTML5]
<textarea name="feedback" rows="4" cols="50" maxlength="200" 
          placeholder="(*@Nhập phản hồi của bạn...@*)"></textarea>
\end{lstlisting}
\end{codeWindow}

\begin{browserWindow}[Kết Quả Trình Duyệt: Ô Nhập Phản Hồi Giới Hạn 200 Ký Tự]
\sffamily\small
\tcbox[arc=2pt, colback=white, colframe=primaryBlue!50, boxsep=2pt, width=0.8\linewidth, height=55pt]{
  \color{gray!70}\small\textit{Nhập phản hồi của bạn...}
} \par\vspace{4pt}
\textcolor{codeComment}{\small\textit{$\rightarrow$ Hiển thị ô nhập rộng 4 dòng $\times$ 50 ký tự mỗi dòng, giới hạn tối đa 200 ký tự.}}
\end{browserWindow}

\paragraph{Ví Dụ 2: Ô nhập chứa văn bản giá trị mặc định sẵn có:}

\begin{codeWindow}[Form với nội dung mặc định đặt giữa thẻ mở và đóng]
\begin{lstlisting}[language=HTML5]
<textarea name="bio" rows="3" cols="40">(*@Tôi là sinh viên CNTT.@*)</textarea>
\end{lstlisting}
\end{codeWindow}

\begin{browserWindow}[Kết Quả Trình Duyệt: Textarea Có Giá Trị Mặc Định]
\sffamily\small
\tcbox[arc=2pt, colback=white, colframe=primaryBlue!50, boxsep=2pt, width=0.7\linewidth, height=45pt]{
  \small\textcolor{black}{Tôi là sinh viên CNTT.}
} \par\vspace{4pt}
\textcolor{codeComment}{\small\textit{$\rightarrow$ Khi tải trang, văn bản "Tôi là sinh viên CNTT." đã xuất hiện sẵn trong ô nhập.}}
\end{browserWindow}

\paragraph{Ví Dụ 3: Ô nhập trạng thái chỉ đọc (readonly):}

\begin{codeWindow}[Form cấu hình thuộc tính readonly không cho sửa nội dung]
\begin{lstlisting}[language=HTML5]
<textarea name="info" rows="3" cols="40" readonly>(*@Thông tin chỉ đọc, không thể sửa...@*)</textarea>
\end{lstlisting}
\end{codeWindow}

\begin{browserWindow}[Kết Quả Trình Duyệt: Textarea Trạng Thái Readonly]
\sffamily\small
\tcbox[arc=2pt, colback=gray!15, colframe=gray!40, boxsep=2pt, width=0.7\linewidth, height=45pt]{
  \small\color{darkText}Thông tin chỉ đọc, không thể sửa...
} \par\vspace{4pt}
\textcolor{codeComment}{\small\textit{$\rightarrow$ Người dùng có thể bôi đen copy nhưng hoàn toàn KHÔNG THỂ chỉnh sửa hay xóa văn bản.}}
\end{browserWindow}

\paragraph{Ví Dụ 4: Ô nhập trạng thái bị vô hiệu hóa (disabled):}

\begin{codeWindow}[Form cấu hình thuộc tính disabled vô hiệu hóa tương tác]
\begin{lstlisting}[language=HTML5]
<textarea name="note" rows="3" cols="40" disabled>(*@Không thể nhập dữ liệu vào đây...@*)</textarea>
\end{lstlisting}
\end{codeWindow}

\begin{browserWindow}[Kết Quả Trình Duyệt: Textarea Trạng Thái Disabled]
\sffamily\small
\tcbox[arc=2pt, colback=gray!25, colframe=gray!30, boxsep=2pt, width=0.7\linewidth, height=45pt]{
  \small\color{gray!80}Không thể nhập dữ liệu vào đây...
} \par\vspace{4pt}
\textcolor{codeComment}{\small\textit{$\rightarrow$ Ô nhập bị mờ đi, không thể tương tác và dữ liệu sẽ KHÔNG ĐƯỢC gửi đi khi submit form.}}
\end{browserWindow}

\paragraph{Ví Dụ 5: Ô nhập bắt buộc điền dữ liệu (required):}

\begin{codeWindow}[Form bắt buộc nhập bình luận trước khi submit]
\begin{lstlisting}[language=HTML5]
<textarea name="comment" rows="4" cols="40" required 
          placeholder="(*@Vui lòng nhập bình luận...@*)"></textarea>
\end{lstlisting}
\end{codeWindow}

% ------------------------------------------------------------------------------
\subsection{Các Lưu Ý Quan Trọng Cần Ghi Nhớ Nằm Lòng}

\begin{warnBox}[Các Lưu Ý Bắt Buộc Phải Thuộc Khi Làm Việc Với Thẻ Textarea]
\begin{enumerate}[leftmargin=1.5em, itemsep=6pt, parsep=0pt]
    \item \textbf{CỰC KỲ QUAN TRỌNG: Không dùng thuộc tính \codeinline{value="..."}}:
    Không giống như các thẻ \codeinline{<input>}, thẻ \codeinline{<textarea>} \textbf{HOÀN TOÀN KHÔNG HỖ TRỢ thuộc tính \codeinline{value}}. Muốn đặt giá trị văn bản mặc định khi tải trang, bạn \textbf{bắt buộc phải đặt văn bản nằm giữa cặp thẻ mở và thẻ đóng}:
    \par\vspace{4pt}
    \begin{center}
    \tcbox[on line, arc=3pt, colback=primaryBlue!8, colframe=primaryBlue!40, boxsep=3pt, left=10pt, right=10pt, top=4pt, bottom=4pt]{%
      \small\ttfamily\bfseries \textcolor{primaryBlue}{<textarea>}\textcolor{black}{Văn bản giá trị mặc định}\textcolor{primaryBlue}{</textarea>}%
    }
    \end{center}
    \par\vspace{2pt}
    \textcolor{red!80!black}{\textbf{Cảnh báo lỗi thường gặp}}: Cú pháp \tcbox[on line, arc=2pt, colback=red!10, colframe=red!30, boxsep=1.5pt, left=4pt, right=4pt, top=1.5pt, bottom=1.5pt]{\small\ttfamily <textarea value="abc">} là \textbf{SAI HOÀN TOÀN} và sẽ bị trình duyệt bỏ qua!

    \item \textbf{Cấu hình kích thước với CSS (\codeinline{resize})}:
    Mặc định, trình duyệt cho phép người dùng kéo chuột ở góc dưới ô để thay đổi kích thước (\codeinline{resize: both}). Trong thiết kế thực tế, để tránh làm phá vỡ giao diện web (\term{Layout Break}), lập trình viên thường dùng CSS vô hiệu hóa tính năng này hoặc chỉ cho phép kéo theo chiều dọc:
    \par\vspace{4pt}
    \begin{center}
    \begin{tcolorbox}[arc=3pt, colback=white, colframe=primaryBlue!30, boxsep=5pt, left=12pt, right=12pt, top=5pt, bottom=5pt, width=0.85\linewidth]
      \small\ttfamily
      \textcolor{codeComment}{/* Không cho phép người dùng thay đổi kích thước */}\\
      \textcolor{primaryBlue}{\textbf{textarea}} \{ \textcolor{codeComment}{resize}: \textcolor{red!80!black}{none}; \}\\[6pt]
      \textcolor{codeComment}{/* Chỉ cho phép thay đổi kích thước theo chiều dọc */}\\
      \textcolor{primaryBlue}{\textbf{textarea}} \{ \textcolor{codeComment}{resize}: \textcolor{red!80!black}{vertical}; \}
    \end{tcolorbox}
    \end{center}
    \par\vspace{2pt}

    \item \textbf{Tự động xuất hiện thanh cuộn (\codeinline{overflow})}:
    Khi nội dung văn bản do người dùng nhập vào vượt quá số dòng (\codeinline{rows}) hoặc số cột (\codeinline{cols}) hiển thị, trình duyệt sẽ tự động thêm thanh cuộn đứng (\term{Vertical Scrollbar}) để cho phép cuộn xem toàn bộ văn bản.

    \item \textbf{Bắt buộc phải có thuộc tính \codeinline{name}}:
    Nếu bạn muốn thu thập và gửi dữ liệu văn bản trong \codeinline{<textarea>} về máy chủ Backend, thẻ \codeinline{<textarea>} bắt buộc phải có thuộc tính \codeinline{name}. Nếu thiếu \codeinline{name}, dữ liệu sẽ bị bỏ qua khi submit form.
\end{enumerate}
\end{warnBox}

% ------------------------------------------------------------------------------
\subsection{Đọc Thêm: Kỹ Thuật Nâng Cao \& Tối Ưu Trải Nghiệm Lập Trình Textarea}

\begin{tipBox}[Tự Động Giãn Chiều Cao Textarea Theo Nội Dung Nhập (Auto-expand)]
Trong các ứng dụng nhắn tin hoặc mạng xã hội hiện đại (như Facebook, Messenger, Zalo), ô nhập văn bản thường tự động cao lên khi người dùng gõ xuống dòng mà không xuất hiện thanh cuộn xấu xí. Ta có thể đạt được hiệu ứng này bằng 1 dòng JavaScript đơn giản:
\begin{codeWindow}[Mã JavaScript tự động thay đổi chiều cao Textarea theo scrollHeight]
\begin{lstlisting}[language=JavaScript]
const textarea = document.querySelector('textarea');
textarea.addEventListener('input', function() {
  this.style.height = 'auto'; // Reset chieu cao ban dau
  this.style.height = (this.scrollHeight) + 'px'; // Cap nhat chieu cao theo noi dung
});
\end{lstlisting}
\end{codeWindow}
\end{tipBox}

\begin{noteBox}[Kỹ Thuật Đếm Số Ký Tự Còn Lại (Real-time Character Counter)]
Kết hợp thuộc tính \codeinline{maxlength="200"} với JavaScript để hiển thị bộ đếm số ký tự thời gian thực giúp người dùng biết mình còn bao nhiêu ký tự được phép gõ:
\begin{codeWindow}[Mã nguồn đếm ký tự thời gian thực cho Textarea]
\begin{lstlisting}[language=HTML5]
<textarea id="myText" maxlength="200" placeholder="Nhap noi dung..."></textarea>
<p>So ky tu con lai: <span id="charCount">200</span>/200</p>

<script>
const myText = document.getElementById('myText');
const charCount = document.getElementById('charCount');

myText.addEventListener('input', function() {
  const remaining = 200 - this.value.length;
  charCount.textContent = remaining;
});
</script>
\end{lstlisting}
\end{codeWindow}
\end{noteBox}

% ==============================================================================
% MỤC GIÁO TRÌNH: CHI TIẾT VỀ THẺ SELECT & OPTION trong HTML (<select>, <option>, <optgroup>)
% ==============================================================================
\section[Thẻ Danh Sách Thả Xuống Select]{Thẻ Danh Sách Thả Xuống Select (<select>, <option> \& <optgroup>)}
\label{sec:html_select}

\begin{conceptBox}[Định Nghĩa Thẻ Select \& Option]
Thẻ \codeinline{<select>} trong HTML được sử dụng để tạo một \textbf{danh sách thả xuống} (\term{Dropdown List}) hoặc một danh sách chọn (\term{List Box}), cho phép người dùng lựa chọn một hoặc nhiều tùy chọn từ danh sách các mục có sẵn.
\begin{itemize}[leftmargin=1.5em, itemsep=2pt, parsep=0pt]
    \item \textbf{Thẻ \codeinline{<option>}}: Được đặt bên trong thẻ \codeinline{<select>} để định nghĩa từng mục lựa chọn cụ thể trong danh sách.
    \item \textbf{Thẻ \codeinline{<optgroup>}}: Dùng để phân nhóm các tùy chọn có cùng chủ đề liên quan, giúp giao diện danh sách ngắn gọn và dễ tìm kiếm hơn.
\end{itemize}

\tcblower
\textbf{ Lợi Ích Cốt Lõi Khi Sử Dụng Thẻ Select}:
\begin{enumerate}[leftmargin=1.2em, itemsep=2pt, parsep=0pt]
    \item \textbf{Tiết kiệm không gian giao diện}: Danh sách ẩn đi và chỉ mở ra khi người dùng kích chuột vào.
    \item \textbf{Tránh nhập sai dữ liệu}: Ép người dùng chỉ được phép chọn từ danh sách chuẩn cố định (ví dụ: Chọn Tỉnh/Thành phố, Quốc gia, Giới tính).
    \item \textbf{Dễ quản lý và mở rộng}: Dễ dàng bổ sung hoặc chỉnh sửa danh sách các tùy chọn trên Frontend hoặc lấy từ Cơ sở dữ liệu Backend.
\end{enumerate}
\end{conceptBox}

% ------------------------------------------------------------------------------
\subsection{Cú Pháp Cơ Bản \& Cấu Trúc Khai Báo}

\begin{codeWindow}[Cú pháp cơ bản tạo danh sách thả xuống Select]
\begin{lstlisting}[language=HTML5]
<label for="country">Chon quoc gia:</label>
<select id="country" name="country">
  <option value="vn">(*@Việt@*) Nam@*)</option>
  <option value="us">(*@Hoa (*@Kỳ@*)</option>
  <option value="jp">(*@Nhật *@Bản@*)</option>
</select>
\end{lstlisting}
\end{codeWindow}

\begin{browserWindow}[Mô Phỏng Trình Duyệt: Danh Sách Thả Xuống Select Cơ Bản]
\sffamily\small
\textbf{Chọn quốc gia:} \par\vspace{4pt}
\tcbox[on line, arc=3pt, colback=white, colframe=primaryBlue!50, boxsep=3pt, left=8pt, right=20pt, top=3pt, bottom=3pt]{%
  \small \textcolor{black}{Việt Nam} \;\; \textcolor{primaryBlue}{\textbf{$\boldsymbol{\vee}$}}%
}
\end{browserWindow}

% ------------------------------------------------------------------------------
\subsection{Các Thuộc Tính Quan Trọng Của Thẻ <select>}

Dưới đây là bảng tổng hợp chi tiết tất cả các thuộc tính quan trọng của thẻ \codeinline{<select>}:

\begin{prettyTableBox}[Bảng Tổng Hợp Thuộc Tính Của Thẻ Select]
\rowcolors{2}{white}{primaryBlue!4}
\begin{tabularx}{\linewidth}{K{2.4cm} Z K{2.5cm} Z}
\rowcolor{primaryBlue}
\thd{Thuộc Tính} & \thd{Ý Nghĩa \& Chức Năng} & \thd{Giá Trị Mặc Định} & \thd{Ví Dụ Cú Pháp} \\
\codeinline{name} & Tên của trường dữ liệu để đóng gói gửi về máy chủ Backend. & Không có & \codeinline{name="country"} \\
\codeinline{id} & Định danh duy nhất để liên kết với thuộc tính \codeinline{for} của thẻ \codeinline{<label>}. & Không có & \codeinline{id="countrySelect"} \\
\codeinline{multiple} & Cho phép người dùng chọn nhiều tùy chọn cùng lúc (dùng Ctrl/Shift hoặc Cmd). & \codeinline{false} & \codeinline{multiple} \\
\codeinline{size} & Số lượng tùy chọn hiển thị trực tiếp cùng lúc (chuyển sang dạng hộp danh sách Listbox). & 1 (hiển thị dạng Dropdown) & \codeinline{size="3"} \\
\codeinline{disabled} & Vô hiệu hóa toàn bộ danh sách select, người dùng không thể tương tác. & \codeinline{false} & \codeinline{disabled} \\
\codeinline{required} & Bắt buộc người dùng phải chọn một tùy chọn trước khi submit form. & \codeinline{false} & \codeinline{required} \\
\codeinline{autofocus} & Tự động đặt con trỏ tiêu điểm vào danh sách khi trang vừa tải xong. & \codeinline{false} & \codeinline{autofocus} \\
\end{tabularx}
\end{prettyTableBox}

% ------------------------------------------------------------------------------
\subsection{Các Thuộc Tính Quan Trọng Của Thẻ <option> Bên Trong}

Mỗi thẻ \codeinline{<option>} đại diện cho một mục lựa chọn và sở hữu các thuộc tính sau:

\begin{prettyTableBox}[Bảng Thuộc Tính Của Thẻ Option]
\rowcolors{2}{white}{primaryBlue!4}
\begin{tabularx}{\linewidth}{K{2.4cm} Z K{2.2cm} Z}
\rowcolor{primaryBlue}
\thd{Thuộc Tính} & \thd{Ý Nghĩa \& Chức Năng} & \thd{Giá Trị Mặc Định} & \thd{Ví Dụ Cú Pháp} \\
\codeinline{value} & Giá trị thực tế được đóng gói gửi về máy chủ khi tùy chọn đó được người dùng chọn. & Nếu không có, lấy văn bản hiển thị làm value & \codeinline{value="vn"} \\
\codeinline{selected} & Đặt tùy chọn này làm mặc định được chọn sẵn khi trang vừa tải xong. & Tùy chọn đầu tiên & \codeinline{selected} \\
\codeinline{disabled} & Làm mờ tùy chọn này, khiến người dùng không thể bấm chọn. & \codeinline{false} & \codeinline{disabled} \\
\codeinline{label} & Nhãn văn bản ngắn gọn hiển thị thay thế nội dung văn bản bên trong thẻ \codeinline{<option>}. & Nội dung bên trong thẻ & \codeinline{label="VN"} \\
\end{tabularx}
\end{prettyTableBox}

% ------------------------------------------------------------------------------
\subsection{Cơ Chế Hoạt Động Của name và value Khi Submit Form}

Sự kết hợp giữa thuộc tính \codeinline{name} trên thẻ \codeinline{<select>} và thuộc tính \codeinline{value} trên các thẻ \codeinline{<option>} quyết định chính xác dữ liệu nhận được trên máy chủ Backend:

\begin{conceptBox}[Cơ Chế Đóng Gói Dữ Liệu Form (Name \& Value)]
\begin{itemize}[leftmargin=1.5em, itemsep=3pt, parsep=0pt]
    \item \textbf{\codeinline{name}}: Định danh tên trường dữ liệu. Giúp máy chủ phân biệt dữ liệu gửi về thuộc về ô nhập nào (ví dụ: \codeinline{fruit}, \codeinline{country}, \codeinline{color}).
    \item \textbf{\codeinline{value}}: Giá trị dữ liệu thực tế được gửi đi khi một tùy chọn được chọn.
\end{itemize}
\end{conceptBox}

\begin{codeWindow}[Form chọn trái cây gửi dữ liệu bằng phương thức GET]
\begin{lstlisting}[language=HTML5]
<form action="submit.php" method="GET">
  <label for="fruit">Chon trai cay:</label>
  <select name="fruit" id="fruit">
    <option value="apple">(*@Táo@*)</option>
    <option value="banana">(*@Chuối@*)</option>
    <option value="mango">(*@Xoài@*)</option>
  </select>
  <button type="submit">(*@Gửi@*)</button>
</form>
\end{lstlisting}
\end{codeWindow}

\begin{tipBox}[Phân Tích URL Chuỗi Truy Vấn (Query String)]
\textbf{Kịch bản}: Người dùng chọn mục \textbf{"Chuối"} (\codeinline{value="banana"}) và nhấn nút \textbf{Gửi}.
\par\vspace{4pt}
Trình duyệt sẽ chuyển hướng tới địa chỉ URL:
\begin{center}
\tcbox[on line, arc=3pt, colback=primaryBlue!8, colframe=primaryBlue!40, boxsep=3pt, left=8pt, right=8pt, top=3pt, bottom=3pt]{%
  \small\ttfamily\bfseries submit.php?fruit=banana%
}
\end{center}
Trong đó: \codeinline{fruit} là giá trị thuộc tính \codeinline{name} của thẻ \codeinline{<select>}, và \codeinline{banana} là giá trị thuộc tính \codeinline{value} của thẻ \codeinline{<option>} được chọn.
\end{tipBox}

\paragraph{Trường Hợp Có value vs Không Có value:}
\begin{enumerate}[leftmargin=1.5em, itemsep=4pt, parsep=0pt]
    \item \textbf{Nếu thẻ có thuộc tính \codeinline{value}}:
    \begin{codeWindow}[Ví dụ tùy chọn khai báo value rõ ràng]
\begin{lstlisting}[language=HTML5]
<select name="color">
  <option value="red">(*@Đỏ@*)</option>
  <option value="blue">Xanh</option>
  <option value="green">(*@Lục@*)</option>
</select>
\end{lstlisting}
    \end{codeWindow}

    \begin{browserWindow}[Kết Quả Hiển Thị Trình Duyệt: Tùy Chọn Khai Báo Value Rõ Ràng]
    \sffamily\small
    \textbf{Chọn màu sắc:} \par\vspace{4pt}
    \tcbox[on line, arc=3pt, colback=white, colframe=primaryBlue!50, boxsep=3pt, left=8pt, right=20pt, top=3pt, bottom=3pt]{%
      \small \textcolor{black}{Xanh} \;\; \textcolor{primaryBlue}{\textbf{$\boldsymbol{\vee}$}}%
    } \par\vspace{4pt}
    \textcolor{codeComment}{\small\textit{$\rightarrow$ Khi người dùng chọn "Xanh", dữ liệu gửi về server là value="blue".}}
    \end{browserWindow}

    \item \textbf{Nếu thẻ KHÔNG KHAI BÁO thuộc tính \codeinline{value}}:
    \begin{codeWindow}[Ví dụ tùy chọn không có thuộc tính value]
\begin{lstlisting}[language=HTML5]
<select name="color">
  <option>(*@Đỏ@*)</option>
  <option>Xanh</option>
  <option>(*@Lục@*)</option>
</select>
\end{lstlisting}
    \end{codeWindow}

    \begin{browserWindow}[Kết Quả Hiển Thị Trình Duyệt: Tùy Chọn Không Khai Báo Value]
    \sffamily\small
    \textbf{Chọn màu sắc:} \par\vspace{4pt}
    \tcbox[on line, arc=3pt, colback=white, colframe=primaryBlue!50, boxsep=3pt, left=8pt, right=20pt, top=3pt, bottom=3pt]{%
      \small \textcolor{black}{Xanh} \;\; \textcolor{primaryBlue}{\textbf{$\boldsymbol{\vee}$}}%
    } \par\vspace{4pt}
    \textcolor{codeComment}{\small\textit{$\rightarrow$ Khi người dùng chọn "Xanh", trình duyệt tự lấy chữ "Xanh" làm value để gửi về server.}}
    \end{browserWindow}
\end{enumerate}

% ------------------------------------------------------------------------------
\subsection{Thiết Lập Tùy Chọn Mặc Định \& Tùy Chọn Trống (Placeholder Option)}

\paragraph{1. Tùy chọn mặc định tự động của trình duyệt:}
Theo mặc định, nếu không có bất kỳ thuộc tính \codeinline{selected} nào được chỉ định, trình duyệt sẽ tự động chọn \textbf{mục tùy chọn đầu tiên} trong danh sách:

\begin{codeWindow}[Mục đầu tiên (Táo) tự động được chọn mặc định]
\begin{lstlisting}[language=HTML5]
<select name="fruit">
  <option value="apple">(*@Táo@*)</option>  <!-- Muc nay se duoc chon mac dinh -->
  <option value="banana">(*@Chuối@*)</option>
  <option value="mango">(*@Xoài@*)</option>
</select>
\end{lstlisting}
\end{codeWindow}

\begin{browserWindow}[Kết Quả Hiển Thị Trình Duyệt: Mục Đầu Tiên Tự Động Chọn Mặc Định]
\sffamily\small
\textbf{Chọn trái cây:} \par\vspace{4pt}
\tcbox[on line, arc=3pt, colback=white, colframe=primaryBlue!50, boxsep=3pt, left=8pt, right=20pt, top=3pt, bottom=3pt]{%
  \small \textcolor{black}{Táo} \;\; \textcolor{primaryBlue}{\textbf{$\boldsymbol{\vee}$}}%
} \par\vspace{4pt}
\textcolor{codeComment}{\small\textit{$\rightarrow$ Mặc định khi nạp trang, mục đầu tiên ("Táo") tự động hiển thị.}}
\end{browserWindow}

\paragraph{2. Đặt một tùy chọn khác làm mặc định với thuộc tính \codeinline{selected}:}
Nếu muốn chọn một mục nằm ở vị trí khác làm mặc định khi tải trang, hãy thêm thuộc tính \codeinline{selected} vào thẻ \codeinline{<option>} đó:

\begin{codeWindow}[Đặt mục Chuối làm tùy chọn mặc định sẵn]
\begin{lstlisting}[language=HTML5]
<select name="fruit">
  <option value="apple">(*@Táo@*)</option>
  <option value="banana" selected>(*@Chuối@*)</option>  <!-- Duoc chon mac dinh -->
  <option value="mango">(*@Xoài@*)</option>
</select>
\end{lstlisting}
\end{codeWindow}

\begin{browserWindow}[Kết Quả Hiển Thị Trình Duyệt: Chọn Sẵn Mục Chuối Bằng Thuộc Tính Selected]
\sffamily\small
\textbf{Chọn trái cây:} \par\vspace{4pt}
\tcbox[on line, arc=3pt, colback=white, colframe=primaryBlue!50, boxsep=3pt, left=8pt, right=20pt, top=3pt, bottom=3pt]{%
  \small \textcolor{black}{Chuối} \;\; \textcolor{primaryBlue}{\textbf{$\boldsymbol{\vee}$}}%
} \par\vspace{4pt}
\textcolor{codeComment}{\small\textit{$\rightarrow$ Nhờ thuộc tính selected, mục "Chuối" hiển thị sẵn ngay khi trang vừa nạp xong.}}
\end{browserWindow}

\paragraph{3. Kỹ thuật tạo dòng lựa chọn hướng dẫn (Placeholder Option) ép người dùng phải chọn:}
Trong thực tế phát triển phần mềm, ta thường tạo một mục đầu tiên rỗng kết hợp với thuộc tính \codeinline{required} để buộc người dùng phải chủ động thao tác chọn một giá trị hợp lệ trước khi gửi form:

\begin{codeWindow}[Kỹ thuật Placeholder Option bắt buộc chọn]
\begin{lstlisting}[language=HTML5]
<select name="fruit" required>
  <option value="" selected disabled>-- Chon trai cay --</option>
  <option value="apple">(*@Táo@*)</option>
  <option value="banana">(*@Chuối@*)</option>
  <option value="mango">(*@Xoài@*)</option>
</select>
\end{lstlisting}
\end{codeWindow}

\begin{browserWindow}[Kết Quả Trình Duyệt: Tùy Chọn Hướng Dẫn Ban Đầu]
\sffamily\small
\textbf{Chọn trái cây:} \par\vspace{4pt}
\tcbox[on line, arc=3pt, colback=white, colframe=primaryBlue!50, boxsep=3pt, left=8pt, right=20pt, top=3pt, bottom=3pt]{%
  \small \textcolor{gray!70}{-- Chọn trái cây --} \;\; \textcolor{primaryBlue}{\textbf{$\boldsymbol{\vee}$}}%
} \par\vspace{4pt}
\textcolor{codeComment}{\small\textit{$\rightarrow$ Dòng "-- Chọn trái cây --" hiển thị ban đầu, bị mờ (disabled) và có value="" khiến form bắt buộc chọn mục khác.}}
\end{browserWindow}

% ------------------------------------------------------------------------------
\subsection{Các Ví Dụ Minh Họa Chi Tiết Thực Tế}

\paragraph{Ví Dụ 1: Select cho phép chọn nhiều giá trị cùng lúc (Multiple Select):}

\begin{codeWindow}[Mã HTML Select hỗ trợ chọn nhiều mục bằng thuộc tính multiple]
\begin{lstlisting}[language=HTML5]
<label for="pets">Chon thu cung yeu thich:</label>
<select id="pets" name="pets[]" multiple size="3">
  <option value="dog">(*@Chó@*)</option>
  <option value="cat">(*@Mèo@*)</option>
  <option value="bird">Chim</option>
  <option value="fish">(*@Cá@*)</option>
</select>
\end{lstlisting}
\end{codeWindow}

\begin{browserWindow}[Kết Quả Trình Duyệt: Hộp Chọn Nhiều Mục (Multiple Selection Listbox)]
\sffamily\small
\textbf{Chọn thú cưng yêu thích (Giữ Ctrl / Cmd để chọn nhiều):} \par\vspace{4pt}
\begin{tcolorbox}[arc=2pt, colback=white, colframe=primaryBlue!50, boxsep=3pt, width=0.55\linewidth]
  \small
  \tcbox[on line, arc=1pt, colback=primaryBlue, colframe=primaryBlue, boxsep=1.5pt, left=6pt, right=6pt]{\color{white}Chó} \par\vspace{2pt}
  \tcbox[on line, arc=1pt, colback=primaryBlue, colframe=primaryBlue, boxsep=1.5pt, left=6pt, right=6pt]{\color{white}Mèo} \par\vspace{2pt}
  \textcolor{black}{Chim} \par\vspace{2pt}
  \textcolor{black}{Cá}
\end{tcolorbox}
\par\vspace{4pt}
\textcolor{codeComment}{\small\textit{$\rightarrow$ Người dùng giữ phím Ctrl (Windows) hoặc Cmd (Mac) để bấm chọn cùng lúc Chó và Mèo.}}
\end{browserWindow}

% ------------------------------------------------------------------------------
\subsection{Chi Tiết Về Thẻ Phân Nhóm Tùy Chọn <optgroup>}

\begin{conceptBox}[Định Nghĩa Thẻ <optgroup>]
Thẻ \codeinline{<optgroup>} (\term{Option Group}) trong HTML được sử dụng để \textbf{gom nhóm các tùy chọn} (\codeinline{<option>}) có cùng chủ đề lại với nhau bên trong một thẻ \codeinline{<select>}.
\par\vspace{4pt}
\textbf{ Mục Đích}: Giúp danh sách thả xuống có cấu trúc phân tầng rõ ràng, mạch lạc, dễ đọc và dễ tìm kiếm, đặc biệt khi danh sách chứa rất nhiều mục lựa chọn.
\end{conceptBox}

\paragraph{1. Cú pháp cơ bản của thẻ <optgroup>:}

\begin{codeWindow}[Mã HTML phân nhóm danh sách thực đơn với optgroup]
\begin{lstlisting}[language=HTML5]
<select name="food">
  <optgroup label="Trai cay">
    <option value="apple">(*@Táo@*)</option>
    <option value="banana">(*@Chuối@*)</option>
  </optgroup>
  <optgroup label="Do uong">
    <option value="coffee">(*@Cà *@phê@*)</option>
    <option value="tea">(*@Trà@*)</option>
  </optgroup>
</select>
\end{lstlisting}
\end{codeWindow}

\begin{browserWindow}[Kết Quả Hiển Thị Trình Duyệt: Dropdown Hiển Thị 2 Nhóm (Trái cây \& Đồ uống)]
\sffamily\small
\textbf{Chọn thực đơn:} \par\vspace{4pt}
\begin{tcolorbox}[arc=2pt, colback=white, colframe=primaryBlue!50, boxsep=4pt, width=0.6\linewidth]
  \small
  \textbf{\textcolor{gray!80}{TRÁI CÂY}} \par\vspace{2pt}
  \quad Táo \par\vspace{2pt}
  \quad Chuối \par\vspace{3pt}
  \textbf{\textcolor{gray!80}{ĐỒ UỐNG}} \par\vspace{2pt}
  \quad Cà phê \par\vspace{2pt}
  \quad Trà
\end{tcolorbox}
\par\vspace{4pt}
\textcolor{codeComment}{\small\textit{$\rightarrow$ Dropdown hiển thị rõ ràng 2 nhóm: Trái cây và Đồ uống với các tiêu đề nhóm được in đậm.}}
\end{browserWindow}

\paragraph{2. Bảng thuộc tính quan trọng của thẻ <optgroup>:}

\begin{prettyTableBox}[Bảng Thuộc Tính Của Thẻ Optgroup]
\rowcolors{2}{white}{primaryBlue!4}
\begin{tabularx}{\linewidth}{K{2.2cm} Z K{2.5cm} Z}
\rowcolor{primaryBlue}
\thd{Thuộc Tính} & \thd{Ý Nghĩa \& Chức Năng} & \thd{Giá Trị Mặc Định} & \thd{Ví Dụ Cú Pháp} \\
\codeinline{label} & Tiêu đề văn bản của nhóm (\textbf{bắt buộc phải có} để hiển thị tên nhóm). & Không có & \codeinline{label="Xe Đức"} \\
\codeinline{disabled} & Vô hiệu hóa toàn bộ nhóm tùy chọn (người dùng không thể chọn các mục trong nhóm này). & \codeinline{false} & \codeinline{disabled} \\
\end{tabularx}
\end{prettyTableBox}

\paragraph{3. Các ví dụ minh họa chi tiết:}

\textbf{Ví dụ 1: Dropdown nhóm tùy chọn xe thương hiệu Đức và Nhật}:
\begin{codeWindow}[Mã HTML phân nhóm danh sách hãng xe bằng optgroup]
\begin{lstlisting}[language=HTML5]
<select name="car">
  <optgroup label="Xe Duc">
    <option value="bmw">BMW</option>
    <option value="audi">Audi</option>
  </optgroup>
  <optgroup label="Xe Nhat">
    <option value="toyota">Toyota</option>
    <option value="honda">Honda</option>
  </optgroup>
</select>
\end{lstlisting}
\end{codeWindow}

\begin{browserWindow}[Kết Quả Hiển Thị Trình Duyệt: Danh Sách Nhóm Xe Đức \& Xe Nhật]
\sffamily\small
\textbf{Chọn hãng xe:} \par\vspace{4pt}
\begin{tcolorbox}[arc=2pt, colback=white, colframe=primaryBlue!50, boxsep=4pt, width=0.6\linewidth]
  \small
  \textbf{\textcolor{gray!80}{XE ĐỨC}} \par\vspace{2pt}
  \quad BMW \par\vspace{2pt}
  \quad Audi \par\vspace{3pt}
  \textbf{\textcolor{gray!80}{XE NHẬT}} \par\vspace{2pt}
  \quad Toyota \par\vspace{2pt}
  \quad Honda
\end{tcolorbox}
\end{browserWindow}

\textbf{Ví dụ 2: Nhóm bị vô hiệu hóa bằng thuộc tính disabled}:
\begin{codeWindow}[Mã HTML nhóm sản phẩm Hết hàng bị vô hiệu hóa bằng disabled]
\begin{lstlisting}[language=HTML5]
<select name="drink">
  <optgroup label="Con hang">
    <option value="coffee">(*@Cà *@phê@*)</option>
    <option value="tea">(*@Trà@*)</option>
  </optgroup>
  <optgroup label="Het hang" disabled>
    <option value="beer">Bia</option>
    <option value="wine">(*@Rượu@*)</option>
  </optgroup>
</select>
\end{lstlisting}
\end{codeWindow}

\begin{browserWindow}[Kết Quả Hiển Thị Trình Duyệt: Nhóm Hết Hàng Bị Vô Hiệu Hóa]
\sffamily\small
\textbf{Chọn đồ uống tại quán:} \par\vspace{4pt}
\begin{tcolorbox}[arc=2pt, colback=white, colframe=primaryBlue!50, boxsep=4pt, width=0.6\linewidth]
  \small
  \textbf{\textcolor{gray!80}{CÒN HÀNG}} \par\vspace{2pt}
  \quad Cà phê \par\vspace{2pt}
  \quad Trà \par\vspace{3pt}
  \textbf{\textcolor{gray!40}{HẾT HÀNG (Vô hiệu hóa)}} \par\vspace{2pt}
  \quad \textcolor{gray!40}{Bia} \par\vspace{2pt}
  \quad \textcolor{gray!40}{Rượu}
\end{tcolorbox}
\par\vspace{4pt}
\textcolor{codeComment}{\small\textit{$\rightarrow$ Nhóm "Hết hàng" (bao gồm Bia và Rượu) bị mờ đi và người dùng hoàn toàn không thể bấm chọn.}}
\end{browserWindow}

\begin{warnBox}[4 Lưu Ý Cốt Lõi Khi Sử Dụng Thẻ <optgroup>]
\begin{enumerate}[leftmargin=1.5em, itemsep=4pt, parsep=0pt]
    \item \textbf{Phải nằm bên trong thẻ \codeinline{<select>}}: Thẻ \codeinline{<optgroup>} chỉ hoạt động khi nằm trực tiếp bên trong thẻ \codeinline{<select>}, không thể đứng độc lập ngoài form.
    \item \textbf{Thuộc tính \codeinline{label} là bắt buộc}: Nếu không khai báo thuộc tính \codeinline{label}, nhóm sẽ không hiển thị tên tiêu đề trên giao diện.
    \item \textbf{Khả năng kết hợp với \codeinline{multiple}}: Có thể kết hợp thuộc tính \codeinline{multiple} trên thẻ \codeinline{<select>} để cho phép người dùng chọn cùng lúc nhiều tùy chọn nằm ở các nhóm khác nhau.
    \item \textbf{CSS tùy biến hạn chế}: Các thuộc tính CSS dành cho \codeinline{<optgroup>} bị trình duyệt giới hạn rất nhiều (không thể thay đổi kiểu phông, màu nền phức tạp trên một số trình duyệt mặc định).
\end{enumerate}
\end{warnBox}

% ------------------------------------------------------------------------------
\subsection{Các Lưu Ý Quan Trọng Cần Ghi Nhớ Nằm Lòng}

\begin{warnBox}[Các Lưu Ý Cốt Lõi Khi Sử Dụng Thẻ Select \& Option]
\begin{enumerate}[leftmargin=1.5em, itemsep=6pt, parsep=0pt]
    \item \textbf{Tên trường dữ liệu cho Multiple Select (\codeinline{name="pets[]"})}:
    Khi sử dụng thuộc tính \codeinline{multiple}, để các ngôn ngữ xử lý Backend (như PHP, Node.js) nhận đầy đủ tất cả các giá trị được chọn dưới dạng \textbf{một mảng (Array)}, bạn nên đặt dấu ngoặc vuông vuông sau tên thuộc tính \codeinline{name}: \codeinline{name="pets[]"}.

    \item \textbf{Giá trị mặc định của thuộc tính \codeinline{value}}:
    Nếu bỏ quên không ghi thuộc tính \codeinline{value} trong thẻ \codeinline{<option>}, trình duyệt sẽ tự động lấy toàn bộ nội dung văn bản hiển thị làm giá trị gửi đi. Hãy luôn chủ động khai báo \codeinline{value} bằng chữ thường không dấu để dễ xử lý dữ liệu Backend.

    \item \textbf{Giới hạn khả năng tùy biến CSS của thẻ <select> mặc định}:
    Giao diện thẻ \codeinline{<select>} mặc định do hệ điều hành và trình duyệt render nên rất khó để tùy biến CSS phức tạp (như thêm icon, hiệu ứng hover đẹp). Trong các dự án thực tế đòi hỏi UX/UI cao cấp, lập trình viên thường dùng các thư viện UI JavaScript (như \term{Select2}, \term{Bootstrap Select}, \term{React-Select}).
\end{enumerate}
\end{warnBox}

% ------------------------------------------------------------------------------
\subsection{Đọc Thêm: Kỹ Thuật Nâng Cao \& Tối Ưu Lập Trình Dropdown}

\begin{tipBox}[Kỹ Thuật Tạo Dropdown Động Tự Động Thay Đổi Theo Danh Sách Cha (Cascading Select)]
Trong các form đăng ký địa chỉ (Tỉnh/Thành phố $\rightarrow$ Quận/Huyện $\rightarrow$ Phường/Xã), khi người dùng chọn Tỉnh/Thành phố ở thẻ `<select>` thứ nhất, JavaScript sẽ lắng nghe sự kiện \codeinline{change} để nạp danh sách Quận/Huyện tương ứng vào thẻ `<select>` thứ hai:
\begin{codeWindow}[Mã JavaScript tự động cập nhật danh sách Select con]
\begin{lstlisting}[language=JavaScript]
const provinceSelect = document.getElementById('province');
const districtSelect = document.getElementById('district');

provinceSelect.addEventListener('change', function() {
  const selectedProvince = this.value;
  // (*@Gọi API hoặc lấy dữ liệu Quận/Huyện dựa theo Tỉnh/Thành đã chọn@*)
  updateDistrictOptions(selectedProvince);
});
\end{lstlisting}
\end{codeWindow}
\end{tipBox}

\begin{noteBox}[Tối Ưu Truy Cập (a11y) Cho Thẻ Select]
\begin{itemize}[leftmargin=1.5em, itemsep=2pt, parsep=0pt]
    \item Luôn luôn liên kết thẻ \codeinline{<select>} với thẻ \codeinline{<label>} bằng cặp thuộc tính \codeinline{for} - \codeinline{id}.
    \item Khi dùng thuộc tính \codeinline{multiple}, nên bổ sung văn bản hướng dẫn rõ ràng (ví dụ: \textit{"Giữ phím Ctrl hoặc Cmd để chọn nhiều mục"}) giúp người dùng thiết bị di động hoặc trình đọc màn hình dễ dàng thao tác.
\end{itemize}
\end{noteBox}

\begin{prettyTableBox}[Bảng So Sánh Chi Tiết: <select> vs <datalist>]
\small
\rowcolors{2}{white}{primaryBlue!4}
\begin{tabularx}{\linewidth}{K{2.8cm} Z Z}
\rowcolor{primaryBlue}
\thd{Tiêu Chí So Sánh} & \thd{Thẻ Danh Sách Thả Xuống <select>} & \thd{Thẻ Gợi Ý Tự Động <datalist>} \\
\textbf{Giới hạn lựa chọn} & Ép buộc người dùng \textbf{chỉ được chọn} các mục trong danh sách có sẵn. & Cho phép vừa chọn mục gợi ý vừa \textbf{nhập văn bản tự do}. \\
\textbf{Trải nghiệm tìm kiếm} & Mặc định không hỗ trợ ô gõ từ khóa tìm kiếm lọc dữ liệu. & Tự động lọc danh sách gợi ý theo ký tự người dùng vừa gõ. \\
\textbf{Đóng gói dữ liệu} & Gửi giá trị \codeinline{value} của thẻ \codeinline{<option>} được chọn. & Gửi chuỗi ký tự hiển thị trong ô \codeinline{<input>}. \\
\textbf{Trường hợp sử dụng} & Chọn Tỉnh/Thành, Quốc gia, Giới tính, Trạng thái đơn hàng. & Ô tìm kiếm từ khóa, Chọn thương hiệu, Ô chọn có từ gợi ý. \\
\end{tabularx}
\end{prettyTableBox}

\begin{tipBox}[Mẹo Phỏng Vấn Tuyển Dụng: Xử Lý Select Đa Giá Trị Trong JavaScript]
\textbf{Câu hỏi phỏng vấn}: Làm thế nào để trích xuất mảng tất cả các giá trị được chọn từ thẻ \codeinline{<select multiple>} bằng JavaScript?
\par\vspace{4pt}
\textbf{Cú pháp tối ưu}:
\begin{codeWindow}[Trích xuất mảng giá trị từ Select Multiple trong JS]
\begin{lstlisting}[language=JavaScript]
const selectElem = document.getElementById('pets');
// (*@Chuyển danh sách HTMLOptionsCollection thành mảng Array và lọc các mục selected@*)
const selectedValues = Array.from(selectElem.selectedOptions).map(opt => opt.value);
console.log(selectedValues); // (*@Kết quả mảng: ['dog', 'cat']@*)
\end{lstlisting}
\end{codeWindow}
\end{tipBox}

\begin{warnBox}[Lưu Ý Bảo Mật: Kiểm Tra Dữ Liệu Select Phía Server (Backend Validation)]
Tuyệt đối \textbf{không tin tưởng dữ liệu từ thẻ \codeinline{<select>} gửi về}! Dù trên giao diện Frontend bạn chỉ cho chọn 3 tùy chọn \codeinline{["vn", "us", "jp"]}, người dùng xấu hoàn toàn có thể mở Developer Tools (F12) sửa giá trị \codeinline{value="hacker"} trước khi bấm Gửi. Máy chủ Backend luôn luôn phải xác thực giá trị nhận được thuộc Danh sách Trắng (\term{Whitelist Validation}).
\end{warnBox}

% ------------------------------------------------------------------------------
% ==============================================================================
% MỤC: THAY ĐỔI FONT CHỮ MẶC ĐỊNH TRONG FORM HTML
% ==============================================================================
\section{Thay Đổi Font Chữ Mặc Định Trong Form HTML}
\label{sec:form_font_customization}

\begin{conceptBox}[Vấn Đề Kế Thừa Font Trong Form HTML]
Trong HTML, các phần tử điều khiển biểu mẫu như \codeinline{<input>}, \codeinline{<textarea>}, \codeinline{<select>} và \codeinline{<button>} \textbf{không tự động kế thừa \codeinline{font-family} từ thẻ \codeinline{<body>}}.
\par\vspace{4pt}
Nếu bạn chỉ khai báo font chữ cho thẻ \codeinline{<body>}, các ô nhập liệu trong form vẫn sẽ hiển thị phông chữ mặc định của trình duyệt/hệ điều hành (như Times New Roman hoặc Segoe UI), làm mất tính đồng bộ thiết kế giao diện!
\end{conceptBox}

\subsection{1. Phương Pháp Thay Đổi Font Đồng Bộ Cực Kỳ Đơn Giản}
Để ép tất cả các phần tử trong form dùng chung một phông chữ chuẩn với trang web, bạn có thể áp dụng thuộc tính CSS \codeinline{font-family: inherit;}:

\begin{codeWindow}[Mã CSS đồng bộ phông chữ toàn bộ form bằng font-family: inherit]
\begin{lstlisting}[language=HTML5]
/* Khai bao font chinh cho toan trang */
body {
  font-family: "Roboto", "Segoe UI", sans-serif;
}

/* Ep tat ca phan tu form ke thua font tu body */
form, input, textarea, select, button, label {
  font-family: inherit;
  font-size: 16px;
  color: #333;
}
\end{lstlisting}
\end{codeWindow}

\subsection{2. Bảng Các Thuộc Tính Liên Quan Đến Font Chữ Trong Form}

\begin{prettyTableBox}[Bảng Thuộc Tính Định Kiểu Font Chữ Trong Form]
\rowcolors{2}{white}{primaryBlue!4}
\begin{tabularx}{\linewidth}{K{2.6cm} Z K{3.2cm}}
\rowcolor{primaryBlue}
\thd{Thuộc Tính CSS} & \thd{Ý Nghĩa \& Tác Dụng} & \thd{Ví Dụ Cú Pháp} \\
\codeinline{font-family} & Chọn họ phông chữ hiển thị. & \codeinline{font-family: Arial, sans-serif;} \\
\codeinline{font-size} & Điều chỉnh kích thước văn bản. & \codeinline{font-size: 16px;} \\
\codeinline{font-weight} & Độ đậm của chữ (100-900 hoặc \codeinline{normal}, \codeinline{bold}). & \codeinline{font-weight: bold;} \\
\codeinline{font-style} & Kiểu dáng chữ (\codeinline{normal}, \codeinline{italic}, \codeinline{oblique}). & \codeinline{font-style: italic;} \\
\codeinline{line-height} & Chiều cao dòng (giúp văn bản thoáng và dễ đọc hơn). & \codeinline{line-height: 1.5;} \\
\codeinline{color} & Màu sắc văn bản nội dung. & \codeinline{color: \#333333;} \\
\end{tabularx}
\end{prettyTableBox}

\subsection{3. Ví Dụ Mã Nguồn HTML \& CSS Đầy Đủ}

\begin{codeWindow}[Mã HTML \& CSS hoàn chỉnh đổi phông chữ form với Google Fonts]
\begin{lstlisting}[language=HTML5]
<!DOCTYPE html>
<html lang="vi">
<head>
  <meta charset="UTF-8">
  <title>Form Doi Font</title>
  <!-- Nhap Google Font Roboto -->
  <link href="https://fonts.googleapis.com/css2?family=Roboto&display=swap" rel="stylesheet">
  <style>
    body {
      font-family: "Roboto", sans-serif;
    }
    form, input, textarea, select, button, label {
      font-family: inherit; /* Ke thua tu body */
      font-size: 16px;
      color: #333;
    }
  </style>
</head>
<body>
  <h2>Form Vi Du</h2>
  <form>
    <label for="name">(*@Họ tên:@*)</label><br>
    <input type="text" id="name" name="name"><br><br>
    
    <label for="email">Email:</label><br>
    <input type="email" id="email" name="email"><br><br>
    
    <label for="msg">Tin nhan:</label><br>
    <textarea id="msg" name="msg"></textarea><br><br>
    
    <button type="submit">(*@Gửi@*)</button>
  </form>
</body>
</html>
\end{lstlisting}
\end{codeWindow}

\begin{browserWindow}[Kết Quả Hiển Thị Trình Duyệt: Form Đồng Bộ Font Chữ Roboto]
\sffamily\small
{\bfseries\large Form Ví Dụ (Đồng Bộ Font Roboto)}\\[6pt]
\textbf{Họ tên:}\\[2pt]
\tcbox[on line, arc=3pt, colback=white, colframe=gray!40, left=6pt, right=50pt, top=3pt, bottom=3pt]{\small Nguyễn Trọng Tấn}\\[6pt]
\textbf{Email:}\\[2pt]
\tcbox[on line, arc=3pt, colback=white, colframe=gray!40, left=6pt, right=50pt, top=3pt, bottom=3pt]{\small contact@domain.com}\\[6pt]
\textbf{Tin nhắn:}\\[2pt]
\tcbox[on line, arc=3pt, colback=white, colframe=gray!40, left=6pt, right=60pt, top=3pt, bottom=20pt]{\small Nội dung tin nhắn...}\\[8pt]
\tcbox[on line, arc=3pt, colback=primaryBlue, colframe=primaryBlue, left=12pt, right=12pt, top=3pt, bottom=3pt]{\color{white}\bfseries\small Gửi}
\end{browserWindow}

\begin{warnBox}[4 Lưu Ý Cốt Lõi Khi Thay Đổi Font Chữ Trong Form]
\begin{enumerate}[leftmargin=1.5em, itemsep=4pt, parsep=0pt]
    \item \textbf{Dùng \codeinline{font-family: inherit;}}: Giúp các phần tử form tự động kế thừa font từ \codeinline{<body>}, giúp dễ quản lý trung tâm (chỉ cần đổi font 1 nơi ở body là toàn bộ form đổi theo).
    \item \textbf{Tránh chỉ đặt font cho \codeinline{<body>}}: Nếu quên bộ chọn \codeinline{input, select, textarea, button}, trình duyệt sẽ rơi về phông mặc định của hệ điều hành.
    \item \textbf{Nạp font trước khi dùng}: Nếu dùng Google Fonts hoặc Custom Font, phải nạp thẻ \codeinline{<link>} hoặc \codeinline{@import} trước khi áp dụng CSS.
    \item \textbf{Ghi đè cho nút bấm}: Nút bấm (\codeinline{<button>}, \codeinline{input[type="submit"]}) luôn sở hữu kiểu dáng mặc định riêng của trình duyệt, do đó luôn phải ghi đè \codeinline{font-family: inherit;}.
\end{enumerate}
\end{warnBox}

% ==============================================================================
% MỤC: GIẢ LỚP CSS :placeholder-shown & KỸ THUẬT FLOATING LABEL
% ==============================================================================
\section{Giả Lớp CSS :placeholder-shown \& Kỹ Thuật Floating Label}
\label{sec:css_placeholder_shown}

\begin{conceptBox}[Giả Lớp :placeholder-shown Là Gì?]
\codeinline{:placeholder-shown} là một CSS Pseudo-class (Giả lớp) được trình duyệt tự động kích hoạt cho các phần tử \codeinline{<input>} hoặc \codeinline{<textarea>} \textbf{khi phần tử đó đang hiển thị văn bản hướng dẫn placeholder} (tức là ô input đang rỗng).
\par\vspace{4pt}
\begin{itemize}[leftmargin=1.5em, itemsep=2pt, parsep=0pt]
    \item \textbf{Input trống} $\rightarrow$ Placeholder hiển thị $\rightarrow$ \codeinline{:placeholder-shown} được kích hoạt.
    \item \textbf{Input có dữ liệu} $\rightarrow$ Placeholder biến mất $\rightarrow$ \codeinline{:placeholder-shown} bị vồ hiệu hóa.
\end{itemize}
\end{conceptBox}

\subsection{1. Cú Pháp Cơ Bản \& Ví Dụ Minh Họa}

\begin{codeWindow}[Mã CSS thay đổi viền và nền ô Input khi placeholder đang hiển thị]
\begin{lstlisting}[language=HTML5]
<!DOCTYPE html>
<html lang="vi">
<head>
  <meta charset="UTF-8">
  <title>Vi du placeholder-shown</title>
  <style>
    input {
      padding: 8px;
      border: 2px solid #ccc;
      border-radius: 4px;
    }
    /* Khi placeholder con hien thi (o input trong) */
    input:placeholder-shown {
      border-color: orange;
      background-color: #fff8e1;
    }
    /* Khi nguoi dung da nhap text vao */
    input:not(:placeholder-shown) {
      border-color: green;
      background-color: #f0fff4;
    }
  </style>
</head>
<body>
  <form>
    <label for="name">(*@Tên:@*)</label><br>
    <input type="text" id="name" placeholder="(*@Nhập *@tên của bạn@*)">
  </form>
</body>
</html>
\end{lstlisting}
\end{codeWindow}

\begin{browserWindow}[Kết Quả Hiển Thị Trình Duyệt: Trạng Thái Ban Đầu vs Khi Nhập Dữ Liệu]
\sffamily\small
\textbf{1. Trạng thái ban đầu (Input trống - Placeholder hiện):}\\[2pt]
\tcbox[on line, arc=3pt, colback=yellow!10, colframe=accentOrange, boxsep=2pt, left=8pt, right=40pt, top=3pt, bottom=3pt]{\color{gray!70}\small Nhập tên của bạn} \quad \textcolor{accentOrange}{\bfseries (Viền cam, Nền vàng nhạt)}\\[8pt]
\textbf{2. Khi người dùng nhập dữ liệu (Placeholder biến mất):}\\[2pt]
\tcbox[on line, arc=3pt, colback=green!10, colframe=green!60!black, boxsep=2pt, left=8pt, right=40pt, top=3pt, bottom=3pt]{\color{black}\small Nguyễn Trọng Tấn} \quad \textcolor{green!60!black}{\bfseries (Viền xanh lá, Nền xanh nhạt)}
\end{browserWindow}

\begin{warnBox}[4 Lưu Ý Cốt Lõi Khi Sử Dụng :placeholder-shown]
\begin{enumerate}[leftmargin=1.5em, itemsep=4pt, parsep=0pt]
    \item \textbf{Chỉ dùng cho trường có placeholder}: Giả lớp này chỉ hoạt động với các thẻ \codeinline{<input>} và \codeinline{<textarea>} có khai báo thuộc tính \codeinline{placeholder="..."}.
    \item \textbf{Không áp dụng cho thẻ <select>}: Thẻ \codeinline{<select>} không có thuộc tính placeholder nên không thể dùng \codeinline{:placeholder-shown}.
    \item \textbf{Phân biệt với ::placeholder}: Selector \codeinline{::placeholder} dùng để định kiểu văn bản chữ gợi ý bên trong, còn \codeinline{:placeholder-shown} dùng để định kiểu toàn bộ ô input dựa theo trạng thái dữ liệu.
    \item \textbf{Tối ưu UX Form Validation}: Giúp đánh dấu trực quan ô chưa nhập dữ liệu mà không cần viết mã JavaScript.
\end{enumerate}
\end{warnBox}

% ==============================================================================
% MỤC: GIẢ LỚP CSS :not(:placeholder-shown) & ỨNG DỤNG THỰC TẾ
% ==============================================================================
\section{Giả Lớp CSS :not(:placeholder-shown) \& Ứng Dụng Floating Label}
\label{sec:css_not_placeholder_shown}

\begin{conceptBox}[Ý Nghĩa Của :not(:placeholder-shown)]
Đây là sự kết hợp giữa hai giả lớp \codeinline{:not()} và \codeinline{:placeholder-shown}.
\par\vspace{4pt}
\textbf{ Ý Nghĩa}: Chọn các phần tử \codeinline{<input>} hoặc \codeinline{<textarea>} \textbf{không còn hiển thị placeholder} (tức là khi người dùng đã chủ động nhập dữ liệu vào ô).
\end{conceptBox}

\subsection{1. Mã Nguồn Form Validation Thuần CSS với :not(:placeholder-shown)}

\begin{codeWindow}[Mã HTML \& CSS đổi màu viền khi ô Input có dữ liệu]
\begin{lstlisting}[language=HTML5]
<!DOCTYPE html>
<html lang="vi">
<head>
  <meta charset="UTF-8">
  <title>Vi du :not(:placeholder-shown)</title>
  <style>
    input {
      padding: 8px;
      border: 2px solid #ccc;
      border-radius: 4px;
    }
    /* Khi field con trong (*@viền đỏ@*) */
    input:placeholder-shown {
      border-color: red;
      background-color: #fff5f5;
    }
    /* Khi field da nhap du lieu (*@viền xanh lá@*) */
    input:not(:placeholder-shown) {
      border-color: green;
      background-color: #f0fff4;
    }
  </style>
</head>
<body>
  <form>
    <label for="email">Email:</label><br>
    <input type="email" id="email" placeholder="(*@Nhập email...@*)">
  </form>
</body>
</html>
\end{lstlisting}
\end{codeWindow}

\begin{browserWindow}[Kết Quả Hiển Thị Trình Duyệt: Validation Thuần CSS]
\sffamily\small
\textbf{1. Khi chưa nhập (Input trống):}\\[2pt]
\tcbox[on line, arc=3pt, colback=red!10, colframe=red!70!black, boxsep=2pt, left=8pt, right=40pt, top=3pt, bottom=3pt]{\color{gray!70}\small Nhập email...} \quad \textcolor{red!70!black}{\bfseries (Viền đỏ - Chưa nhập)}\\[8pt]
\textbf{2. Khi đã nhập text (Dữ liệu sẵn sàng):}\\[2pt]
\tcbox[on line, arc=3pt, colback=green!10, colframe=green!60!black, boxsep=2pt, left=8pt, right=40pt, top=3pt, bottom=3pt]{\color{black}\small tan@domain.com} \quad \textcolor{green!60!black}{\bfseries (Viền xanh lá - Đã nhập)}
\end{browserWindow}

\subsection{2. Các Ứng Dụng Thực Tế Đỉnh Cao}

\begin{tipBox}[3 Ứng Dụng Thực Tế Đỉnh Cao Của :not(:placeholder-shown)]
\begin{enumerate}[leftmargin=1.5em, itemsep=4pt, parsep=0pt]
    \item \textbf{Form Validation Nhẹ Bằng CSS}: Đánh dấu trạng thái ô nhập dữ liệu ngay khi người dùng hoàn thành gõ chữ mà không cần lắng nghe sự kiện JavaScript.
    \item \textbf{Kỹ Thuật Floating Label (Nhãn Trượt Nổi Thuần CSS)}: Khi người dùng bấm vào ô và bắt đầu nhập text, thẻ \codeinline{<label>} tự động thu nhỏ và bay lên mép trên ô input:
    \begin{codeWindow}[Mã CSS tạo hiệu ứng Floating Label không dùng JS]
\begin{lstlisting}[language=HTML5]
/* Khi input da co du lieu -> nhan label bay len tren */
input:not(:placeholder-shown) + label {
  transform: translateY(-20px);
  font-size: 12px;
  color: #1a73e8;
}
\end{lstlisting}
    \end{codeWindow}
    \item \textbf{Tối Ưu Hiệu Năng}: Loại bỏ mã JavaScript lắng nghe sự kiện \codeinline{input} hay \codeinline{change}, giúp trình duyệt render giao diện mượt mà và tiết kiệm RAM.
\end{enumerate}
\end{tipBox}

\begin{warnBox}[Lưu Ý Quan Trọng Với Thuộc Tính Placeholder]
\begin{itemize}[leftmargin=1.5em, itemsep=3pt, parsep=0pt]
    \item \textbf{Yêu cầu thuộc tính placeholder}: Thuộc tính \codeinline{:placeholder-shown} bắt buộc ô input phải có thuộc tính \codeinline{placeholder="..."}. Nếu bạn quên khai báo thuộc tính này, \codeinline{:placeholder-shown} sẽ không bao giờ khớp, khiến giả lớp \codeinline{:not(:placeholder-shown)} luôn luôn đúng (luôn có tác dụng ngay cả khi input trống!).
    \item \textbf{Mẹo nhỏ}: Nếu muốn dùng kỹ thuật Floating Label nhưng không muốn hiện chữ gợi ý, hãy khai báo thuộc tính rỗng: \codeinline{placeholder=" "}.
\end{itemize}
\end{warnBox}

\subsection{3. Đọc Thêm: Phân Biệt Phân Cấp Pseudo-class :placeholder-shown vs Pseudo-element ::placeholder}

\begin{prettyTableBox}[Bảng So Sánh Chi Tiết: :placeholder-shown vs ::placeholder]
\small
\rowcolors{2}{white}{primaryBlue!4}
\begin{tabularx}{\linewidth}{K{2.8cm} Z Z}
\rowcolor{primaryBlue}
\thd{Tiêu Chí So Sánh} & \thd{Giả Lớp :placeholder-shown (Pseudo-class)} & \thd{Phần Tử Giả ::placeholder (Pseudo-element)} \\
\textbf{Đối tượng tác động} & Định kiểu cho \textbf{toàn bộ ô input / textarea} (viền, nền, bóng đổ, vị trí label con). & Chỉ định kiểu cho \textbf{chữ văn bản gợi ý} nằm bên trong ô. \\
\textbf{Cú pháp CSS} & Sử dụng dấu hai chấm đơn \codeinline{:placeholder-shown}. & Sử dụng dấu hai chấm kép \codeinline{::placeholder}. \\
\textbf{Các thuộc tính CSS hỗ trợ} & Hỗ trợ toàn bộ thuộc tính CSS (\codeinline{border}, \codeinline{background}, \codeinline{transform}, \codeinline{padding}). & Chỉ hỗ trợ các thuộc tính văn bản (\codeinline{color}, \codeinline{font-style}, \codeinline{opacity}). \\
\textbf{Mục đích sử dụng} & Kiểm tra trạng thái dữ liệu đã nhập hay chưa để tạo hiệu ứng UI/UX. & Đổi màu chữ gợi ý mờ hơn hoặc nghiêng để tránh nhầm với chữ user gõ. \\
\end{tabularx}
\end{prettyTableBox}

\begin{tipBox}[Kỹ Thuật Kết Hợp Định Kiểu Hoàn Hảo Cả 2 Giả Lớp]
\begin{codeWindow}[Mã CSS kết hợp định kiểu cả văn bản gợi ý lẫn trạng thái ô input]
\begin{lstlisting}[language=HTML5]
/* 1. Dinh kieu chu goi y nghieng va mo 50% */
input::placeholder {
  color: #888888;
  font-style: italic;
  opacity: 0.7;
}

/* 2. Dinh kieu o input khi dang cho nhap du lieu */
input:placeholder-shown {
  border: 1.5px solid #1a73e8;
  box-shadow: 0 0 5px rgba(26, 115, 232, 0.2);
}
\end{lstlisting}
\end{codeWindow}
\end{tipBox}





% ==============================================================================
% MỤC: GIẢ LỚP CSS :focus-within & HIỆU ỨNG FOCUS NHÓM
% ==============================================================================
\section{Giả Lớp CSS :focus-within \& Hiệu Ứng Focus Nhóm}
\label{sec:css_focus_within}

\begin{conceptBox}[Giả Lớp :focus-within Trong CSS Là Gì?]
\codeinline{:focus-within} là một CSS Pseudo-class (Giả lớp). Nó chọn một phần tử khi **chính bản thân nó** hoặc **bất kỳ phần tử con nào bên trong nó** đang được nhận tiêu điểm focus.
\par\vspace{4pt}
\begin{itemize}[leftmargin=1.5em, itemsep=2pt, parsep=0pt]
    \item \textbf{Cơ chế hoạt động}: Nếu người dùng nhấp chuột hoặc dùng phím Tab focus vào một ô \codeinline{<input>} nằm bên trong một khối \codeinline{<form>} hoặc \codeinline{<div class="form-group">}, thì toàn bộ khối cha đó cũng sẽ lập tức khớp với giả lớp \codeinline{:focus-within}.
\end{itemize}
\end{conceptBox}

\subsection{1. Cú Pháp Cơ Bản}
\begin{codeWindow}[Cú pháp CSS đổi viền khối container khi có con được focus]
\begin{lstlisting}[language=HTML5]
.container:focus-within {
  border: 2px solid blue;
}
\end{lstlisting}
\end{codeWindow}

\subsection{2. Ví Dụ Minh Họa Chi Tiết}

\begin{codeWindow}[Mã HTML \& CSS hoàn chỉnh ứng dụng :focus-within làm nổi bật khối Form Group]
\begin{lstlisting}[language=HTML5]
<!DOCTYPE html>
<html lang="vi">
<head>
  <meta charset="UTF-8">
  <title>Vi du :focus-within</title>
  <style>
    .form-group {
      border: 2px solid #ccc;
      padding: 10px;
      border-radius: 6px;
      transition: 0.3s;
    }
    /* Khi co input ben trong duoc focus */
    .form-group:focus-within {
      border-color: blue;
      background-color: #f0f8ff;
    }
    input {
      padding: 6px;
      border: 1px solid #aaa;
    }
  </style>
</head>
<body>
  <div class="form-group">
    <label for="name">(*@Họ tên:@*)</label><br>
    <input type="text" id="name" placeholder="(*@Nhập tên...@*)">
  </div>
</body>
</html>
\end{lstlisting}
\end{codeWindow}

\begin{browserWindow}[Kết Quả Hiển Thị Trình Duyệt: Trạng Thái Focus Nhóm Form]
\sffamily\small
\textbf{1. Khi chưa focus (Trạng thái đơn thuần bình thường):}\\[2pt]
\tcbox[on line, arc=6pt, colback=white, colframe=gray!40, left=8pt, right=40pt, top=6pt, bottom=6pt]{
  \textbf{Họ tên:}\\[2pt]
  \tcbox[on line, arc=3pt, colback=white, colframe=gray!50, left=6pt, right=30pt, top=2pt, bottom=2pt]{\color{gray!70}\small Nhập tên...}
} \quad \textcolor{gray!80}{\bfseries (Viền xám, Nền trắng)}\\[10pt]
\textbf{2. Khi người dùng focus - nhấp chuột vào ô Input bên trong:}\\[2pt]
\tcbox[on line, arc=6pt, colback=blue!5, colframe=blue, left=8pt, right=40pt, top=6pt, bottom=6pt]{
  \textbf{Họ tên:}\\[2pt]
  \tcbox[on line, arc=3pt, colback=white, colframe=blue, left=6pt, right=30pt, top=2pt, bottom=2pt]{\color{black}\small Nguyễn Trọng Tấn $|$}
} \quad \textcolor{blue}{\bfseries (Khối cha đổi viền xanh, Nền xanh nhạt \#f0f8ff)}
\end{browserWindow}

% ------------------------------------------------------------------------------
\subsection{3. Ví Dụ 2: Khung Tìm Kiếm Tự Động Hiển Thị Menu Gợi Ý (Search Box Dropdown)}

\begin{codeWindow}[Mã HTML \& CSS hiển thị danh sách từ khóa gợi ý khi ô Input Search được focus]
\begin{lstlisting}[language=HTML5]
<!DOCTYPE html>
<html lang="vi">
<head>
  <meta charset="UTF-8">
  <title>Search Box :focus-within</title>
  <style>
    .search-box {
      position: relative;
      width: 320px;
      border: 2px solid #ddd;
      border-radius: 20px;
      padding: 8px 16px;
      transition: all 0.3s ease;
      background-color: #fff;
    }
    /* Khi input ben trong duoc focus -> viem xanh & do bong Glow */
    .search-box:focus-within {
      border-color: #1a73e8;
      box-shadow: 0 4px 16px rgba(26, 115, 232, 0.25);
    }
    .search-input {
      border: none;
      outline: none;
      width: 100%;
      font-size: 14px;
    }
    /* Dynamic Dropdown */
    .search-dropdown {
      display: none; /* Mac dinh an */
      position: absolute;
      top: 100%; left: 0; right: 0;
      background: #fff;
      border: 1px solid #ddd;
      border-radius: 8px;
      margin-top: 6px;
      box-shadow: 0 4px 12px rgba(0,0,0,0.1);
    }
    /* Tu dong hien thi dropdown khi co focus (*@bên@*) trong */
    .search-box:focus-within .search-dropdown {
      display: block;
    }
  </style>
</head>
<body>
  <div class="search-box">
    <input type="search" class="search-input" placeholder="(*@Tìm kiếm từ khóa...@*)">
    <div class="search-dropdown">
      <div style="padding: 8px;">(*@1. Lập trình Web HTML5/CSS3@*)</div>
      <div style="padding: 8px;">(*@2. Giả lớp CSS :focus-within@*)</div>
    </div>
  </div>
</body>
</html>
\end{lstlisting}
\end{codeWindow}

\begin{browserWindow}[Kết Quả Hiển Thị Trình Duyệt: Thanh Tìm Kiếm Đổ Bóng \& Mở Menu Gợi Ý]
\sffamily\small
\textbf{1. Trạng thái ban đầu (Chưa focus):}\\[2pt]
\tcbox[on line, arc=15pt, colback=white, colframe=gray!40, left=10pt, right=40pt, top=4pt, bottom=4pt]{\color{gray!70}\small [Kính lúp] Tìm kiếm từ khóa...} \quad \textcolor{gray!70}{\bfseries (Menu ẩn)}\\[8pt]
\textbf{2. Khi người dùng nhấp chuột focus vào ô tìm kiếm:}\\[2pt]
\tcbox[on line, arc=15pt, colback=white, colframe=primaryBlue, boxsep=2pt, left=10pt, right=40pt, top=4pt, bottom=4pt]{\color{black}\small [Kính lúp] Lập trình web $|$}\\[-4pt]
\tcbox[on line, arc=4pt, colback=white, colframe=gray!30, left=10pt, right=20pt, top=4pt, bottom=4pt]{
  \begin{minipage}{5.5cm}
  \textcolor{primaryBlue}{\bfseries [Kính lúp] Từ khóa gợi ý tự động:}\\[2pt]
  1. Lập trình Web HTML5/CSS3\\[2pt]
  2. Giả lớp CSS :focus-within
  \end{minipage}
} \quad \textcolor{primaryBlue}{\bfseries (Viền xanh sáng, Tự hiện Menu Gợi ý không cần JS)}
\end{browserWindow}

% ------------------------------------------------------------------------------
\subsection{4. Ví Dụ 3: Floating Label Màu Tím Neon Tối Ưu Tương Tác}

\begin{codeWindow}[Mã HTML \& CSS kỹ thuật Floating Label nhãn trượt nổi màu tím neon]
\begin{lstlisting}[language=HTML5]
<!DOCTYPE html>
<html lang="vi">
<head>
  <meta charset="UTF-8">
  <style>
    .field-wrapper {
      position: relative;
      margin-top: 15px;
      border: 2px solid #ccc;
      border-radius: 8px;
      padding: 12px;
      transition: 0.3s;
    }
    .field-wrapper input {
      border: none;
      outline: none;
      width: 100%;
    }
    .field-wrapper label {
      position: absolute;
      top: 12px; left: 12px;
      color: #777;
      transition: 0.2s ease-all;
      background-color: white;
      padding: 0 4px;
    }
    /* Khi nhap text HOAC dang focus -> Nhan label bay len mep tren */
    .field-wrapper:focus-within label,
    .field-wrapper input:not(:placeholder-shown) + label {
      top: -10px;
      left: 10px;
      font-size: 12px;
      color: #6200ee;
      font-weight: bold;
    }
    .field-wrapper:focus-within {
      border-color: #6200ee;
    }
  </style>
</head>
<body>
  <div class="field-wrapper">
    <input type="text" id="phone" placeholder=" ">
    <label for="phone">(*@Số điện thoại liên hệ@*)</label>
  </div>
</body>
</html>
\end{lstlisting}
\end{codeWindow}

\begin{browserWindow}[Kết Quả Hiển Thị Trình Duyệt: Floating Label Màu Tím Neon]
\sffamily\small
\textbf{1. Trạng thái rỗng chưa focus:}\\[2pt]
\tcbox[on line, arc=4pt, colback=white, colframe=gray!40, left=8pt, right=50pt, top=8pt, bottom=8pt]{\color{gray!70}\small Số điện thoại liên hệ}\\[8pt]
\textbf{2. Khi focus nhấp chuột (Label trượt nổi lên mép trên):}\\[2pt]
\tcbox[on line, arc=4pt, colback=white, colframe=purple!80!black, left=8pt, right=50pt, top=8pt, bottom=8pt]{
  \textcolor{purple!80!black}{\bfseries\scriptsize [Số điện thoại liên hệ]}\\[-2pt]
  {\small 0988 123 456 $|$}
} \quad \textcolor{purple!80!black}{\bfseries (Label nổi viền tím Neon)}
\end{browserWindow}

% ------------------------------------------------------------------------------
\subsection{5. Ví Dụ 4: Giao Diện Multi-Card Form -- Nổi Bật Thẻ Card Đang Thao Tác}

\begin{codeWindow}[Mã HTML \& CSS làm nổi bật Card form đang được người dùng thao tác]
\begin{lstlisting}[language=HTML5]
<!DOCTYPE html>
<html lang="vi">
<head>
  <style>
    .form-card {
      border: 1.5px solid #e0e0e0;
      border-radius: 10px;
      padding: 16px;
      margin-bottom: 15px;
      opacity: 0.65; /* Mo nhe khi chua thao tac */
      transition: all 0.3s ease;
      background-color: #fafafa;
    }
    /* Card nao chua o input dang focus se sang ro va noi len */
    .form-card:focus-within {
      opacity: 1;
      border-color: #00c853;
      background-color: #ffffff;
      box-shadow: 0 6px 20px rgba(0, 200, 83, 0.15);
      transform: translateY(-2px);
    }
  </style>
</head>
<body>
  <div class="form-card">
    <h3>(*@Thẻ 1: Thông tin cá nhân@*)</h3>
    <input type="text" placeholder="(*@Họ và tên@*)">
  </div>
  <div class="form-card">
    <h3>(*@Thẻ 2: Địa chỉ giao hàng@*)</h3>
    <input type="text" placeholder="(*@Số nhà, Tên đường@*)">
  </div>
</body>
</html>
\end{lstlisting}
\end{codeWindow}

\begin{browserWindow}[Kết Quả Hiển Thị Trình Duyệt: Nổi Bật Card Tương Tác]
\sffamily\small
\textbf{1. Card 1 (Chưa focus):}\\[2pt]
\tcbox[on line, arc=6pt, colback=gray!10, colframe=gray!30, left=8pt, right=50pt, top=4pt, bottom=4pt]{\color{gray!70}\small Thẻ 1: Thông tin cá nhân (Mờ 65\%)}\\[8pt]
\textbf{2. Card 2 (Đang thao tác gõ số nhà - Card tự nổi sáng 100\%):}\\[2pt]
\tcbox[on line, arc=6pt, colback=white, colframe=green!70!black, boxsep=2pt, left=10pt, right=60pt, top=6pt, bottom=6pt]{
  \textcolor{green!70!black}{\bfseries Thẻ 2: Địa chỉ giao hàng}\\[2pt]
  \tcbox[on line, arc=3pt, colback=white, colframe=green!70!black, left=6pt, right=30pt, top=2pt, bottom=2pt]{\small 123 Đường Nguyễn Huệ, Q1 $|$}
} \quad \textcolor{green!70!black}{\bfseries (Nổi sáng 100\%, Viền xanh lá, Bóng mờ Card)}
\end{browserWindow}

\begin{tipBox}[3 Ứng Dụng Thực Tế Đỉnh Cao Của :focus-within]
\begin{enumerate}[leftmargin=1.5em, itemsep=4pt, parsep=0pt]
    \item \textbf{Highlight nhóm form}: Làm nổi bật toàn bộ thẻ cha chứa ô nhập liệu khi người dùng thao tác, giúp cải thiện tầm nhìn trực quan.
    \item \textbf{Floating label}: Thẻ label tự động đổi màu hoặc nổi lên khi input bên trong nhận focus.
    \item \textbf{Tạo hiệu ứng trực quan cho UX}: Xây dựng các hiệu ứng tương tác cao cấp hoàn toàn bằng CSS mà không cần viết mã JavaScript.
\end{enumerate}
\end{tipBox}

\begin{warnBox}[Các Lưu Ý Khi Sử Dụng Giả Lớp :focus-within]
\begin{itemize}[leftmargin=1.5em, itemsep=4pt, parsep=0pt]
    \item \textbf{Khác biệt hoàn toàn với \codeinline{:focus}}: Giả lớp \codeinline{:focus} chỉ áp dụng trực tiếp cho chính bản thân phần tử đang nhận tiêu điểm, trong khi \codeinline{:focus-within} áp dụng cho phần tử cha chứa nó (rất tiện lợi khi muốn định kiểu toàn bộ khối \term{Block}).
    \item \textbf{Độ tương thích trình duyệt}: Được hỗ trợ đầy đủ trên tất cả các trình duyệt hiện đại (Chrome, Firefox, Edge, Safari).
\end{itemize}
\end{warnBox}


\section{Truy Cập Web Cho Biểu Mẫu (Form Accessibility - a11y)}
\label{sec:form_accessibility}

\begin{conceptBox}[Tiêu Chuẩn Accessibility (a11y) Trong HTML Form]
\term{Web Accessibility (a11y)} trong biểu mẫu là việc thiết kế và lập trình sao cho tất cả người dùng, bao gồm cả **người khiếm thị hoặc người khuyết tật sử dụng trình đọc màn hình (Screen Reader)** hoặc chỉ dùng bàn phím (Keyboard navigation), đều có thể hiểu và tương tác với form dễ dàng.
\end{conceptBox}

\subsection{1. Các Quy Tắc Vàng Giúp Form Đạt Chuẩn WCAG}
\begin{enumerate}[leftmargin=1.5em, itemsep=4pt, parsep=0pt]
    \item \textbf{Liên kết thẻ \codeinline{<label>} chuẩn xác}: Mỗi ô input bắt buộc phải có thẻ \codeinline{<label>} liên kết bằng cặp thuộc tính \codeinline{for="..."} và \codeinline{id="..."}.
    \item \textbf{Sử dụng thuộc tính ARIA}:
    \begin{itemize}[leftmargin=1.2em, itemsep=2pt]
        \item \codeinline{aria-describedby="..."}: Liên kết ô input với đoạn văn bản hướng dẫn hoặc thông báo lỗi bên dưới.
        \item \codeinline{aria-invalid="true"}: Thông báo cho trình đọc màn hình biết ô input đang bị lỗi dữ liệu.
        \item \codeinline{aria-required="true"}: Đánh dấu trường bắt buộc nhập cho công cụ trợ năng.
    \end{itemize}
    \item \textbf{Cung cấp nhãn ẩn cho Screen Reader (\codeinline{.sr-only})}: Đối với các nút bấm chỉ có Icon (như nút Tìm kiếm kính lúp), hãy thêm văn bản ẩn để Screen Reader đọc được tên nút.
\end{enumerate}

\begin{codeWindow}[Mã HTML form chuẩn truy cập Accessibility a11y]
\begin{lstlisting}[language=HTML5]
<form>
  <!-- 1. Lien ket label bang for va id -->
  <label for="username">Ten dang nhap:</label>
  <input type="text" id="username" name="username" 
         aria-describedby="username-hint" aria-required="true">
  <small id="username-hint">Ten dang nhap tu 6-20 ky tu.</small>
  
  <!-- 2. Nung Nut bam Icon co nhan an Screen Reader -->
  <button type="submit" aria-label="Tim kiem san pham">
    <svg width="16" height="16"><path d="..."/></svg>
    <span class="sr-only">Tim kiem</span>
  </button>
</form>
\end{lstlisting}
\end{codeWindow}

\begin{browserWindow}[Kết Quả Trình Duyệt: Form Đạt Chuẩn Trợ Năng a11y]
\sffamily\small
\textbf{Tên đăng nhập:}\\[2pt]
\tcbox[on line, arc=3pt, colback=white, colframe=primaryBlue!40, left=6pt, right=50pt, top=3pt, bottom=3pt]{\small nguyenvana} \\[2pt]
\textcolor{gray!70}{\small\textit{Tên đăng nhập từ 6-20 ký tự (Đã liên kết aria-describedby)}}\\[8pt]
\tcbox[on line, arc=3pt, colback=primaryBlue, colframe=primaryBlue, left=10pt, right=10pt, top=3pt, bottom=3pt]{\color{white}\bfseries\small [SVG Icon] Tìm kiếm}
\end{browserWindow}

\subsection{2. Thuộc Tính Điều Hướng Bàn Phím tabindex Trong HTML}
\label{subsec:html_tabindex}

\begin{conceptBox}[Thuộc Tính tabindex Trong HTML Là Gì?]
\codeinline{tabindex} là một thuộc tính toàn cục (\term{Global Attribute}) trong HTML. Nó xác định thứ tự tiêu điểm focus của các phần tử khi người dùng nhấn phím **Tab** trên bàn phím.
\par\vspace{4pt}
\begin{itemize}[leftmargin=1.5em, itemsep=2pt, parsep=0pt]
    \item \textbf{Tác dụng cốt lõi}: Giúp cải thiện trải nghiệm điều hướng bằng bàn phím (\term{Keyboard Navigation}), là yêu cầu sống còn cho tiêu chuẩn truy cập Web Accessibility (WCAG, ARIA).
\end{itemize}
\end{conceptBox}

\begin{prettyTableBox}[Bảng Chi Tiết Các Giá Trị Của Thuộc Tính tabindex]
\small
\rowcolors{2}{white}{primaryBlue!4}
\begin{tabularx}{\linewidth}{K{2.8cm} Z K{3.5cm}}
\rowcolor{primaryBlue}
\thd{Giá Trị tabindex} & \thd{Ý Nghĩa \& Hành Vi Tiêu Điểm} & \thd{Ví Dụ HTML} \\
\codeinline{tabindex="0"} & Phần tử có thể focus được, và thứ tự focus tuân theo thứ tự xuất hiện tự nhiên trong mã DOM HTML. & \codeinline{<div tabindex="0">Có thể focus</div>} \\
\codeinline{tabindex="-1"} & Phần tử có thể focus bằng mã JavaScript (qua phương thức \codeinline{element.focus()}), nhưng sẽ **bị bỏ qua** khi người dùng bấm phím Tab. & \codeinline{<div tabindex="-1">Bỏ qua khi Tab</div>} \\
\codeinline{tabindex="n"} ($n > 0$) & Đặt thứ tự focus cụ thể ưu tiên trước thứ tự DOM tự nhiên (số nhỏ hơn được focus trước). & \codeinline{<input type="text" tabindex="1">} \\
\end{tabularx}
\end{prettyTableBox}

\subsection{3. Ví Dụ Minh Họa Chi Tiết Về Thứ Tự Tab}

\begin{codeWindow}[Mã HTML minh họa các giá trị tabindex khác nhau]
\begin{lstlisting}[language=HTML5]
<!DOCTYPE html>
<html lang="vi">
<head>
  <meta charset="UTF-8">
  <title>Vi du tabindex</title>
</head>
<body>
  <h2>(*@Thứ *@tự@*) Tab@*)</h2>
  <!-- Theo DOM mac dinh (tabindex=0) -->
  <input type="text" placeholder="(*@Tên (DOM)@*)" tabindex="0"><br>
  <input type="text" placeholder="Email (DOM)" tabindex="0"><br>
  
  <!-- Bo qua khi nhan phim Tab (tabindex=-1) -->
  <button tabindex="-1">(*@Nút bị bỏ qua@*)</button><br>
  
  <!-- Custom thu tu uu tien (tabindex > 0) -->
  <input type="text" placeholder="Custom 1" tabindex="2"><br>
  <input type="text" placeholder="Custom 2" tabindex="1"><br>
</body>
</html>
\end{lstlisting}
\end{codeWindow}

\begin{browserWindow}[Kết Quả Hành Vi Điều Hướng Phím Tab Khi Người Dùng Thao Tác]
\sffamily\small
\textbf{Thứ tự di chuyển tiêu điểm khi nhấn phím Tab liên tiếp:}\\[4pt]
1. \tcbox[on line, arc=3pt, colback=white, colframe=blue, left=4pt, right=20pt, top=2pt, bottom=2pt]{\small Custom 2} \quad \textcolor{blue}{\bfseries (tabindex=1 được focus đầu tiên)}\\[4pt]
2. \tcbox[on line, arc=3pt, colback=white, colframe=blue, left=4pt, right=20pt, top=2pt, bottom=2pt]{\small Custom 1} \quad \textcolor{blue}{\bfseries (tabindex=2 được focus thứ hai)}\\[4pt]
3. \tcbox[on line, arc=3pt, colback=white, colframe=blue, left=4pt, right=20pt, top=2pt, bottom=2pt]{\small Tên (DOM)} \quad \textcolor{primaryBlue}{\bfseries (Tiếp tục sang phần tử DOM mặc định 1)}\\[4pt]
4. \tcbox[on line, arc=3pt, colback=white, colframe=blue, left=4pt, right=20pt, top=2pt, bottom=2pt]{\small Email (DOM)} \quad \textcolor{primaryBlue}{\bfseries (Tiếp tục sang phần tử DOM mặc định 2)}\\[4pt]
5. \tcbox[on line, arc=3pt, colback=gray!20, colframe=gray!40, left=4pt, right=20pt, top=2pt, bottom=2pt]{\small Nút bị bỏ qua} \quad \textcolor{red!70!black}{\bfseries (Đã bị BỎ QUA hoàn toàn do tabindex=-1)}
\end{browserWindow}

\begin{warnBox}[4 Lưu Ý Nằm Lòng Khi Sử Dụng Thuộc Tính tabindex]
\begin{enumerate}[leftmargin=1.5em, itemsep=4pt, parsep=0pt]
    \item \textbf{Tránh lạm dụng số dương (\codeinline{tabindex > 0})}: Việc đặt các số nguyên dương sẽ ghi đè thứ tự DOM tự nhiên, dễ gây rối loạn và mất phương hướng cho người dùng chỉ dùng bàn phím. **Khuyên dùng**: Chỉ sử dụng \codeinline{0} hoặc \codeinline{-1}.
    \item \textbf{Biến thẻ không tương tác thành focusable (\codeinline{tabindex="0"})}: Dùng cho các thẻ mặc định không thể focus (như \codeinline{<div>}, \codeinline{<span>}) để người dùng bàn phím có thể Tab tới và tương tác.
    \item \textbf{Focus bằng JavaScript với \codeinline{tabindex="-1"}}: Hữu ích khi cần chuyển tiêu điểm bằng script: \codeinline{document.getElementById("myDiv").focus();}, nhưng không muốn người dùng tự Tab tới phần tử đó.
    \item \textbf{Các thẻ form mặc định không cần \codeinline{tabindex="0"}}: Các phần tử như \codeinline{<input>}, \codeinline{<button>}, \codeinline{<select>}, \codeinline{<a href>} đã mặc định có thể nhận focus, do đó không cần khai báo thêm \codeinline{tabindex="0"}.
\end{enumerate}
\end{warnBox}



